\section{Propiedades construibles}
\label{section:IV.9}

Sea $S$ un preesquema, sea $f:X\to S$ un morfismo de \emph{presentación finita}~\sref{IV.1.6.1},
y sea $\sh{F}$ un $\sh{O}_X$-Módulo cuasicoherente de \emph{presentación finita}.
Nos proponemos en este parágrafo dar criterios que aseguren, por ejemplo, que el conjunto de los $s\in S$ tales que el $\mathbf{k}(s)$-preesquema $X_s=f^{-1}(s)=X\times_S\operatorname{Spec}(\mathbf{k}(s))$ posee cierta propiedad, o tales que el $\sh{O}_{X_s}$-Módulo $\sh{F}_s=\sh{F}\otimes_{\sh{O}_X}\mathbf{k}(s)$ posee cierta propiedad, es \emph{localmente construible}, o al menos \emph{ind-construible}~\sref{IV.1.9.4}; veremos que este es el caso para la mayoría de las propiedades que aparecen en Geometría Algebraica.
Suponiendo únicamente que $f$ es \emph{localmente de presentación finita}, daremos también en~{\hyperlink{ega4.c04.subsec.9-9}{(9.9)}} criterios para que el conjunto de puntos $x\in X$ en los que la fibra $X_{f(x)}$ (o el $\sh{O}_{X_{f(x)}}$-Módulo $\sh{F}_{f(x)}$) posee cierta propiedad sea \emph{localmente construible}.
Veremos en \textsection12 que estos resultados, junto con la hipótesis adicional de que $f$ sea \emph{plano} (resp.\ \emph{propio y plano}), permiten demostrar que los conjuntos considerados en $X$ (resp.\ en $S$) son incluso \emph{abiertos} en numerosos casos.

\subsection{El principio de extensión finita}
\label{subsection:IV.9.1}

\begin{proposition}[9.1.1]
\label{IV.9.1.1}
\textnormal{(Principio de extensión finita).}
Sea $k$ un cuerpo y sea $\mathcal{C}$ un conjunto de extensiones de $k$.
Supongamos que se satisfacen las condiciones siguientes:
\begin{enumerate}[label=(\roman*)]
  \item Si $K\in\mathcal{C}$ y existe un $k$-homomorfismo $K\to K'$ (donde $K'$ pertenece al universo en el que nos hemos situado), entonces $K'\in\mathcal{C}$.
  \item Si $K\in\mathcal{C}$, existe una subextensión $K'$ de $K$, de tipo finito sobre $k$, tal que $K'\in\mathcal{C}$.
  \item Si $K\in\mathcal{C}$ es el cuerpo de fracciones de una $k$-álgebra íntegra de tipo finito $A$, existe $f\in A-\{0\}$ tal que, para todo ideal maximal $\mathfrak{m}$ de $A$, se tiene $A_f/\mathfrak{m}\in\mathcal{C}$.
\end{enumerate}

\oldpage[IV-3]{55}

Sea entonces $\Omega$ una extensión algebraicamente cerrada de $k$ (en el universo considerado).
Las condiciones siguientes son equivalentes:
\begin{enumerate}[label=\emph{\alph*)}]
  \item $\mathcal{C}$ no es vacío.
  \item Existe $K'\in\mathcal{C}$ que es una extensión finita de $k$.
  \item Se tiene $\Omega\in\mathcal{C}$.
\end{enumerate}
\end{proposition}

La condición~(i) implica evidentemente que \emph{b)} entraña \emph{c)}, y \emph{c)} entraña trivialmente \emph{a)};
demostremos que \emph{a)} entraña \emph{b)}.
En virtud de~(ii) y~(iii), existe una extensión $K'\in\mathcal{C}$ de la forma $A/\mathfrak{m}$, donde $A$ es una $k$-álgebra de tipo finito sobre $k$ y $\mathfrak{m}$ es un ideal maximal de $A$.
Sabemos, por el Nullstellensatz de Hilbert (Bourbaki, \emph{Alg.\ comm.}, cap.~V, \textsection3, n\textsuperscript{o}~1, cor.~2 del teo.~1), que $K'$ es una extensión finita de $k$.

\begin{corollary}[9.1.2]
\label{IV.9.1.2}
Bajo las hipótesis~\emph{(i)}, \emph{(ii)}, \emph{(iii)} de~\sref{IV.9.1.1}, si $k$ es algebraicamente cerrado y $\mathcal{C}$ no es vacío, $\mathcal{C}$ contiene todas las extensiones de $k$ del universo considerado.
\end{corollary}

En efecto, como \emph{a)} entraña \emph{c)}, se tiene $k\in\mathcal{C}$, y la conclusión se sigue de la hipótesis~(i).

\begin{remark}[9.1.3]
\label{IV.9.1.3}
En la práctica se verificará la condición~(ii) de~\sref{IV.9.1.1} observando que $K$ es el límite directo de sus subextensiones de tipo finito $K_\alpha$, y aplicando los resultados de \textsection8, teniendo en cuenta eventualmente que, para $K_\alpha\subset K_\beta$, $K_\beta$ es fielmente plana sobre $K_\alpha$.
Con frecuencia, el conjunto $\mathcal{C}$ está formado por \emph{cuerpos} que pertenecen a un conjunto $\mathcal{C}'$ de $k$-álgebras que satisface la condición siguiente:
\begin{enumerate}[label=(\roman*~bis)]
  \item Si $A\in\mathcal{C}'$ y existe un homomorfismo de $k$-álgebras $A\to A'$ (donde $A'$ pertenece al universo en el que nos hemos situado), entonces $A'\in\mathcal{C}'$.
\end{enumerate}
\end{remark}

En tal caso, la condición~(i) se satisface trivialmente, y se verificará en general la condición~(iii) de~\sref{IV.9.1.1} observando que el cuerpo de fracciones $K$ de $A$ es el límite directo de las álgebras $A_f$ (para la relación $D(f)\supset D(g)$), y aplicando los resultados de \textsection8, teniendo en cuenta eventualmente que los morfismos $D(g)\to D(f)$ son inmersiones abiertas.

Por otra parte, cuando no se satisface (i~bis), la demostración de~(iii) suele ser más delicada y está ligada a criterios de construibilidad que se desarrollarán más adelante.

He aquí algunos ejemplos típicos de aplicación del principio de extensión finita:

\begin{proposition}[9.1.4]
\label{IV.9.1.4}
Sea $k$ un cuerpo, sea $\Omega$ una extensión algebraicamente cerrada de $k$, y sean $X$ e $Y$ dos preesquemas de tipo finito sobre $k$.
Las condiciones siguientes son equivalentes:
\begin{enumerate}[label=\emph{\alph*)}]
  \item Existe un $\Omega$-morfismo $X_{(\Omega)}\to Y_{(\Omega)}$ \textnormal{(}resp.\ un $\Omega$-morfismo que posee una de las propiedades~\emph{(i)} a~\emph{(xiv)} especificada de~\sref{IV.8.10.5}\textnormal{)}.
  \item Existe una extensión finita $k'$ de $k$ y un $k'$-morfismo $X_{(k')}\to Y_{(k')}$ \textnormal{(}resp.\ un $k'$-morfismo que posee la propiedad considerada\textnormal{)}.
  \item Existe una extensión $K$ de $k$ y un $K$-morfismo $X_{(K)}\to Y_{(K)}$ \textnormal{(}resp.\ un $K$-morfismo que posee la propiedad considerada\textnormal{)}.
\end{enumerate}
\end{proposition}

Aplicamos la observación~\sref{IV.9.1.3} tomando como $\mathcal{C}'$ el conjunto de todas las $k$-álgebras $A$ (del universo en el que nos hemos situado) tales que existe un $A$-morfismo $X\otimes_k A\to Y\otimes_k A$ (resp.\ un morfismo que posee la propiedad particular de~\sref{IV.8.10.5} considerada).
La condición~(i~bis) de~\sref{IV.9.1.3} se satisface entonces gracias a que las
\oldpage[IV-3]{56}
propiedades consideradas en~\sref{IV.8.10.5} son todas estables por cambio de base.
La condición~(iii) de~\sref{IV.9.1.1} se satisface así en virtud de~\sref{IV.8.8.2}[i] (resp.\ \sref{IV.8.10.5}), pues $\operatorname{Spec}(k)$ es cuasicompacto y cuasiseparado.
Queda por verificar la condición~(ii) de~\sref{IV.9.1.1}, que se sigue también de~\sref{IV.8.8.2}[i] y~\sref{IV.8.10.5}, considerando $K$ como límite directo de sus subextensiones de tipo finito.
Concluimos, por tanto, mediante~\sref{IV.9.1.1}.

En particular, si existe una extensión $K$ de $k$ y un $K$-isomorfismo $X_{(K)}\xrightarrow{\sim} Y_{(K)}$, decimos que $X$ e $Y$ son \emph{geométricamente isomorfos}.

El corolario siguiente generaliza~(\textbf{II},~6.6.5):

\begin{corollary}[9.1.5]
\label{IV.9.1.5}
Sea $k$ un cuerpo y sea $X$ un $k$-preesquema.
Si existe una extensión $K$ de $k$ tal que $X_{(K)}$ es proyectivo \textnormal{(}resp.\ cuasiproyectivo\textnormal{)} sobre $K$, entonces $X$ es proyectivo \textnormal{(}resp.\ cuasiproyectivo\textnormal{)} sobre $k$.
\end{corollary}

Como el morfismo $\operatorname{Spec}(K)\to\operatorname{Spec}(k)$ es fielmente plano y cuasicompacto, ya se sigue de~\sref{IV.2.7.1}[v] que $X$ es de tipo finito sobre $k$.
La hipótesis significa que existe una inmersión cerrada (resp.\ una inmersión) $X_{(K)}\to\mathbf{P}_K^n=\mathbf{P}_k^n\otimes_k K$ (\textbf{II},~5.5.4,~(ii) y~5.5.3); aplicando~\sref{IV.9.1.4} a la propiedad~(iv) (resp.~(ii)) de~\sref{IV.8.10.5}, deducimos que, para alguna extensión finita $k'$ de $k$, existe una inmersión cerrada (resp.\ una inmersión) $X_{(k')}\to\mathbf{P}_{k'}^n$; en otras palabras, $X_{(k')}$ es proyectivo (resp.\ cuasiproyectivo) sobre $k'$.
Concluimos entonces mediante~(\textbf{II},~6.6.5).

\begin{proposition}[9.1.6]
\label{IV.9.1.6}
Sea $k$ un cuerpo, sea $\Omega$ una extensión algebraicamente cerrada de $k$, sea $X$ un preesquema de tipo finito sobre $k$, y sean $\sh{F}$, $\sh{G}$ dos $\sh{O}_X$-Módulos coherentes.
Supongamos que existe un isomorfismo $\sh{F}\otimes_k\Omega\xrightarrow{\sim}\sh{G}\otimes_k\Omega$.
Entonces existe una extensión finita $k'$ de $k$ y un isomorfismo $\sh{F}\otimes_k k'\xrightarrow{\sim}\sh{G}\otimes_k k'$.
\end{proposition}

El razonamiento es el mismo que en~\sref{IV.9.1.4}, aplicando~\sref{IV.8.5.2}[i] (en la demostración de la propiedad~(iii) de~\sref{IV.9.1.1}, utilizamos aquí el hecho de que los morfismos $D(g)\to D(f)$ (con las notaciones de~\sref{IV.9.1.3}) son inmersiones abiertas y, \emph{a fortiori}, morfismos planos).

\subsection{Propiedades construibles e ind-construibles}
\label{subsection:IV.9.2}
\label{IV.9.2}

\begin{definition}[9.2.1]
\label{IV.9.2.1}
Sea $\mathbf{P}(X,\sh{F},k)$ una relación.
Decimos que $\mathbf{P}$ es una propiedad \emph{construible} (resp.\ \emph{ind-construible}) de los preesquemas algebraicos si se satisfacen las dos condiciones siguientes:
\begin{enumerate}[label=(\roman*)]
  \item Si $k$ es un cuerpo, $X$ un preesquema algebraico sobre $k$, $\sh{F}$ un $\sh{O}_X$-Módulo coherente y $k'$ una extensión de $k$, para que $\mathbf{P}(X,\sh{F},k)$ sea cierta es necesario y suficiente que $\mathbf{P}(X_{(k')},\sh{F}\otimes_k k',k')$ sea cierta.
  \item Sea $S$ un preesquema noetheriano íntegro, de punto genérico $\eta$, sea $u:X\to S$ un morfismo de tipo finito y sea $\sh{F}$ un $\sh{O}_X$-Módulo coherente.
Para todo $s\in S$, pongamos $X_s=u^{-1}(s)=X\times_S\operatorname{Spec}(\mathbf{k}(s))$, $\sh{F}_s=\sh{F}\otimes_{\sh{O}_S}\mathbf{k}(s)$.
Sea $E$ el conjunto de los $s\in S$ tales que $\mathbf{P}(X_s,\sh{F}_s,\mathbf{k}(s))$ es cierta.
Entonces uno de los dos conjuntos $E$, $S-E$ (resp.\ el conjunto $E$) contiene un abierto no vacío (y, por consiguiente, es un entorno de $\eta$) (resp.\ contiene un abierto no vacío si contiene a $\eta$).
\end{enumerate}
\end{definition}

\oldpage[IV-3]{57}

\begin{remarks}[9.2.2]
\label{IV.9.2.2}
\begin{enumerate}[label=(\roman*)]
  \item Se trata, desde luego, de una convención de lenguaje metamatemático y no de una definición matemática propiamente dicha.
Tenemos ``definiciones'' análogas para relaciones entre $k$, uno o varios $k$-preesquemas algebraicos, $k$-morfismos entre estos preesquemas, Módulos coherentes sobre estos preesquemas u homomorfismos entre estos Módulos.
  \item También tendremos que considerar relaciones en las que intervengan \emph{partes construibles} de preesquemas.
Por ejemplo, sea $\mathbf{P}(X,Z,k)$ una relación; diremos (por abuso de lenguaje) que $\mathbf{P}$ es una propiedad construible (resp.\ ind-construible) de la parte construible $Z$ de $X$ si se satisfacen las dos condiciones siguientes:
\begin{enumerate}[label=$\arabic*^{\mathrm{o}}$]
  \item Si $k$ es un cuerpo, $X$ un preesquema algebraico sobre $k$, $Z$ una \emph{parte construible} de $X$ y $k'$ una extensión de $k$, para que $\mathbf{P}(X,Z,k)$ sea cierta es necesario y suficiente que $\mathbf{P}(X_{(k')},p^{-1}(Z),k')$ sea cierta (siendo $p:X_{(k')}\to X$ la proyección canónica).
  \item Sea $S$ un preesquema noetheriano íntegro, de punto genérico $\eta$, sea $u:X\to S$ un morfismo de tipo finito y sea $Z$ una \emph{parte construible de $X$}.
Para todo $s\in S$, pongamos $X_s=u^{-1}(s)$, $Z_s=Z\cap X_s$.
Sea $E$ el conjunto de los $s\in S$ tales que $\mathbf{P}(X_s,Z_s,\mathbf{k}(s))$ es cierta.
Entonces uno de los dos conjuntos $E$, $S-E$ (resp.\ el conjunto $E$) contiene un abierto no vacío (resp.\ contiene un abierto no vacío si contiene a $\eta$).
\end{enumerate}

Observemos que en la condición~$2^\circ$ hay que suponer que $Z$ es una parte construible de $X$, y no solamente que $Z_s$ es una parte construible de $X_s$ para todo $s$; la primera de estas dos propiedades entraña la segunda~\sref{IV.1.8.2}, pero no recíprocamente.

  \item Si $\mathbf{P}$ es una propiedad construible, está claro que lo mismo ocurre con ``no $\mathbf{P}$''.
Si $\mathbf{P}$, $\mathbf{Q}$ son dos propiedades construibles (resp.\ ind-construibles), lo mismo ocurre con las propiedades ``$\mathbf{P}$ o $\mathbf{Q}$'' y ``$\mathbf{P}$ y $\mathbf{Q}$''; en efecto, bajo las hipótesis de~\sref{IV.9.2.1}[ii], si $E$, $E'$ son dos partes de $S$ y $E$ contiene un abierto no vacío, lo mismo ocurre con $E\cup E'$, y si $S-E$ y $S-E'$ contienen sendos abiertos no vacíos, lo mismo ocurre con $S-(E\cup E')=(S-E)\cap(S-E')$.

  \item Sea $\mathbf{P}(X,\sh{F},k)$ una relación que satisface la condición~\sref{IV.9.2.1}[i]; sea $S$ un preesquema y sea $u:X\to S$ un morfismo de tipo finito; con las notaciones de~\sref{IV.9.2.1}[ii], sea $E$ el conjunto de los $s\in S$ tales que $\mathbf{P}(X_s,\sh{F}_s,\mathbf{k}(s))$ es cierta.
Sea, por otra parte, $g:S'\to S$ un morfismo cualquiera, y pongamos $X'=X_{(S')}$, $\sh{F}'=\sh{F}\otimes_{\sh{O}_S}\sh{O}_{S'}$; de la transitividad de las fibras~(\textbf{I},~3.6.4) y de la condición~\sref{IV.9.2.1}[i] se sigue entonces que \emph{el conjunto de los $s'\in S'$ tales que $\mathbf{P}(X'_{s'},\sh{F}'_{s'},\mathbf{k}(s'))$ es cierta es igual a $g^{-1}(E)$}.
Esto se extiende inmediatamente al caso en que aparecen en $\mathbf{P}$ varios preesquemas, Módulos, morfismos de preesquemas u homomorfismos de Módulos, así como a las propiedades del tipo considerado en~(ii).

  \item Como veremos en el resto de este parágrafo y en el resto del cap.~IV, la mayoría de las propiedades $\mathbf{P}$ que satisfacen la condición~\sref{IV.9.2.1}[i] satisfacen también la condición~\sref{IV.9.2.1}[ii].
Como posibles excepciones, señalemos la propiedad de ser \emph{proyectivo}, o \emph{cuasiproyectivo}, o \emph{afín}, o \emph{cuasiafín} sobre el cuerpo base (para un preesquema algebraico); veremos en~\sref{IV.9.6.2} que estas propiedades son \emph{ind-construibles}, pero daremos más adelante un ejemplo en el que $S$ es una parte abierta no vacía de $\operatorname{Spec}(\mathbf{Z})$ (o una parte abierta de una curva elíptica sobre un cuerpo finito) y en el que todas las fibras $X_s$

\oldpage[IV-3]{58}
salvo la fibra genérica $X_\eta$ son proyectivas sobre $\mathbf{k}(s)$ (siendo todos los preesquemas $X_s$ de dimensión~2).
\end{enumerate}
\end{remarks}

\begin{proposition}[9.2.3]
\label{IV.9.2.3}
Sea $\mathbf{P}$ una propiedad construible \textnormal{(}resp.\ ind-construible\textnormal{)} de los preesquemas algebraicos, sea $S$ un preesquema, sea $X$ un preesquema de presentación finita sobre $S$ y sea $\sh{F}$ un $\sh{O}_X$-Módulo cuasicoherente de presentación finita.
Entonces el conjunto $E$ de los $s\in S$ tales que $\mathbf{P}(X_s,\sh{F}_s,\mathbf{k}(s))$ es cierta es localmente construible \textnormal{(}resp.\ ind-construible\textnormal{)}.
Además, si $S$ es irreducible de punto genérico $\eta$, uno de los dos conjuntos $E$, $S-E$ es un entorno de $\eta$ \textnormal{(}resp.\ $E$ es un entorno de $\eta$ si contiene este punto\textnormal{)}.
\end{proposition}

Para demostrar estas afirmaciones, podemos reducirnos al caso en que $S=\operatorname{Spec}(A)$ es afín.
Sabemos entonces que existe un subanillo $A_0$ de $A$ que es una $\mathbf{Z}$-álgebra de tipo finito, un $A_0$-preesquema de tipo finito $X_0$ y un $\sh{O}_{X_0}$-Módulo coherente $\sh{F}_0$ tales que $X$ es isomorfo a $X_0\otimes_{A_0}A$ y $\sh{F}$ a $\sh{F}_0\otimes_{A_0}A$~\sref{IV.8.9.1}.
Sea $p:S\to S_0=\operatorname{Spec}(A_0)$ el morfismo correspondiente a la inyección $A_0\to A$, y sea $E_0$ el conjunto de los $s_0\in S_0$ tales que
\[
  \mathbf{P}((X_0)_{s_0},(\sh{F}_0)_{s_0},\mathbf{k}(s_0))
\]
es cierta; entonces, en virtud de~\sref{IV.9.2.2}[iv], se tiene $E=p^{-1}(E_0)$; por~\sref{IV.1.8.2}, podemos por tanto reducirnos al caso en que $S$ es el espectro de una $\mathbf{Z}$-álgebra de tipo finito y, por tanto, un esquema noetheriano.
Utilicemos el criterio de construibilidad~($\mathbf{0}_{\mathrm{III}}$,~9.2.3) (resp.\ el criterio de ind-construibilidad~\sref{IV.1.9.10}); empleando~\sref{IV.9.2.2}[iv] como antes y sustituyendo $S$ por un subpreesquema cerrado íntegro de $S$, nos reducimos al caso en que $S$ es noetheriano e íntegro y en que hay que demostrar que $E$ no es denso en ninguna parte de $S$ o contiene un subconjunto abierto no vacío de $S$ (resp.\ que $E$ contiene un subconjunto abierto no vacío de $S$ si contiene el punto genérico); esto queda garantizado por la condición~\sref{IV.9.2.1}[ii].

Observemos que hemos utilizado~\sref{IV.9.2.1}[ii] únicamente cuando $S$ es el espectro de un anillo íntegro de tipo finito sobre $\mathbf{Z}$.
Por otra parte, está claro que el enunciado de~\sref{IV.9.2.3} se aplica también cuando aparecen en $\mathbf{P}$ varios preesquemas, Módulos sobre estos preesquemas, morfismos de preesquemas u homomorfismos de Módulos.
Se aplica asimismo cuando aparecen partes (en número finito) de los preesquemas considerados, siempre que se exija que estas partes sean \emph{localmente construibles}.
En efecto, la reducción al caso en que $S$ es afín muestra que podemos reducirnos al caso en que estas partes son \emph{construibles}: se aplica entonces~{\hyperlink{ega4.c04.thm.8-3-11}{(8.3.11)}}, que muestra (con la notación precedente) que una parte construible de $X$ es la imagen inversa de una parte construible de $X_0$ para una elección adecuada de $A_0$.

\begin{corollary}[9.2.4]
\label{IV.9.2.4}
Sea $\mathbf{P}$ una propiedad construible \textnormal{(}resp.\ ind-construible\textnormal{)} de los preesquemas algebraicos, sean $X$, $Y$ dos $S$-preesquemas de presentación finita y sea $f:X\to Y$ un $S$-morfismo.
Para todo $s\in S$, pongamos $X_s=X\times_S\operatorname{Spec}(\mathbf{k}(s))$, $Y_s=Y\times_S\operatorname{Spec}(\mathbf{k}(s))$, $f_s=f\times 1:X_s\to Y_s$.
Entonces el conjunto $E$ de los $s\in S$ tales que, para todo $y\in Y_s$, la propiedad $\mathbf{P}(f_s^{-1}(y),\mathbf{k}(y))$ es cierta, es localmente construible \textnormal{(}resp.\ ind-construible\textnormal{)}.
\end{corollary}

En efecto, sea $Z$ el conjunto de los $y\in Y$ tales que $\mathbf{P}(f^{-1}(y),\mathbf{k}(y))$ es cierta.
Como las fibras $f^{-1}(y)$ y $f_s^{-1}(y)$ son isomorfas, vemos que $E$ es el conjunto de los $s\in S$ tales que $Y_s\subset Z$; si $g:Y\to S$ es el morfismo estructural, se tiene por tanto $E=S-g(Y-Z)$.

\oldpage[IV-3]{59}

Ahora bien, $f$ es de presentación finita~\sref{IV.1.6.2}[v], luego se sigue de~\sref{IV.9.2.3} que $Z$ es localmente construible (resp.\ ind-construible) en $Y$, y por tanto $Y-Z$ es localmente construible (resp.\ pro-construible) en $Y$.
Como $g$ es de presentación finita, $g(Y-Z)$ es localmente construible (resp.\ pro-construible) en $S$, en virtud del teorema de Chevalley~\sref{IV.1.8.4} (resp.\ de~\sref{IV.1.9.5}[vii]); por tanto, $E$ es localmente construible (resp.\ ind-construible) en $S$.

\begin{remark}[9.2.5]
\label{IV.9.2.5}
Observemos que, si $\mathbf{P}$ es una propiedad de los preesquemas algebraicos para la cual es cierta la prop.~\sref{IV.9.2.3}, entonces $\mathbf{P}$ satisface también la condición~\sref{IV.9.2.1}[ii]: esto resulta, en efecto, de que, en un espacio noetheriano irreducible, un conjunto construible o bien no es denso en ninguna parte o bien contiene un abierto no vacío~($\mathbf{0}_{\mathrm{III}}$,~9.2.3).
\end{remark}

\begin{proposition}[9.2.6]
\label{IV.9.2.6}
Sea $\mathbf{P}$ una de las propiedades siguientes de un $k$-preesquema algebraico $X$:
\begin{enumerate}[label=(\roman*)]
  \item $X$ es vacío.
  \item $X$ es finito sobre $k$.
  \item $X$ es radicial sobre $k$.
  \item $\dim(X)$ pertenece a una parte dada $\Phi$ del conjunto $\overline{\mathbf{Z}}=\mathbf{Z}\cup\{-\infty\}$.
\end{enumerate}
Entonces $\mathbf{P}$ es construible.
\end{proposition}

Está claro que (i) y (ii) son casos particulares de (iv), tomando como $\Phi$, respectivamente, el conjunto $\{-\infty\}$ y el conjunto $\{-\infty,0\}$.
Por tanto, solo tenemos que demostrar (iii) y (iv).
En cada uno de estos dos casos, la condición~(i) de~\sref{IV.9.2.1} se satisface en virtud de~\sref{IV.2.7.1}[xv] y~\sref{IV.4.1.4}.
Por otra parte, en el caso (iii), la propiedad $\mathbf{P}$ satisface la conclusión de~\sref{IV.9.2.3} en virtud de~\sref{IV.1.8.7}; queda, pues, por comprobar que lo mismo ocurre en el caso (iv).
Esto se seguirá de la proposición más precisa que damos a continuación:

\begin{proposition}[9.2.6.1]
\label{IV.9.2.6.1}
Si $f:X\to S$ es un morfismo de presentación finita, la función $s\mapsto\dim(f^{-1}(s))$ es localmente construible.
\end{proposition}

La cuestión es local sobre $S$, por lo que podemos suponer que $S=\operatorname{Spec}(A)$ es afín y demostrar que, para todo $n$, el conjunto $E$ de los $s\in S$ tales que $\dim(X_s)=n$ es construible.
El mismo razonamiento que en~\sref{IV.9.2.3} nos reduce al caso en que $A$ es noetheriano e íntegro, y basta entonces demostrar el siguiente

\begin{corollary}[9.2.6.2]
\label{IV.9.2.6.2}
Sea $S$ un preesquema noetheriano íntegro de punto genérico $\eta$, y sea $f:X\to S$ un morfismo de tipo finito.
Entonces existe un entorno de $\eta$ en $S$ tal que la función $s\mapsto\dim(X_s)$ es constante en este entorno.
\end{corollary}

Las imágenes por $f$ de las componentes irreducibles (en número finito) de $X$ que no cortan a $X_\eta$ están contenidas en partes cerradas de $S$ que no contienen a $\eta$ (pues $S$ es íntegro~($\mathbf{0}_{\mathrm{I}}$,~2.1.5)); por tanto, tras sustituir $S$ por un entorno abierto de $\eta$, podemos reducirnos al caso en que todas las componentes irreducibles $X_i$ de $X$ cortan a $X_\eta$.
Denotemos de nuevo por $X_i$ el subpreesquema cerrado reducido de $X$ que tiene a $X_i$ como espacio subyacente; como $\dim(X_s)=\sup_i\dim((X_i)_s)$~\sref{IV.4.1.1}, podemos reducirnos al caso en que $X$ es irreducible.
Existe entonces un recubrimiento finito $(U_j)$ de $X$ por

\oldpage[IV-3]{60}
abiertos afines densos, y los números $\dim((U_j)_\eta)$ son todos iguales a $n=\dim(X_\eta)$~\sref{IV.4.1.1.3}; podemos, por tanto, reducirnos al caso en que $X$ es afín, y en consecuencia también $X_\eta$.
Existe entonces, en virtud de~\sref{IV.4.1.2}, un abierto no vacío $W$ de $X$ tal que $W\supset X_\eta$, y un $\mathbf{k}(\eta)$-morfismo finito y sobreyectivo $h:W_\eta\to\mathbf{V}_{\mathbf{k}(\eta)}^n(=\operatorname{Spec}(\mathbf{k}(\eta)[T_1,\ldots,T_n]))$; aplicando~\sref{IV.8.8.2}[i] y el método de~\sref{IV.8.1.2}[a], deducimos (tras sustituir $S$ por un entorno de $\eta$ si es necesario) que $h=g_\eta$, donde $g:W\to\mathbf{V}_A^n(=\operatorname{Spec}(A[T_1,\ldots,T_n]))$ es un $S$-morfismo, y podemos suponer que este morfismo es finito y sobreyectivo en virtud de~\sref{IV.8.10.5}[vi] y~\sref{IV.8.10.5}[x].
Concluimos, por tanto, que para todo $s\in S$, el morfismo $g_s:W_s\to\mathbf{V}_{\mathbf{k}(s)}^n$ es finito y sobreyectivo, y por consiguiente $\dim(W_s)=n$~\sref{IV.4.1.2}.

\subsection{Propiedades construibles de los morfismos de preesquemas algebraicos}
\label{subsection:IV.9.3}

\begin{proposition}[9.3.1]
\label{IV.9.3.1}
Sea $\mathbf{P}(X,k)$ una propiedad construible \textnormal{(}resp.\ ind-construible\textnormal{)} de los preesquemas algebraicos.
Denotemos por $\mathbf{P}'(f,X,Y,k)$ la relación siguiente: $f:X\to Y$ es un $k$-morfismo de $k$-preesquemas algebraicos tal que, para todo $y\in Y$, la propiedad $\mathbf{P}(f^{-1}(y),\mathbf{k}(y))$ es cierta.
Entonces $\mathbf{P}'$ es una propiedad construible \textnormal{(}resp.\ ind-construible\textnormal{)}.
\end{proposition}

En efecto, como $\mathbf{P}$ satisface la condición~\sref{IV.9.2.1}[i], lo mismo ocurre con $\mathbf{P}'$ en virtud de la transitividad de las fibras~(\textbf{I},~3.6.4); además, el hecho de que $\mathbf{P}'$ satisface la condición~\sref{IV.9.2.1}[ii] se sigue de~\sref{IV.9.2.4}, en razón de la observación~\sref{IV.9.2.5}.

\begin{proposition}[9.3.2]
\label{IV.9.3.2}
Sea $\mathbf{P}$ una de las propiedades siguientes de un $k$-morfismo $f:X\to Y$ de preesquemas algebraicos:
\begin{enumerate}[label=(\roman*)]
  \item $f$ es sobreyectivo.
  \item $f$ es cuasifinito.
  \item $f$ es radicial.
  \item Para todo $y\in Y$, $\dim(f^{-1}(y))$ pertenece a $\Phi$ \textnormal{(}notación de~\sref{IV.9.2.6}\textnormal{)}.
\end{enumerate}
Entonces $\mathbf{P}$ es construible.
\end{proposition}

Esto se sigue inmediatamente de~\sref{IV.9.3.1} y~\sref{IV.9.2.6}, teniendo en cuenta que $f$ es de tipo finito~\sref[I]{I.1.5.4}[v], la caracterización de los morfismos radiciales~\sref[I]{I.3.5.8} y la de los morfismos cuasifinitos~\sref[II]{II.6.2.2}.

\begin{proposition}[9.3.3]
\label{IV.9.3.3}
Supongamos satisfechas las hipótesis de~\sref{IV.8.8.1} y conservemos las mismas notaciones; supongamos además que $S_\alpha$ es cuasicompacto, que $X_\alpha$ e $Y_\alpha$ son de presentación finita sobre $S_\alpha$, y sea $f_\alpha:X_\alpha\to Y_\alpha$ un $S_\alpha$-morfismo.
Sea $\mathbf{P}$ una propiedad ind-construible de los morfismos de preesquemas algebraicos.
Para todo $s\in S$ \textnormal{(}resp.\ $s_\lambda\in S_\lambda$\textnormal{)}, pongamos $X_s=X\times_S\operatorname{Spec}(\mathbf{k}(s))$, $Y_s=Y\times_S\operatorname{Spec}(\mathbf{k}(s))$, $f_s=f\times 1:X_s\to Y_s$ \textnormal{(}resp.\ $X_{\lambda,s_\lambda}=X_\lambda\times_{S_\lambda}\operatorname{Spec}(\mathbf{k}(s_\lambda))$, $Y_{\lambda,s_\lambda}=Y_\lambda\times_{S_\lambda}\operatorname{Spec}(\mathbf{k}(s_\lambda))$, $f_{\lambda,s_\lambda}=f_\lambda\times 1:X_{\lambda,s_\lambda}\to Y_{\lambda,s_\lambda}$\textnormal{)}.
Entonces, para que para todo $s\in S$ la propiedad $\mathbf{P}(f_s)$ sea cierta, es necesario y suficiente que exista $\lambda\geqslant\alpha$ tal que para todo $s_\lambda\in S_\lambda$ se tenga $\mathbf{P}(f_{\lambda,s_\lambda})$.
\end{proposition}

En efecto, sea $E$ (resp.\ $E_\lambda$) el conjunto de los $s\in S$ (resp.\ $s_\lambda\in S_\lambda$) tales que la propiedad $\mathbf{P}(f_s)$ (resp.\ $\mathbf{P}(f_{\lambda,s_\lambda})$) es cierta; se sigue de~\sref{IV.9.2.2}[iv] que $E_\mu=u_{\lambda\mu}^{-1}(E_\lambda)$ para $\lambda\leqslant\mu$, y que $E=u_\lambda^{-1}(E_\lambda)$; además, en virtud de~\sref{IV.9.2.3}, $E$ (resp.\ $E_\lambda$) es ind-construible en $S$ (resp.\ $S_\lambda$); la proposición se sigue entonces de~\sref{IV.8.3.4} aplicado a la parte ind-construible $E$ de $S$.

\oldpage[IV-3]{61}

Este resultado se extiende sin dificultad a las propiedades $\mathbf{P}$ del tipo considerado en (9.2.3, (i) y (ii)).

\begin{remark}[9.3.4]
\label{IV.9.3.4}
La conjunción de~\sref{IV.9.3.3} y~\sref{IV.9.3.2}[ii] demuestra la afirmación~\sref{IV.8.10.5}[xi].
\end{remark}

\begin{proposition}[9.3.5]
\label{IV.9.3.5}
Sea $\mathbf{P}(X,Y,k)$ la propiedad: ``$X$ e $Y$ son dos preesquemas de tipo finito sobre el cuerpo $k$, y existe una extensión $k'$ de $k$ y un $k'$-morfismo $g:X_{(k')}\to Y_{(k')}$ que satisface $\mathbf{Q}(g)$'', donde $\mathbf{Q}$ es una de las propiedades~\emph{(i)} a~\emph{(xiv)} de~\sref{IV.8.10.5}.
Entonces $\mathbf{P}$ es una propiedad ind-construible.
\end{proposition}

La definición de $\mathbf{P}$ muestra, en efecto, que se satisface la condición~\sref{IV.9.2.1}[i], teniendo en cuenta que la propiedad $\mathbf{Q}(g)$ es estable por cambio de base y que dos extensiones de $k$ siempre pueden considerarse como subextensiones de una tercera extensión de $k$.
Para verificar~\sref{IV.9.2.1}[ii], podemos reducirnos al caso en que $S=\operatorname{Spec}(A)$ es afín; si $K=\mathbf{k}(\eta)$ es el cuerpo de fracciones de $A$, existen, por hipótesis y en virtud de~\sref{IV.9.1.4}, una extensión finita $K'$ de $K$ y un $K'$-morfismo $g':(X_\eta)_{(K')}\to(Y_\eta)_{(K')}$ que satisface $\mathbf{Q}(g')$, y $K'$ es evidentemente el cuerpo de fracciones de una $A$-álgebra íntegra finita $A'$.
Si ponemos $S'=\operatorname{Spec}(A')$, se sigue entonces de~\sref{IV.8.10.5} que existe un entorno $U'$ del punto genérico $\eta'$ de $S'$ tal que, si ponemos $X'=X\otimes_A A'$ e $Y'=Y\otimes_A A'$, para todo $s'\in U'$ existe un morfismo $X'_{s'}\to Y'_{s'}$ que posee la propiedad $\mathbf{Q}$.
Pero el morfismo $h:S'\to S$ es finito y, por tanto, cerrado; como $h^{-1}(\eta)=\{\eta'\}$, $h(U')$ contiene un entorno abierto $U$ de $\eta$ en $S$.
Para todo $s\in U$ existe, por tanto, $s'\in U'$ tal que $h(s')=s$, y como $X'_{s'}=X_s\otimes_{\mathbf{k}(s)}\mathbf{k}(s')$ e $Y'_{s'}=Y_s\otimes_{\mathbf{k}(s)}\mathbf{k}(s')$, la propiedad $\mathbf{P}(X_s,Y_s,\mathbf{k}(s))$ es cierta para todo $s\in U$.

\begin{env}[9.3.6]
\label{IV.9.3.6}
\textbf{Ejemplo.}
Tomemos, por ejemplo, como $\mathbf{Q}$ la propiedad de ser un isomorfismo.
Combinando entonces~\sref{IV.9.3.5} y~\sref{IV.9.3.3}, obtenemos la propiedad siguiente: con las notaciones e hipótesis de~\sref{IV.8.8.1}, siendo $S_\alpha$ cuasicompacto y $X_\alpha$ e $Y_\alpha$ de presentación finita sobre $S_\alpha$, para que, para todo $s\in S$, $X_s$ e $Y_s$ sean \emph{geométricamente isomorfos}~\sref{IV.9.1.4}, es necesario y suficiente que exista $\lambda\geqslant\alpha$ tal que, para todo $s_\lambda\in S_\lambda$, $X_{\lambda,s_\lambda}$ e $Y_{\lambda,s_\lambda}$ sean geométricamente isomorfos.
\end{env}

Tenemos un resultado análogo cuando los preesquemas considerados están dotados de ``leyes de composición'' de cierto tipo~($\mathbf{0}_{\mathrm{III}}$,~8.2.1), por ejemplo ``preesquemas en grupos'', ``preesquemas en anillos'', etc.~($\mathbf{0}_{\mathrm{III}}$,~8.2.3).
El enunciado precedente sigue siendo válido cuando por ``isomorfismo'' se entienden isomorfismos de preesquemas que son homomorfismos para las leyes de composición consideradas~($\mathbf{0}_{\mathrm{III}}$,~8.2.2); basta utilizar aquí no solo~\sref{IV.8.10.5}, sino también~\sref{IV.8.8.2}[i], observando que la noción de homomorfismo para una ley de composición se expresa diciendo que ciertos diagramas de morfismos de preesquemas son conmutativos (se sobreentiende que los morfismos de transición $X_\mu\to X_\lambda$ e $Y_\mu\to Y_\lambda$ para $\lambda\leqslant\mu$ son homomorfismos para las leyes de composición consideradas).

También se pueden considerar, en vez de morfismos de preesquemas como en~\sref{IV.9.3.5}, homomorfismos de Módulos, utilizando~\sref{IV.9.1.6} en lugar de~\sref{IV.9.1.5}.

\oldpage[IV-3]{62}

\subsection{Construibilidad de ciertas propiedades de los módulos}
\label{subsection:IV.9.4}

\begin{env}[9.4.1]
\label{IV.9.4.1}
\textbf{Notaciones (9.4.1).}
En este n\textsuperscript{o} y en los siguientes hasta el final de \textsection9, utilizaremos sistemáticamente las notaciones siguientes: dado un morfismo $f:X\to S$, pondremos, para todo $s\in S$, $X_s=f^{-1}(s)=X\times_S\operatorname{Spec}(\mathbf{k}(s))$; para todo $\sh{O}_X$-Módulo cuasicoherente $\sh{F}$, $\sh{F}_s$ denotará el $\sh{O}_{X_s}$-Módulo cuasicoherente $\sh{F}\otimes_{\sh{O}_S}\mathbf{k}(s)$, y para todo homomorfismo $u:\sh{F}\to\sh{G}$ de $\sh{F}$ en un $\sh{O}_X$-Módulo cuasicoherente $\sh{G}$, $u_s:\sh{F}_s\to\sh{G}_s$ será el morfismo $p^*(u)$, designando por $p$ la proyección canónica $X_s\to X$.
Para toda sección $g$ de $\sh{F}$ sobre $X$, denotaremos por $g_s$ la imagen de $g$ por el homomorfismo canónico $\Gamma(X,\sh{F})\to\Gamma(X_s,\sh{F}_s)$.
Para toda parte $Z$ de $X$, denotaremos por $Z_s$ la imagen inversa $p^{-1}(Z)=Z\cap X_s$~(\textbf{I},~3.6.1).
Finalmente, si $Y$ es un segundo $S$-preesquema y $h:X\to Y$ un $S$-morfismo, denotaremos por $h_s$ el morfismo $h\times 1:X_s\to Y_s$.
\end{env}

\begin{proposition}[9.4.2]
\label{IV.9.4.2}
Sea $S$ un preesquema íntegro de punto genérico $\eta$, $f:X\to S$ un morfismo de presentación finita, y sean $\sh{F}$, $\sh{G}$, $\sh{H}$ tres $\sh{O}_X$-Módulos cuasicoherentes de presentación finita.
Sean $u:\sh{F}\to\sh{G}$, $v:\sh{G}\to\sh{H}$ dos homomorfismos de $\sh{O}_X$-Módulos, y supongamos que la sucesión $\sh{F}_\eta\xrightarrow{u_\eta}\sh{G}_\eta\xrightarrow{v_\eta}\sh{H}_\eta$ es exacta.
Entonces existe un entorno $U$ de $\eta$ en $S$ tal que, para todo $s\in U$, la sucesión $\sh{F}_s\xrightarrow{u_s}\sh{G}_s\xrightarrow{v_s}\sh{H}_s$ es exacta.
\end{proposition}

\begin{proof}
Con las notaciones generales de~\sref{IV.9.2.1}, se trata aquí de la relación $\mathbf{P}(X,\sh{F},\sh{G},\sh{H},u,v,k)$: ``$X$ es un preesquema algebraico sobre el cuerpo $k$, $\sh{F}\xrightarrow{u}\sh{G}\xrightarrow{v}\sh{H}$ es una sucesión exacta de $\sh{O}_X$-Módulos cuasicoherentes''.
Como, para toda extensión $k'$ de $k$, la proyección canónica $X_{(k')}\to X$ es un morfismo fielmente plano~\sref{IV.2.2.13}, se satisface la condición~\sref{IV.9.2.1},~(i))~\sref{IV.2.2.7}.
En virtud de~\sref{IV.9.2.3}, podemos restringirnos al caso en que $S$ es íntegro y noetheriano, en cuyo caso $X$ es noetheriano, y $\sh{F}$, $\sh{G}$, $\sh{H}$ son $\sh{O}_X$-Módulos coherentes.
La hipótesis implica que existe un entorno $U$ de $\eta$ en $S$ tal que $\sh{F}|f^{-1}(U)\to\sh{G}|f^{-1}(U)\to\sh{H}|f^{-1}(U)$ es \emph{exacta}, en virtud de~\sref{IV.8.5.8},~(ii)) aplicado según el método general de~\sref{IV.8.1.2},~$a$)), y por tanto podemos suponer ya que $\sh{F}\xrightarrow{u}\sh{G}\xrightarrow{v}\sh{H}$ es exacta; basta evidentemente demostrar que $\operatorname{Ker}(v_s)=(\operatorname{Ker} v)_s$ y $\operatorname{Im}(u_s)=(\operatorname{Im} u)_s$ para $s$ próximo a $\eta$ en $S$.
Teniendo en cuenta la coherencia de $\sh{F}$, $\sh{G}$ y $\sh{H}$ y de
($\mathbf{0}_{\mathrm{I}}$,~5.3.4), nos reducimos al caso particular en que
\[
 0\longrightarrow\sh{F}\longrightarrow\sh{G}\longrightarrow\sh{H}
 \longrightarrow0
\]
es exacta.
Existe entonces un entorno abierto $U$ de $\eta$ en $S$ tal que
$\sh{H}|f^{-1}(U)$ es plano sobre $U$ \sref{IV.6.9.1}; por tanto, \sref{IV.2.1.8}
muestra que para todo $s\in U$ la sucesión
\[
 0\longrightarrow\sh{F}_s\longrightarrow\sh{G}_s\longrightarrow\sh{H}_s
 \longrightarrow0
\]
es exacta.
\end{proof}

\begin{corollary}[9.4.3]
\label{IV.9.4.3}
Sea $S$ un preesquema íntegro, de punto genérico $\eta$, y sea $f:X\to S$ un morfismo de presentación finita.
Sea $\sh{L}^\bullet=(\sh{L}^i)_{i\in\mathbf{Z}}$ un complejo cuyos términos son $\sh{O}_X$-Módulos cuasicoherentes de presentación finita.
Para todo $i$, existe un entorno abierto $U$ de $\eta$ en $S$ tal que los homomorfismos canónicos
\begin{equation}
\label{IV.9.4.3.1}
  (\sh{H}^i(\sh{L}^\bullet))_s\to\sh{H}^i(\sh{L}^\bullet_s) \tag{9.4.3.1}
\end{equation}
son biyectivos para todo $s\in U$.
\end{corollary}

\oldpage[IV-3]{63}

\begin{proof}
Es suficiente considerar un complejo de tres términos en grados $-1$, $0$, $+1$,
\[
 \sh{M}\xrightarrow{u}\sh{N}\xrightarrow{v}\sh{P},
\]
y $i=0$.
El homomorfismo canónico en cuestión es entonces
\[
 (\operatorname{Ker}v/\operatorname{Im}u)_s
 \longrightarrow \operatorname{Ker}(v_s)/\operatorname{Im}(u_s).
\]
Utilizando \sref{IV.8.9.1} y \sref{IV.8.5.2},~(i), podemos reducirnos al caso
en que $S$ (y por tanto $X$) es noetheriano; los Módulos $\sh{M}$, $\sh{N}$ y
$\sh{P}$ son entonces coherentes.
Así, $\operatorname{Im}u$ y $\operatorname{Ker}v$ son coherentes
($\mathbf{0}_{\mathrm{I}}$,~5.3.4), y \sref{IV.9.4.2} proporciona un entorno
$U$ de $\eta$ en el cual
\[
 \operatorname{Ker}(v_s)=(\operatorname{Ker}v)_s,
 \qquad
 \operatorname{Im}(u_s)=(\operatorname{Im}u)_s.
\]
La conclusión se sigue aplicando \sref{IV.9.4.2} a la sucesión exacta
\[
 0\longrightarrow\operatorname{Im}u\longrightarrow\operatorname{Ker}v
 \longrightarrow\operatorname{Ker}v/\operatorname{Im}u\longrightarrow0,
\]
puesto que $\sh{O}_\eta=\boldsymbol{k}(\eta)$ y, por consiguiente, la sucesión
correspondiente sobre $\eta$ es exacta.
\end{proof}

\begin{proposition}[9.4.4]
\label{IV.9.4.4}
Sea $f:X\to S$ un morfismo de presentación finita, y sean $\sh{F}$, $\sh{G}$, $\sh{H}$ tres $\sh{O}_X$-Módulos cuasicoherentes de presentación finita.
Sean $u:\sh{F}\to\sh{G}$, $v:\sh{G}\to\sh{H}$ dos homomorfismos de $\sh{O}_X$-Módulos.
Entonces el conjunto $E$ de los $s\in S$ tales que la sucesión $\sh{F}_s\xrightarrow{u_s}\sh{G}_s\xrightarrow{v_s}\sh{H}_s$ es exacta es localmente construible.
\end{proposition}

Teniendo en cuenta~\sref{IV.9.2.3}, basta establecer que la propiedad $\mathbf{P}$ considerada en~\sref{IV.9.4.2} es construible.
Ya hemos observado en la demostración de~\sref{IV.9.4.2} que la condición~\sref{IV.9.2.1},~(i)) se satisface para esta propiedad, y queda verificar la condición~\sref{IV.9.2.1},~(ii)).
Supongamos, pues, que $S$ es íntegro noetheriano, de punto genérico $\eta$, y demostremos que $E$ o $S-E$ es un entorno de $\eta$.
Si $\eta\in E$, nuestra afirmación se sigue de~\sref{IV.9.4.2}, y por tanto podemos restringirnos al caso en que $\eta\notin E$, es decir, la sucesión $\sh{F}_\eta\xrightarrow{u_\eta}\sh{G}_\eta\xrightarrow{v_\eta}\sh{H}_\eta$ no es exacta.
Distinguimos entonces dos casos:

$1^\circ$ Pongamos $w=v\circ u$, y supongamos primero que $w_\eta=v_\eta\circ u_\eta\neq 0$.
Como $\sh{F}$, $\sh{G}$, $\sh{H}$ son coherentes, también lo es $\sh{N}=\operatorname{Ker}(w)$~($\mathbf{0}_{\mathrm{I}}$,~5.3.4); se sigue entonces de~\sref{IV.9.4.2}, aplicado a la sucesión exacta $0\to\sh{N}\to\sh{F}\xrightarrow{w}\sh{H}$, que existe un entorno $U$ de $\eta$ en $S$ tal que $\operatorname{Ker}(w_s)=\sh{N}_s$; restringiendo $S$, podemos por tanto suponer que esta relación se satisface para todo $s\in S$.
Sea $j$ la inyección canónica $\sh{N}\to\sh{F}$, y pongamos $\sh{M}=\operatorname{Coker}(j)$; la exactitud por la derecha del funtor $\sh{F}\mapsto\sh{F}_s$ implica que $\sh{M}_s=\operatorname{Coker}(j_s)$ para todo $s\in S$.
La hipótesis $w_\eta\neq 0$ significa que $\sh{M}_\eta\neq 0$; como $\sh{M}$ es coherente~($\mathbf{0}_{\mathrm{I}}$,~5.3.4), se sigue de~\sref{IV.1.8.6} que existe un entorno $U$ de $\eta$ en $S$ tal que $\sh{M}_s\neq 0$ para todo $s\in U$, de donde $w_s\neq 0$ para todo $s\in U$, y \emph{a fortiori} $S-E$ es un entorno de $\eta$.

$2^\circ$ Supongamos que $w_\eta=0$; en virtud de~\sref{IV.8.5.2},~(i)), aplicado según el método general de~\sref{IV.8.1.2},~$a$)), existe un entorno $U$ de $\eta$ tal que $w|f^{-1}(U)=0$; sustituyendo $S$ por $U$, podemos suponer ya que $w=0$ en $X$.
Entonces $\sh{F}\xrightarrow{u}\sh{G}\xrightarrow{v}\sh{H}$ es un complejo de tres términos, al que podemos aplicar~\sref{IV.9.4.3}; por hipótesis tenemos $(\sh{H}^0(\sh{L}^\bullet))_\eta=\sh{H}^0(\sh{L}^\bullet_\eta)\neq 0$, y $\sh{H}^0(\sh{L}^\bullet)$ es coherente~($\mathbf{0}_{\mathrm{I}}$,~5.3.4 y~5.3.3), por tanto se sigue de~\sref{IV.1.8.6} que existe un entorno $U$ de $\eta$ tal que $(\sh{H}^0(\sh{L}^\bullet))_s\neq 0$ para todo $s\in U$; pero como podemos suponer que $\sh{H}^0(\sh{L}^\bullet_s)=(\sh{H}^0(\sh{L}^\bullet))_s$ para todo $s\in U$ por~\sref{IV.9.4.3}, vemos de nuevo que $S-E$ es un entorno de $\eta$ en $S$.

\oldpage[IV-3]{64}

\begin{corollary}[9.4.5]
\label{IV.9.4.5}
Sea $f:X\to S$ un morfismo de presentación finita, sean $\sh{F}$, $\sh{G}$ dos $\sh{O}_X$-Módulos cuasicoherentes de presentación finita, y sea $u:\sh{F}\to\sh{G}$ un homomorfismo de $\sh{O}_X$-Módulos.
Entonces el conjunto de los $s\in S$ tales que $u_s$ es inyectivo \textnormal{(}resp.\ sobreyectivo, biyectivo, nulo\textnormal{)} es localmente construible.
\end{corollary}

Basta aplicar~\sref{IV.9.4.4} a las sucesiones $0\to\sh{F}\xrightarrow{u}\sh{G}$, $\sh{F}\xrightarrow{u}\sh{G}\to 0$, $0\to\sh{F}\xrightarrow{u}\sh{G}\to 0$, $\sh{F}\xrightarrow{u}\sh{G}\xrightarrow{1_\sh{G}}\sh{G}$.

\begin{corollary}[9.4.6]
\label{IV.9.4.6}
Sea $f:X\to S$ un morfismo de presentación finita y sea $\sh{F}$ un $\sh{O}_X$-Módulo cuasicoherente de presentación finita.
Sea $h$ una sección de $\sh{F}$ sobre $X$; para todo $s\in S$, sea $h_s$ la sección correspondiente de $\sh{F}_s$ sobre $X_s$ \textnormal{(}para el morfismo de proyección $X_s\to X$\textnormal{)}.
Entonces el conjunto de los $s\in S$ tales que $h_s=0$ es localmente construible.
\end{corollary}

Basta observar que $h$ corresponde a un homomorfismo $u:\sh{O}_X\to\sh{F}$~($\mathbf{0}_{\mathrm{I}}$,~5.1.1) y $h_s$ al homomorfismo $u_s$.

\begin{proposition}[9.4.7]
\label{IV.9.4.7}
Sea $f:X\to S$ un morfismo de presentación finita y sea $\sh{F}$ un $\sh{O}_X$-Módulo cuasicoherente de presentación finita.
El conjunto $E$ \textnormal{(}resp.\ $E'$\textnormal{)} de los $s\in S$ tales que $\sh{F}_s$ es un $\sh{O}_{X_s}$-Módulo localmente libre \textnormal{(}resp.\ localmente libre de rango $n$\textnormal{)} es localmente construible.
\end{proposition}

Si $X$ es un preesquema algebraico sobre un cuerpo $k$, $\sh{F}$ un $\sh{O}_X$-Módulo coherente y $k'$ una extensión de $k$, para que $\sh{F}$ sea localmente libre (resp.\ localmente libre de rango $n$), es necesario y suficiente que lo sea $\sh{F}\otimes_k k'$, puesto que la proyección $X_{(k')}\to X$ es un morfismo fielmente plano~\sref{IV.2.2.7}.
En otras palabras, la condición~\sref{IV.9.2.1},~(i)) se satisface para las propiedades cuya construibilidad queremos demostrar, y queda verificar~\sref{IV.9.2.1},~(ii)); podemos, por tanto, suponer aún que $S$ es afín, noetheriano e íntegro.
Hay entonces cuatro casos que considerar:

$1^\circ$ $\eta\in E$.
En virtud de~\sref{IV.8.5.5}, aplicado según el método general de~\sref{IV.8.1.2},~$a$)), existe un entorno $U$ de $\eta$ en $S$ tal que $\sh{F}|f^{-1}(U)$ es localmente libre; \emph{a fortiori} $\sh{F}_s$ es localmente libre para todo $s\in U$.

$2^\circ$ $\eta\in E'$.
El mismo razonamiento que en $1^\circ$.

$3^\circ$ $\eta\in S-E$.
Como $\sh{F}_\eta$ es un $\sh{O}_{X_\eta}$-Módulo coherente, decir que no es localmente libre significa que no es plano sobre $\sh{O}_{X_\eta}$ (Bourbaki, \emph{Alg.\ comm.}, cap.~II, \textsection5, n\textsuperscript{o}~2, cor.~2 del teorema~1).
El hecho de que $S-E$ sea un entorno de $\eta$ se seguirá entonces del lema más general siguiente:

\begin{lemma}[9.4.7.1]
\label{IV.9.4.7.1}
Sea $S$ un preesquema íntegro noetheriano, de punto genérico $\eta$, sean $X$, $Y$ dos $S$-preesquemas de tipo finito sobre $S$, $g:X\to Y$ un $S$-morfismo y $\sh{F}$ un $\sh{O}_X$-Módulo coherente.
Si $\sh{F}_\eta$ no es $g_\eta$-plano, entonces existe un entorno $U$ de $\eta$ en $S$ tal que, para todo $s\in U$, $\sh{F}_s$ no es $g_s$-plano.
\end{lemma}

Teniendo en cuenta~\sref{IV.2.1.2} y Bourbaki, \emph{Alg.\ comm.}, cap.~I\textsuperscript{er}, \textsection2, n\textsuperscript{o}~3, Remarque~1, la hipótesis significa que existe un abierto no vacío $V$ de $Y_\eta$, y un homomorfismo inyectivo $v:\sh{M}\to\sh{N}$ de $\sh{O}_V$-Módulos \emph{coherentes}, tal que el homomorfismo $1\otimes v:\sh{F}_\eta\otimes_{\sh{O}_V}\sh{M}\to\sh{F}_\eta\otimes_{\sh{O}_V}\sh{N}$ no es inyectivo.
Tenemos $V=Y_\eta\cap W$, donde $W$ es abierto en $Y$~(\textbf{I},~3.6.1), y se sigue de~\sref{IV.8.5.2},~(i) y~(ii)), aplicado según el método de~\sref{IV.8.1.2},~$a$)), que existe un entorno $U_0$ de $\eta$ en $S$, dos $\sh{O}_Z$-Módulos
\oldpage[IV-3]{65}
coherentes $\sh{G}$, $\sh{H}$ (donde $Z=W\cap h^{-1}(U_0)$, siendo $h:Y\to S$ el morfismo estructural) y un $\sh{O}_Z$-homomorfismo $u:\sh{G}\to\sh{H}$ tales que $\sh{M}=\sh{G}_\eta$, $\sh{N}=\sh{H}_\eta$, y $v=u_\eta$; podemos por tanto tomar $U_0$ de modo que $u_s:\sh{G}_s\to\sh{H}_s$ sea inyectivo para todo $s\in U_0$~\sref{IV.9.4.5}.
Pero para todo $s\in U_0$, el homomorfismo $1\otimes u_s:\sh{F}_s\otimes_{\sh{O}_{Z_s}}\sh{G}_s\to\sh{F}_s\otimes_{\sh{O}_{Z_s}}\sh{H}_s$ no es otro que $(1\otimes u)_s$; la hipótesis de que $(1\otimes u)_\eta$ no es inyectivo implica, pues, por~\sref{IV.9.4.5}, la existencia de un abierto no vacío $U\subset U_0$ tal que, para todo $s\in U$, $(1\otimes u)_s$ no es inyectivo, y por tanto $\sh{F}_s$ no es $g_s$-plano para todo $s\in U$.

$4^\circ$ $\eta\in S-E'$.
Está claro que $S-E\subset S-E'$, y si $\eta\in S-E$, $S-E'$ es \emph{a fortiori} un entorno de $\eta$ por $3^\circ$.
Supongamos, pues, que $\eta\in E$, de donde $\sh{F}_\eta$ es localmente libre; estas hipótesis implican que $X_\eta$ no es conexo, y que los rangos del $\sh{O}_{X_\eta}$-Módulo localmente libre $\sh{F}_\eta$ no son los mismos sobre las diversas componentes conexas de $X_\eta$.
Ahora bien, se sigue de \sref{IV.8.4.2}, aplicado según el método de \sref{IV.8.1.2},~$a)$, que podemos suponer (tras sustituir $S$ por un entorno de $\eta$) que $X$ y $X_\eta$ tienen el mismo número de componentes conexas, siendo las componentes conexas de $X_\eta$ las intersecciones de $X_\eta$ con las componentes conexas de $X$.
La conclusión se sigue entonces del razonamiento dado en $2^\circ$, aplicado a cada una de las componentes conexas de $X$ (que son en número finito).

\begin{remark}[9.4.7.2]
\label{IV.9.4.7.2}
El lema~\sref{IV.9.4.7.1} será generalizado más adelante y liberado de las hipótesis noetherianas~\sref{IV.11.2.8}.
\end{remark}

\begin{proposition}[9.4.8]
\label{IV.9.4.8}
Sea $S$ un preesquema localmente noetheriano, $f:X\to S$ un morfismo de tipo finito y $\sh{F}$ un $\sh{O}_X$-Módulo coherente.
Supongamos que, para todo $s\in S$, $X_s$ es un preesquema localmente íntegro.
Entonces el conjunto $E$ \textnormal{(}resp.\ $E'$\textnormal{)} de los $s\in S$ tales que $\sh{F}_s$ es un $\sh{O}_{X_s}$-Módulo de torsión \textnormal{(}resp.\ un $\sh{O}_{X_s}$-Módulo sin torsión\textnormal{)} es localmente construible.
\end{proposition}

\begin{proof}
Podemos suponer evidentemente que $S=\operatorname{Spec}(A)$ es afín y noetheriano y demostrar que $E$ (resp.\ $E'$) es construible utilizando el criterio~($\mathbf{0}_{\mathrm{III}}$,~9.2.3); sustituyendo $S$ por el subpreesquema cerrado reducido de $S$ cuyo espacio subyacente es una parte cerrada irreducible $Y$ de $S$, y observando~(\textbf{I},~3.6.4) que para $s\in Y$, la fibra $(X_{(Y)})_s$ se identifica canónicamente con $X_s$ y el haz $(\sh{F}\otimes_{\sh{O}_S}\sh{O}_Y)_s$ con $\sh{F}_s$, vemos que nos reducimos al caso en que $S$ es íntegro de punto genérico $\eta$, y a demostrar que $E$ o $S-E$ (resp.\ $E'$ o $S-E'$) es un entorno de $\eta$ en $S$.
Además, $X$ es una unión finita de abiertos afines $U_i$ de tipo finito sobre $S$, y cada $(U_i)_s$ es un subpreesquema abierto de $X_s$, por tanto localmente íntegro. Si $(U_i)_\eta$ es vacío (resp.\ no vacío), entonces $(U_i)_s$ es también vacío (resp.\ no vacío) en un entorno de $\eta$ \sref{IV.9.2.6}. Podemos, por tanto, suponer que todos los $(U_i)_\eta$ son no vacíos e íntegros; entonces $\sh{F}_s$ es de torsión (resp.\ sin torsión) si y solo si cada restricción $\sh{F}_s|_{(U_i)_s}$ es de torsión (resp.\ sin torsión).

Nos reducimos así al caso en que $X=\operatorname{Spec}(B)$ es afín y $\sh{F}=\widetilde{M}$, siendo $M$ un $B$-módulo finito. Pongamos $B_s=B\otimes_A\mathbf{k}(s)$ y $M_s=M\otimes_A\mathbf{k}(s)$, y supongamos $B_\eta$ íntegro. Hay cuatro casos.

$1^\circ$ Supongamos $\eta\in E$. Entonces $M_\eta$ es un $B_\eta$-módulo finito de torsión, por lo que existe un $h\neq 0$ en $B_\eta$ tal que $hM_\eta=0$. Por \sref{IV.8.5.2},~(i), aplicado como en \sref{IV.8.1.2},~$a)$,

\oldpage[IV-3]{66}

tras restringir $S$ podemos escribir $h=g_\eta$ para cierto $g\in\Gamma(X,\sh{O}_X)$. Sea $u:\sh{F}\to\sh{F}$ la multiplicación por $g$. Como $u_\eta=0$, \sref{IV.9.4.5} muestra que $u_s=0$ cerca de $\eta$. Por otra parte, la multiplicación por $g_\eta=h$ es inyectiva sobre $\sh{O}_{X_\eta}$ porque $B_\eta$ es íntegro y $h\neq 0$; de nuevo por \sref{IV.9.4.5}, la multiplicación por $g_s$ sobre $\sh{O}_{X_s}$ es inyectiva cerca de $\eta$. Así, $g_s$ es $\sh{O}_{X_s}$-regular y anula $\sh{F}_s$, por lo que $\sh{F}_s$ es de torsión cerca de $\eta$.

$2^\circ$ Supongamos $\eta\in S-E$. Si $T$ es el submódulo de torsión de $M_\eta$, entonces $M_\eta/T\neq 0$ y, siendo un $B_\eta$-módulo finito sin torsión, se sumerge en $B_\eta^n$ para cierto $n$. Existe, por tanto, un homomorfismo no nulo $w:M_\eta\to B_\eta$. Por \sref{IV.8.5.2},~(i), tras restringir $S$ se extiende a un homomorfismo $v:\sh{F}\to\sh{O}_X$. Como $v_\eta\neq 0$, \sref{IV.9.4.5} da $v_s\neq 0$ cerca de $\eta$; como $X_s$ es localmente íntegro, $\sh{F}_s$ no es de torsión para estos $s$.

$3^\circ$ Supongamos $\eta\in E'$. Como $M_\eta$ es un $B_\eta$-módulo finito sin torsión, existe un homomorfismo inyectivo $w:M_\eta\to B_\eta^n$. Utilizando \sref{IV.8.5.2},~(i), y \sref{IV.9.4.5} como en $2^\circ$, tras restringir $S$ obtenemos un homomorfismo $v:\sh{F}\to\sh{O}_X^n$ tal que $v_\eta=w$ y $v_s$ es inyectivo cerca de $\eta$. Así, $\sh{F}_s$ no tiene torsión cerca de $\eta$.

$4^\circ$ Supongamos $\eta\in S-E'$. Sea $T$ el submódulo de torsión de $M_\eta$. Entonces $T\neq 0$ y es finito. Por \sref{IV.8.5.2},~(i) y (ii), tras restringir $S$ existen un $\sh{O}_X$-Módulo coherente $\sh{G}$ y un homomorfismo inyectivo $u:\sh{G}\to\sh{F}$ tales que $\sh{G}_\eta=\widetilde{T}$ y $u_\eta$ es la inclusión canónica. Por $1^\circ$ y \sref{IV.1.8.6}, cerca de $\eta$ el módulo $\sh{G}_s$ es un $\sh{O}_{X_s}$-Módulo de torsión no nulo; por \sref{IV.9.4.5}, $u_s$ es inyectivo cerca de $\eta$. Por tanto, el submódulo de torsión de $\sh{F}_s$ es no nulo cerca de $\eta$.
\end{proof}

\begin{remark}[9.4.9]
\label{IV.9.4.9}
La propiedad ``$X$ es un $k$-preesquema algebraico localmente íntegro'' no satisface la condición \sref{IV.9.2.1},~(i), y por tanto no es seguro que el enunciado de \sref{IV.9.4.8} siga siendo válido cuando no se impone ninguna hipótesis sobre $S$ y se supone únicamente que $f$ es un morfismo de presentación finita y $\sh{F}$ es un $\sh{O}_X$-Módulo de presentación finita. Consideremos, sin embargo, el caso particular siguiente. Sea $S_0$ un preesquema localmente noetheriano, sea $f_0:X_0\to S_0$ un morfismo de tipo finito cuyas fibras $(X_0)_{s_0}$ son localmente íntegras para todo $s_0\in S_0$, y sea $\sh{F}_0$ un $\sh{O}_{X_0}$-Módulo coherente. Sea $g:S\to S_0$ un morfismo cualquiera, pongamos
$X=X_0\times_{S_0}S$, $\sh{F}=\sh{F}_0\otimes_{\sh{O}_{S_0}}\sh{O}_S$, y $f=f_0\times 1_S:X\to S$, y supongamos que, para todo $s\in S$, la fibra $X_s$ es también localmente íntegra. Entonces el conjunto $E$ (resp. $E'$) de los $s\in S$ tales que $\sh{F}_s$ es de torsión (resp. sin torsión)

\oldpage[IV-3]{67}

es también localmente construible.

En efecto, sea $s\in S$ y pongamos $s_0=g(s)$. Teniendo en cuenta \sref[I]{I.1.8.2}, basta demostrar que $\sh{F}_s$ es de torsión (resp. sin torsión) si y solo si $(\sh{F}_0)_{s_0}$ lo es. Ahora bien, $\operatorname{Supp}(\sh{F}_s)$ es la imagen inversa de $\operatorname{Supp}((\sh{F}_0)_{s_0})$ por la proyección $p:X_s\to(X_0)_{s_0}$ \sref[I]{I.9.1.13}. Como $p$ es fielmente plano y cuasicompacto, decir que $\operatorname{Supp}(\sh{F}_s)$ contiene un punto maximal de $X_s$ equivale a decir que $\operatorname{Supp}((\sh{F}_0)_{s_0})$ contiene un punto maximal de $(X_0)_{s_0}$ \sref[I]{I.1.1.5} y \sref{IV.2.3.4}; esto demuestra la afirmación relativa a $E$ \sref[I]{I.7.4.6}. Si el submódulo de torsión $\sh{T}$ de $(\sh{F}_0)_{s_0}$ no es nulo, entonces $\sh{T}\otimes_{\boldsymbol{k}(s_0)}\boldsymbol{k}(s)$, que es de torsión por lo que precede, no es nulo y se identifica con un submódulo de $\sh{F}_s$ \sref{IV.2.2.7}; por tanto, el submódulo de torsión de $\sh{F}_s$ no es nulo. Finalmente, si $(\sh{F}_0)_{s_0}$ no tiene torsión, entonces, tras restringirnos a un abierto afín de $(X_0)_{s_0}$, podemos suponer que es isomorfo a un submódulo de algún $\sh{O}_{(X_0)_{s_0}}^n$; por consiguiente, $\sh{F}_s$ es isomorfo a un submódulo de $\sh{O}_{X_s}^n$ \sref{IV.2.2.7}, lo que demuestra la afirmación relativa a $E'$.
\end{remark}

\subsection{Construibilidad de propiedades topológicas}
\label{subsection:IV.9.5}

\begin{proposition}[9.5.1]
\label{IV.9.5.1}
Sea $f:X\to S$ un morfismo de presentación finita y sea $Z$ una parte localmente construible de $X$.
Entonces el conjunto $E$ de los $s\in S$ tales que $Z_s\neq\varnothing$ es localmente construible.
\end{proposition}

\begin{proof}
En efecto, tenemos $E=f(Z)$, y basta aplicar el teorema de Chevalley \sref{IV.1.8.4}.
\end{proof}

\begin{corollary}[9.5.2]
\label{IV.9.5.2}
Si $Z'$, $Z''$ son dos partes localmente construibles de $X$, el conjunto de los $s\in S$ tales que $Z'_s\subset Z''_s$ (resp.\ $Z'_s=Z''_s$) es localmente construible.
\end{corollary}

\begin{proof}
En efecto, la relación $Z'_s\subset Z''_s$ equivale a $(Z'\cap(\complement Z''))_s=\varnothing$, y $Z'\cap\complement Z''$ es localmente construible.
\end{proof}

\begin{proposition}[9.5.3]
\label{IV.9.5.3}
Sea $f:X\to S$ un morfismo de presentación finita, y sean $Z$, $Z'$ dos partes localmente construibles de $X$ tales que $Z\subset Z'$.
Entonces el conjunto $E$ de los $s\in S$ tales que $Z_s$ es denso en $Z'_s$ es localmente construible en $S$.
\end{proposition}

\begin{proof}
Debemos verificar las dos condiciones de \sref{IV.9.2.2},~(iii).
En cuanto a la primera, consideremos un preesquema $X$ algebraico sobre un cuerpo $k$, dos partes construibles $Z$, $Z'$ de $X$ tales que $Z\subset Z'$, y una extensión $k'$ de $k$.
Entonces la proyección canónica $p:X_{(k')}\to X$ es un morfismo fielmente plano y cuasicompacto, y tenemos $p^{-1}(\overline{Z})=\overline{p^{-1}(Z)}$ y $p^{-1}(\overline{Z'})=\overline{p^{-1}(Z')}$; en virtud de \sref{IV.2.3.10}, como $p$ es sobreyectivo, la relación $\overline{Z}=\overline{Z'}$ equivale a $\overline{p^{-1}(Z)}=\overline{p^{-1}(Z')}$.

Verifiquemos la segunda condición y supongamos, por tanto, que $S$ es afín, noetheriano, íntegro, de punto genérico $\eta$.
Distinguimos dos casos:
\begin{enumerate}[label=\upshape{$\arabic{enumi}^\circ$}]
  \item $\eta\in S-E$, es decir, $Z_\eta$ no es denso en $Z'_\eta$; existe entonces en $X$ un abierto $V$ tal que $V\cap Z_\eta=\varnothing$ y $V\cap Z'_\eta\neq\varnothing$.
  Como $X$ es noetheriano, $V$ es localmente construible, y por tanto también lo es $V\cap Z$; en virtud de \sref{IV.9.5.1}, existe un entorno $U$ de $\eta$ en $S$ tal que, para todo $s\in U$, $(V\cap Z)_s=\varnothing$ y $(V\cap Z')_s\neq\varnothing$; esto implica que $Z_s$ no es denso en $Z'_s$ para $s\in U$, es decir, $U\subset S-E$.

\oldpage[IV-3]{68}

  \item\label{IV.9.5.3.case2} $\eta\in E$, por lo que $Z_\eta$ es denso en $Z'_\eta$.
  Demostremos primero que podemos suponer $Z'$ \emph{cerrado}.
  En efecto, $Z_\eta$ es denso en $\overline{Z'_\eta}$ (adherencia tomada en $X_\eta$); poniendo $V_\eta=X_\eta-\overline{Z_\eta}$, que es abierto en $X_\eta$ y no encuentra $Z'_\eta$, podemos suponer $V_\eta$ de la forma $V\cap X_\eta$, donde $V$ es abierto (y por tanto construible) en $X$, y la hipótesis $V_\eta\cap Z'_\eta=\varnothing$ implica entonces $V_s\cap Z'_s=\varnothing$ para todo $s$ en un entorno de $\eta$ en $S$, en virtud de \sref{IV.9.5.1}.
  Sustituyendo $S$ por un entorno abierto de $\eta$, podemos, por tanto, suponer $V\cap Z'=\varnothing$, de donde $V\cap\overline{Z'}=\varnothing$ (adherencia tomada en $X$), y por consiguiente $(\overline{Z'})_\eta=\overline{Z'_\eta}$, lo que demuestra nuestra afirmación.
  Como la descomposición de $Z'$ en sus componentes irreducibles es una unión finita, y restringiendo aún más $S$ a un entorno de $\eta$, podemos suponer que todas las componentes irreducibles $Z'_i$ de $Z'$ encuentran $X_\eta$, de donde se sigue que $X_\eta$ contiene los puntos genéricos de los $Z'_i$ \sref[0]{0.2.1.8}.
  Decir que $Z_s$ es denso en $Z'_s$ equivale a decir que cada uno de los $(Z\cap Z'_i)_s$ es denso en $(Z'_i)_s$, y nos reducimos así al caso en que $Z'$ es \emph{irreducible}.
  Sustituyendo entonces $X$ por el subpreesquema cuyo espacio subyacente reducido es $Z'$, vemos que podemos suponer $Z'=X$ y que $X$ es \emph{íntegro} y domina $S$.
  Finalmente, recubriendo $X$ por un número finito de abiertos afines $W_j$ y sustituyendo $Z$ por $Z\cap W_j$, podemos suponer que $X=\operatorname{Spec}(A)$, donde $A$ es un anillo noetheriano íntegro.
  Como $X_\eta$ es noetheriano íntegro y $Z_\eta$ es construible y denso en $X_\eta$, $Z_\eta$ contiene un abierto no vacío de $X_\eta$~($\mathbf{0}_{\mathrm{III}}$,~9.2.2), que puede tomarse de la forma $(D(t))_\eta$, donde $t$ es un elemento no nulo de $A$.
  Sustituyendo $S$ por un entorno de $\eta$ en $S$ si es necesario, en virtud de la relación $(D(t))_\eta\subset Z_\eta$, podemos suponer además que $D(t)\subset Z$ \sref{IV.9.5.2}.
  Finalmente, como la homotecia de razón $t_\eta$ en $\sh{O}_{X_\eta}$ es inyectiva, se sigue de \sref{IV.9.4.5} que para $s$ próximo a $\eta$, $t_s$ es $\sh{O}_{X_s}$-regular, de modo que $(X_s)_{t_s}$ es denso en $X_s$, y \emph{a fortiori} lo mismo vale para $Z_s$, que contiene $(X_s)_{t_s}$.
\end{enumerate}
\end{proof}

\begin{corollary}[9.5.4]
\label{IV.9.5.4}
Sea $f:X\to S$ un morfismo de presentación finita y sea $Z$ una parte localmente construible de $X$.
El conjunto $E$ de los $s\in S$ tales que $Z_s$ es cerrado (resp.\ abierto, resp.\ localmente cerrado) en $X_s$ es localmente construible en $S$.
\end{corollary}

\begin{proof}
Decir que $Z_s$ es abierto en $X_s$ significa que $(X-Z)_s=X_s-Z_s$ es cerrado en $X_s$, y como $X-Z$ es localmente construible, podemos limitarnos a considerar el conjunto de los $s\in S$ tales que $Z_s$ es cerrado y el conjunto de los $s\in S$ tales que $Z_s$ es localmente cerrado.

Verifiquemos en cada caso las dos condiciones de \sref{IV.9.2.2},~(ii).
La primera resulta del hecho de que $p:X_{(k')}\to X$ es fielmente plano y cuasicompacto, y de \sref{IV.2.3.12} y \sref{IV.2.3.14}.
Verifiquemos, por tanto, la segunda condición, suponiendo $S$ afín, noetheriano, íntegro, de punto genérico $\eta$.
Pongamos $Z'=\overline{Z}$; el mismo razonamiento que en \sref{IV.9.5.3} muestra que $Z'_\eta$ es igual a la adherencia de $Z_\eta$ en $X_\eta$; en virtud de \sref{IV.9.5.3}, existe, pues, un entorno $U$ de $\eta$ tal que para $s\in U$, $Z_s$ es denso en $Z'_s$, siendo este último cerrado.
Decir que $Z_s$ es cerrado en $X_s$ significa que $Z''_s=\varnothing$, donde $Z''=Z'-Z$; se sigue de \sref{IV.9.5.1} que el conjunto $E$ de los $s\in S$ donde $Z''_s=\varnothing$ es tal que $E$ o $S-E$ contiene un entorno de $\eta$.
Decir que $Z_s$ es localmente cerrado en $X_s$ significa que $Z''_s$ es cerrado en $X_s$; basta, por tanto, aplicar el resultado precedente sustituyendo $Z$ por $Z''$, que es localmente construible en $X$.
\end{proof}

\oldpage[IV-3]{69}

\begin{proposition}[9.5.5]
\label{IV.9.5.5}
Sea $f:X\to S$ un morfismo de presentación finita y sea $Z$ una parte localmente construible de $X$ tal que, para todo $s\in S$, $Z_s$ es localmente cerrado en $X_s$.
Para todo $s\in S$, sea $D(s)\subset\mathbf{Z}\cup\{-\infty\}$ el conjunto de las dimensiones de las componentes irreducibles de $Z_s$.
Entonces la función $s\mapsto D(s)$ es localmente construible en $S$.
\end{proposition}

\begin{proof}
Sea $\Phi$ una parte finita de $\mathbf{Z}\cup\{-\infty\}$; se trata de demostrar que el conjunto de los $s\in S$ tales que $D(s)=\Phi$ es localmente construible; teniendo en cuenta \sref{IV.9.2.3}, debemos verificar una vez más las dos condiciones de \sref{IV.9.2.2},~(iii).

En cuanto a la primera, se trata de ver que si $X$ es un preesquema algebraico sobre un cuerpo $k$, $Z$ una parte localmente cerrada de $X$, $k'$ una extensión de $k$, y $p:X_{(k')}\to X$ la proyección canónica, entonces el conjunto de las dimensiones de las componentes irreducibles de $Z$ es el mismo que el de las dimensiones de las componentes irreducibles de $p^{-1}(Z)$; teniendo en cuenta la existencia de un subpreesquema de $X$ cuyo espacio subyacente es $Z$ \sref[I]{I.5.2.1}, esto se sigue de \sref{IV.4.2.8}.

Para la segunda condición de \sref{IV.9.2.2},~(iii), estamos en el caso en que $S$ es noetheriano, íntegro, de punto genérico $\eta$, y $f:X\to S$ es un morfismo de tipo finito, de modo que $X$ es noetheriano.
El subespacio $Z_\eta$ del espacio noetheriano $X_\eta$ es, por hipótesis, localmente cerrado, por lo que tiene un número finito de componentes irreducibles $Z_{i\eta}$, que son localmente cerradas en $X_\eta$.
Se sigue, para cada índice $i$, que existe una parte localmente cerrada $Z_i$ de $X$ tal que $(Z_i)_\eta=Z_{i\eta}$, por lo que si $Z'=\bigcup_i Z_i$, tenemos $Z'_\eta=Z_\eta$.
Pero como $Z$ y $Z'$ son localmente construibles, podemos, sustituyendo $S$ por un entorno de $\eta$, suponer que $Z'=Z$ \sref{IV.9.5.1}.
Además, para $i\neq j$, $Z_{i\eta}\cap Z_{j\eta}$ es raro y \emph{cerrado} en $Z_{i\eta}$; por tanto, mediante \sref{IV.9.5.3} y \sref{IV.9.5.4}, podemos también suponer, restringiendo $S$, que para $j\neq i$, $(Z_i)_s\cap(Z_j)_s$ es raro y cerrado en $(Z_i)_s$.
Como $Z_i$ es localmente cerrado en $X$, existe un subpreesquema de $X$ cuyo espacio subyacente es $Z_i$ (y que seguimos denotando por $Z_i$), que es de tipo finito sobre $S$ \sref[I]{I.6.3.5}.
Pongamos $U_i=Z_i-\bigcup_{j\neq i}(Z_i\cap Z_j)$, que es abierto en $Z_i$ y tal que $(U_i)_s$ es abierto y denso en todas partes en $(Z_i)_s$; además, por construcción, los $U_i$ no tienen dos a dos ningún punto común.
Como $(Z_i)_s$ es un $\boldsymbol{k}(s)$-preesquema algebraico, el conjunto de las dimensiones de las componentes irreducibles de $Z_s=\bigcup_i(Z_i)_s$ es el mismo que el conjunto de las dimensiones de las componentes irreducibles de la reunión de los $(U_i)_s$ \sref{IV.4.1.1.3}, siendo cada una de estas componentes ya una componente irreducible de uno de los $(U_i)_s$.
Podemos, por tanto, limitarnos al caso en que $Z=U_i$; además, como $Z_\eta$ es entonces irreducible, solo hay una componente irreducible de $Z$ que encuentra $X_\eta$~\sref[0]{0.2.1.8}, y podemos evidentemente, restringiendo $Z$, suponer $Z$ irreducible.
Nos reducimos así finalmente a demostrar el siguiente caso particular de \sref{IV.9.5.5}:
\end{proof}

\begin{corollary}[9.5.6]
\label{IV.9.5.6}
Sea $S$ un preesquema noetheriano, irreducible, de punto genérico $\eta$, sea $X$ un preesquema irreducible y sea $f:X\to S$ un morfismo dominante de tipo finito.
Entonces existe un entorno $U$ de $\eta$ en $S$ tal que, para todo $s\in U$, todas las componentes irreducibles de $X_s$ son de dimensión $n=\dim(X_\eta)$.
\end{corollary}

\begin{proof}
Podemos evidentemente reducirnos al caso en que $S=\operatorname{Spec}(A)$ es afín, siendo por tanto $A$ noetheriano; sustituyendo $f$ por $f_{\mathrm{red}}$, que es de tipo finito \sref[I]{I.5.4},~(vii),

\oldpage[IV-3]{70}

podemos suponer $A$ íntegro y $X$ reducido, por tanto íntegro.
Sabemos por \sref{IV.4.1.2} que existe un abierto $W$ denso en $X_\eta$ y un $\boldsymbol{k}(\eta)$-morfismo \emph{finito} y \emph{sobreyectivo} $h:W\to\mathbf{V}((\boldsymbol{k}(\eta))^n)=\operatorname{Spec}(B')$, donde $B'=\boldsymbol{k}(\eta)[\mathrm{T}_1,\ldots,\mathrm{T}_n]$ (siendo los $\mathrm{T}_i$ indeterminados).
Si $V$ es un abierto de $X$ tal que $V\cap X_\eta=W$, sabemos por \sref{IV.9.5.3} que para $s$ próximo a $\eta$ en $S$, $V_s$ es un abierto denso en $X_s$, y podemos por tanto, por \sref{IV.4.1.1.3}, reducirnos al caso en que $V=X$, $W=X_\eta$.
Pongamos $B=A[\mathrm{T}_1,\ldots,\mathrm{T}_n]$ y $Y=\operatorname{Spec}(B)=\mathbf{V}_A^n$, de modo que $\operatorname{Spec}(B')=Y_\eta$; se sigue de \sref{IV.8.8.2},~(i) y \sref{IV.8.10.5},~(vi) y (x), aplicados según el método de \sref{IV.8.1.2},~\emph{a}), que, sustituyendo $S$ por un entorno de $\eta$, podemos suponer que $h=g_\eta$, donde $g:X\to Y$ es un morfismo finito sobreyectivo; en otras palabras, tenemos $X=\operatorname{Spec}(C)$, donde $C$ es una B-álgebra finita y el homomorfismo $B\to C$ correspondiente a $g$ es \emph{inyectivo}; como $C$ es un anillo íntegro, $C$ es por tanto un B-módulo de tipo finito \emph{sin torsión}, y $C_\eta=C\otimes_A\boldsymbol{k}(\eta)$ es un módulo de tipo finito sin torsión sobre $B_\eta=B'=B\otimes_A\boldsymbol{k}(\eta)=\boldsymbol{k}(\eta)[\mathrm{T}_1,\ldots,\mathrm{T}_n]$ (pues es un anillo de fracciones cuyos denominadores pertenecen a $A-\{0\}\subset B-\{0\}$).
Se sigue de \sref{IV.9.4.8} que existe un entorno $U$ de $\eta$ en $S$ tal que para $s\in S$, $C_s=C\otimes_A\boldsymbol{k}(s)$ es un módulo de tipo finito \emph{sin torsión} sobre $B_s=B\otimes_A\boldsymbol{k}(s)=\boldsymbol{k}(s)[\mathrm{T}_1,\ldots,\mathrm{T}_n]$, y en particular el homomorfismo $g_s:B_s\to C_s$ es inyectivo.
Como ningún elemento de $B_s$ es divisor de cero en $C_s$, para todo ideal primo minimal $\mathfrak{p}_i$ de $C_s$ (cuyos elementos son divisores de cero en $C_s$), tenemos necesariamente $\mathfrak{p}_i\cap B_s=0$, de donde el homomorfismo canónico $B_s\to C_s/\mathfrak{p}_i$ es inyectivo.
Deducimos que para cada componente irreducible $Z_i=\operatorname{Spec}(C_s/\mathfrak{p}_i)$ de $X_s$, la restricción a $Z_i$ de $g$ es un morfismo finito y dominante $Z_i\to Y_s$, por tanto \emph{sobreyectivo} \sref[II]{II.6.1.10}.
Concluimos entonces de \sref{IV.4.1.2} que $\dim Z_i=n$, lo que termina la demostración.
\end{proof}

\begin{remark}[9.5.7]
\label{IV.9.5.7}
Conviene observar que bajo las hipótesis de \sref{IV.9.5.6}, puede ocurrir que $X_s$ no sea irreducible para \emph{ningún} $s\neq\eta$ en un entorno de $\eta$; en otras palabras, la propiedad ``$X$ es un $k$-preesquema algebraico irreducible'' \emph{no es construible}.
Tomemos, por ejemplo, $S=\operatorname{Spec}(k[\mathrm{T}])$, donde $k$ es un cuerpo algebraicamente cerrado y $\mathrm{T}$ un indeterminado; tenemos entonces $\boldsymbol{k}(\eta)=K=k(\mathrm{T})$.
Sea $L$ una extensión separable finita de $K$ de grado $>1$, y sea $X$ la \emph{clausura integral} de $S$ en $L$ \sref[II]{II.6.3.4}; tenemos entonces $X=\operatorname{Spec}(B)$, donde $B$ es la clausura integral de $k[\mathrm{T}]$ en $L$.
Sabemos que $B$ es un anillo de Dedekind, y que todos salvo un número finito de ideales maximales de $k[\mathrm{T}]$ no están ramificados en $B$. Además, el cuerpo residual de todo ideal maximal de $B$ es necesariamente $k$, puesto que es una extensión finita del cuerpo algebraicamente cerrado $k$. Así, para casi todos los ideales maximales $\mathfrak{j}_s$ de $k[\mathrm{T}]$, el anillo $B_s=B/\mathfrak{j}_s B$ es un producto directo de $[L:K]$ copias de $k$; en particular, $X_s$ no es irreducible, aunque $X_\eta=\operatorname{Spec}(L)$ sí lo es.
El mismo ejemplo muestra que la propiedad ``$X$ es un $k$-preesquema algebraico íntegro'' no es construible.
Finalmente, la propiedad ``$X$ es un $k$-preesquema algebraico reducido'' tampoco es construible.
Para ello basta tomar de nuevo $S=\operatorname{Spec}(k[\mathrm{T}])$, donde esta vez $k$ es un cuerpo algebraicamente cerrado de característica $p>0$, y como $X$ la clausura integral de $S$ en $L=K^{1/p}$ (donde $K=k(\mathrm{T})$), de modo que $X=\operatorname{Spec}(B)$ con $B=k[\mathrm{T}^{1/p}]$ (siendo $k$ perfecto); todo ideal maximal de $k[\mathrm{T}]$ es de la forma $(\mathrm{T}-\alpha)$ con $\alpha\in k$; el único ideal de $B$ sobre el ideal $(\mathrm{T}-\alpha)$ es el ideal principal $(\mathrm{T}^{1/p}-\alpha^{1/p})$ y es inmediato
\oldpage[IV-3]{71}
que el anillo cociente $B/(\mathrm{T}-\alpha)B$ tiene elementos nilpotentes; en otras palabras, $X_s$ no es reducido para ningún $s\neq\eta$, mientras que $X_\eta=\operatorname{Spec}(L)$ es íntegro.

Veremos un poco más adelante \sref{IV.9.7} que, por el contrario, se obtienen propiedades construibles cuando se consideran las nociones ``geométricas'' correspondientes a las nociones de preesquema irreducible, reducido o íntegro \sref{IV.4.5} y \sref{IV.4.6}.
\end{remark}

\subsection{Construibilidad de ciertas propiedades de morfismos}
\label{subsection:IV.9.6}

\begin{proposition}[9.6.1]
\label{IV.9.6.1}
Sean $X$, $Y$ dos S-preesquemas de presentación finita y sea $f:X\to Y$ un S-morfismo.
Sea $E$ el conjunto de los $s\in S$ para los cuales $f_s$ posee una de las propiedades siguientes: ser
\begin{enumerate}[label=(\roman*),widest=xii]
  \item sobreyectivo;
  \item dominante;
  \item separado;
  \item propio;
  \item radicial;
  \item finito;
  \item cuasifinito;
  \item una inmersión;
  \item una inmersión cerrada;
  \item una inmersión abierta;
  \item un isomorfismo;
  \item un monomorfismo.
\end{enumerate}
Entonces $E$ es localmente construible en $S$.
\end{proposition}

\begin{proof}
Las afirmaciones relativas a (i), (v) y (vii) se incluyen solo como referencia, pues ya fueron establecidas en \sref{IV.9.3.2}.

(ii): Observemos primero que $f$ es de presentación finita \sref[I]{I.6.2},~(v); por tanto, en virtud del teorema de Chevalley \sref{IV.1.8.4}, $f(X)$ es localmente construible en $Y$.
Además, tenemos $f_s(X_s)=(f(X))_s$ \sref[I]{I.3.6.1}; el conjunto de los $s$ tales que $f_s$ es dominante es el conjunto de los $s$ tales que $(f(X))_s$ es denso en $Y_s$; la conclusión se sigue, pues, de \sref{IV.9.5.3}.

(iii): Como $f:X\to Y$ es de presentación finita, la inmersión diagonal $\Delta_f:X\to X\times_Y X$ es de presentación finita \sref[I]{I.6.2},~(iv) y (v); se sigue de \sref{IV.1.8.4.1} que $\Delta_f(X)=\Delta_X$ es una parte localmente construible de $X\times_Y X$.
Decir que $f_s$ es separado significa (teniendo en cuenta \sref[I]{I.5.3.4}) que $(\Delta_X)_s$ es cerrado en $X_s\times_{Y_s} X_s=(X\times_Y X)_s$; la conclusión se sigue de \sref{IV.9.5.4}.

\oldpage[IV-3]{72}

(iv): Demostremos que la propiedad ``$X$ e $Y$ son preesquemas algebraicos sobre un cuerpo $k$, y $f:X\to Y$ un $k$-morfismo propio'' es construible.
La condición \sref{IV.9.2.1},~(i) se verifica en virtud de \sref{IV.2.7.1},~(vii).
Podemos, por tanto, reducirnos al caso en que $S$ es afín, noetheriano, íntegro, de punto genérico $\eta$, y debemos demostrar que $E$ o $S-E$ es un entorno de $\eta$.
Supongamos primero que $\eta\in E$; se sigue de \sref{IV.8.10.5},~(xii), aplicado según el método de \sref{IV.8.1.2},~\emph{a}), que, sustituyendo $S$ por un entorno de $\eta$, podemos
suponer que el propio $f$ es un morfismo propio; sabemos entonces que lo mismo vale para $f_s$ para todo $s\in S$ \sref[II]{II.5.4.2},~(iii).
Supongamos, por el contrario, que $\eta\in S-E$, y distingamos dos casos:
\begin{enumerate}[label=\upshape{$\arabic{enumi}^\circ$}]
  \item Supongamos que $f_\eta$ no es separado; se sigue entonces de (iii) que $f_s$ no es separado (y \emph{a fortiori} no es propio) en un entorno de $\eta$.

  \item Supongamos $f_\eta$ separado; $Y$ es una unión finita de abiertos afines $V_i$, y para que $f_s$ sea propio, basta que cada una de sus restricciones $f_s^{-1}((V_i)_s)\to(V_i)_s$ lo sea \sref[II]{II.5.4.1}; podemos, por tanto, reducirnos al caso en que $Y$ es \emph{afín}, y por tanto un esquema.
  Decir que $f_\eta$ no es propio significa, por \sref[II]{II.5.6.3}, que existe un morfismo \emph{de tipo finito} $h:Z\to Y_\eta$ tal que el morfismo $(f_\eta)_{(Z)}:X_\eta\times_{Y_\eta} Z\to Z$ no es cerrado.
  Como $Y$ es un esquema, deducimos de \sref{IV.8.8.2},~(i) y (ii) (restringiendo $S$, si es necesario, a un entorno de $\eta$) que existe un morfismo de tipo finito $g:Y'\to Y$ tal que $Z$ es isomorfo a $Y'_\eta$ y que $g_\eta=h$;
  si ponemos $X'=X\times_Y Y'$ y $f'=f_{(Y')}:X'\to Y'$, entonces $(f_\eta)_{(Z)}=f'_\eta$, y por hipótesis existe, por tanto, una parte cerrada $M'$ de $X'_\eta$ tal que $f'_\eta(M')$ no es cerrado en $Y'_\eta$.
  Ahora bien, $M'$ es la traza sobre $X'_\eta$ de una parte cerrada $N'$ de $X'$; como $X'$ es noetheriano y $f'$ de tipo finito, $f'(N')$ es construible en $Y'$ \sref{IV.1.8.4} y, por hipótesis, $(f'(N'))_\eta=f'_\eta(N'_\eta)=f'_\eta(M')$ no es cerrado en $Y'_\eta$.
  Concluimos de \sref{IV.9.5.4} que existe un entorno $U$ de $\eta$ en $S$ tal que para $s\in U$, $(f'(N'))_s=f'_s(N'_s)$ no es cerrado en $Y'_s$; en otras palabras, el morfismo $f'_s$ no es cerrado, y \emph{a fortiori} el morfismo $f_s$ no es propio.
\end{enumerate}

(vi): La propiedad es una conjunción de las propiedades (iv) y (vii) \sref{IV.8.11.1}; por tanto, la proposición se sigue de lo ya demostrado.

(ix): La verificación de la condición \sref{IV.9.2.1},~(i) se sigue de \sref{IV.2.7.1},~(xi).
Podemos, por tanto, reducirnos al caso en que $S$ es afín, noetheriano, íntegro, de punto genérico $\eta$, y basta demostrar que $E$ o $S-E$ es un entorno de $\eta$.
Si $\eta\in E$, podemos (sustituyendo $S$ por un entorno de $\eta$) suponer que $f$ es una inmersión cerrada en virtud de \sref{IV.8.10.5},~(iv), y entonces está claro que $f_s$ es una inmersión cerrada para todo $s\in S$ \sref[I]{I.4.3.2}.
Supongamos, pues, que $\eta\in S-E$, y distingamos dos casos:
\begin{enumerate}[label=\upshape{$\arabic{enumi}^\circ$}]
  \item $f_\eta$ no es un morfismo finito.
  Se sigue entonces de (vi) que en un entorno de $\eta$, $f_s$ no es finito, ni \emph{a fortiori} una inmersión cerrada.

  \item $f_\eta$ es finito; entonces, en virtud de \sref{IV.8.10.5},~(x), podemos suponer (restringiendo $S$ si es necesario) que el propio $f$ es un morfismo \emph{finito}.
  En este caso $\mathscr{A}=\mathscr{A}(X)=f_*(\sh{O}_X)$ es un $\sh{O}_Y$-Módulo coherente \sref[II]{II.6.1.3} y $f=\mathscr{A}(u)$, donde $u:\sh{O}_Y\to\mathscr{A}$ es un homomorfismo de $\sh{O}_Y$-Álgebras \sref[II]{II.1.1.2}; como $f_\eta$ no es una inmersión cerrada por hipótesis, $u_\eta$ no es sobreyectivo \sref[II]{II.1.4.10}, por tanto, por \sref{IV.9.4.5}, existe un entorno de $\eta$ en el cual $u_s$ no es sobreyectivo y, en consecuencia, por \sref[I]{I.4.2.3}, $f_s$ no es una inmersión cerrada.
\end{enumerate}

(viii): La verificación de la condición \sref{IV.9.2.1},~(i) se hace como en (ix), utilizando esta vez \sref{IV.2.7.1},~(x) y el hecho de que toda inmersión de un preesquema noetheriano en otro es cuasicompacta.
Nos reducimos, por tanto, al caso en que $S$ es afín, noetheriano, íntegro, de punto genérico $\eta$, y basta demostrar que $E$ o $S-E$ es un entorno de $\eta$.
Si $\eta\in E$, concluimos como en (ix) con ayuda de \sref{IV.8.10.5},~(ii).
Si $\eta\in S-E$, distinguimos de nuevo dos casos:
\begin{enumerate}[label=\upshape{$\arabic{enumi}^\circ$}]
  \item $f_\eta(X_\eta)$ no es una parte localmente cerrada de $Y_\eta$.
  Como $f(X)$ es construible en $Y$ \sref{IV.1.8.4} y $f_s(X_s)=(f(X))_s$, deducimos de \sref{IV.9.5.4} que para $s$ próximo a $\eta$, $f_s(X_s)$ no es localmente cerrado en $Y_s$, y \emph{a fortiori} $f_s$ no es una inmersión.

  \item\oldpage[IV-3]{73} $f_\eta(X_\eta)$ es localmente cerrado en $Y_\eta$.
  Como $f(X)$ es construible en $Y$ \sref{IV.1.8.4} y lo mismo vale para $\overline{f(X)}$, puesto que $Y$ es noetheriano, se sigue de \sref{IV.8.3.11} que, restringiendo $S$ si es necesario, podemos suponer que $f(X)$ es localmente cerrado en $Y$.
  Existe entonces un abierto $V$ de $Y$ que contiene $f(X)$ y en el cual $f(X)$ es cerrado.
  Como $Y$ es noetheriano, $V$ es de tipo finito sobre $S$; sustituyendo $Y$ por $V$, podemos, por tanto, reducirnos al caso en que $f(X)$ es cerrado en $Y$.
  Pero entonces $f_s(X_s)=(f(X))_s$ es cerrado en $Y_s$ para todo $s\in S$, y decir que $f_s$ es una inmersión equivale a decir que es una inmersión cerrada.
  Nos reducimos, por tanto, al caso~(ix).
\end{enumerate}

(x): Utilizando esta vez \sref{IV.2.7.1},~(ix) y \sref{IV.8.10.5},~(iii), nos reducimos al caso en que $S$ es afín, noetheriano, íntegro, de punto genérico $\eta$, y en que $\eta\in S-E$.
Distinguimos tres casos:
\begin{enumerate}[label=\upshape{$\arabic{enumi}^\circ$}]
  \item $f_\eta(X_\eta)$ no es abierto en $Y_\eta$.
  Como $f(X)$ es construible en $Y$ \sref{IV.1.8.4}, deducimos de \sref{IV.9.5.4} que $f_s(X_s)$ no es abierto en $Y_s$ para $s$ próximo a $\eta$, y \emph{a fortiori} $f_s$ no es una inmersión abierta.

  \item $f_\eta(X_\eta)$ es abierto en $Y_\eta$ pero $f_\eta$ no es una inmersión.
  Se sigue entonces de (viii) que para $s$ próximo a $\eta$, $f_s$ no es una inmersión, ni \emph{a fortiori} una inmersión abierta.

  \item $f_\eta(X_\eta)$ es abierto en $Y_\eta$ y $f_\eta$ es una inmersión.
  Como $f(X)$ es construible en $Y$ \sref{IV.1.8.4}, se sigue de \sref{IV.8.3.11} que, restringiendo $S$ si es necesario, podemos suponer ya que $f(X)$ es abierto en $Y$.
  Como $Y$ es noetheriano, el subpreesquema inducido sobre $f(X)$ es de tipo finito sobre $S$, por lo que podemos reducirnos al caso en que $f$ es \emph{sobreyectivo} sustituyendo $Y$ por $f(X)$.
  Por hipótesis, $f_\eta$ es una inmersión cerrada, por lo que, como en (ix), podemos suponer que $f$ es una \emph{inmersión cerrada}; en otras palabras, que $X$ es un subpreesquema cerrado de $Y$ definido por un ideal coherente $\mathscr{J}$ de $\sh{O}_Y$.
  Por hipótesis, $f_\eta$ no es un isomorfismo, de donde $\mathscr{J}_\eta\neq 0$; concluimos por \sref{IV.9.4.5} que en un entorno de $\eta$, $\mathscr{J}_s\neq 0$, y por consiguiente $f_s$ no es una inmersión abierta.
\end{enumerate}

(xi): La propiedad es una conjunción de las propiedades (i) y (x), y por tanto se sigue de lo demostrado.

(xii): En virtud de \sref[I]{I.5.3.8}, decir que $f_s$ es un monomorfismo significa que $\Delta_{f_s}=(\Delta_f)_s:X_s\to X_s\times_{Y_s} X_s=(X\times_Y X)_s$ es un isomorfismo, y como $X\times_Y X$ es un S-preesquema de presentación finita \sref[I]{I.6.2},~(iv), la conclusión se sigue de (xi).
\end{proof}

\begin{proposition}[9.6.2]
\label{IV.9.6.2}
Sean $X$, $Y$ dos S-preesquemas de presentación finita y sea $f:X\to Y$ un S-morfismo.
\begin{enumerate}[label=\upshape{I)}]
  \item Sea $E$ el conjunto de los $s\in S$ para los cuales $f_s$ posee una de las propiedades siguientes: ser
  \begin{enumerate}[label=(\roman*)]
    \item afín;
    \item cuasiafín;
    \item proyectivo;
    \item cuasiproyectivo.
  \end{enumerate}
  Entonces $E$ es ind-construible en $S$.
\end{enumerate}

\oldpage[IV-3]{74}

\begin{enumerate}[label=\upshape{II)}]
  \item Sea $\sh{L}$ un $\sh{O}_X$-Módulo invertible.
  Entonces el conjunto $E'$ de los $s\in S$ tales que $\sh{L}_s$ es un $\sh{O}_{X_s}$-Módulo amplio (resp.\ muy amplio) relativamente a $f_s$, es ind-construible en $S$.
\end{enumerate}
\end{proposition}

\begin{proof}
I) Verifiquemos las condiciones (i) y (ii) de \sref{IV.9.2.1}.
En cuanto a la condición \sref{IV.9.2.1},~(i), se sigue, para las propiedades (i) y (ii), de \sref{IV.2.7.1},~(xiii) y (xiv); para las propiedades (iii) y (iv), se sigue de \sref{IV.9.1.5}.
Verifiquemos entonces la condición \sref{IV.9.2.1},~(ii), suponiendo, por tanto, $S$ noetheriano, íntegro, de punto genérico $\eta$, y que $f_\eta:X_\eta\to Y_\eta$ posee una de las propiedades (i) a (iv) del enunciado.
Aplicando \sref{IV.8.10.5},~(viii), (ix), (xiii) y (xiv) según el método de \sref{IV.8.1.2},~\emph{a}), vemos primero que existe un entorno abierto $U$ de $\eta$ tal que, si $V$ y $W$ son las imágenes inversas respectivas de $U$ en $X$ e $Y$ por los morfismos estructurales, la restricción $V\to W$ de $f$ posee la propiedad (i), (ii), (iii) o (iv) considerada.
La conclusión se sigue entonces del hecho de que todas las propiedades consideradas son estables por cambio de base.

II) Procedemos del mismo modo.
Para la condición \sref{IV.9.2.1},~(i), se sigue esta vez de \sref{IV.2.7.2}.
Para la condición \sref{IV.9.2.1},~(ii), con las mismas notaciones que en I), se sigue de \sref{IV.8.10.5.2} que para un entorno $U$ de $\eta$ en $S$, la restricción $\sh{L}|V$ es amplia (resp.\ muy amplia) relativamente a la restricción $V\to W$ de $f$.
La conclusión se sigue de nuevo de la estabilidad de las propiedades consideradas por cambio de base \sref[II]{II.4.6.13} y \sref[II]{II.4.4.10}.
\end{proof}

Podemos mejorar \sref{IV.9.6.2},~II) bajo ciertas condiciones:

\begin{proposition}[9.6.3]
\label{IV.9.6.3}
Sean $X$, $Y$ dos S-preesquemas de presentación finita, sea $f:X\to Y$ un S-morfismo propio y sea $\sh{L}$ un $\sh{O}_X$-Módulo invertible.
Entonces el conjunto $E$ (resp.\ $E'$) de los $s\in S$ tales que $\sh{L}_s$ es un $\sh{O}_{X_s}$-Módulo amplio (resp.\ muy amplio) relativamente a $f_s:X_s\to Y_s$ es localmente construible en $S$.
\end{proposition}

\begin{proof}
Sea $k$ un cuerpo, sean $Z$, $T$ dos preesquemas algebraicos sobre $k$, sea $\mathscr{M}$ un $\sh{O}_Z$-Módulo invertible y sea $g:Z\to T$ un $k$-morfismo de tipo finito; entonces, si $\mathbf{P}(Z,T,\mathscr{M},g,k)$ denota la relación ``$\mathscr{M}$ es amplio (resp.\ muy amplio) relativamente a $g$'', ya hemos observado en \sref{IV.9.6.2} que $\mathbf{P}(Z,T,\mathscr{M},g,k)$ satisface la condición \sref{IV.9.2.1},~(i) en virtud de \sref{IV.2.7.2}.
Sabemos, por otra parte, que $E$ y $E'$ son ind-construibles.
Queda, por tanto, ver que si $S$ es noetheriano, íntegro, de punto genérico $\eta$ y si $\eta\in S-E$ (resp.\ $\eta\in S-E'$), entonces $S-E$ (resp.\ $S-E'$) contiene un entorno de $\eta$.
Consideraremos separadamente los casos $E$ y $E'$.

\oldpage[IV-3]{75}

\emph{Caso de $E'$.} ---
Como $f$ es separado y $Y$ es cuasicompacto, existe un entero $h$ tal que
$R^qf_*(\sh{L})=0$ para $q>h$ \sref[III]{III.1.4.12}.
Por otra parte, como $f$ es propio y $Y$ es noetheriano, todos los
$R^qf_*(\sh{L})$ son $\sh{O}_Y$-Módulos coherentes \sref[III]{III.3.2.1}.
Se anulan salvo para un número finito de valores de $q$, por lo que el teorema de planitud genérica
\sref{IV.6.9.1} muestra que, tras restringir $S$ a un entorno de $\eta$, podemos
suponer que $\sh{L}$ y todos los $R^qf_*(\sh{L})$ son planos sobre $S$.
Se sigue entonces de \sref[III]{III.6.9.9} que el homomorfismo canónico
\begin{equation}
\label{IV.9.6.3.1}
 f_*(\sh{L})\otimes_{\sh{O}_S}\boldsymbol{k}(s)
 \longrightarrow (f_s)_*(\sh{L}_s)
 \tag{9.6.3.1}
\end{equation}
es un isomorfismo.

Se sigue de \sref[II]{II.4.4.4} que decir que $\sh{L}_s$ no es muy amplio relativamente a $f_s$ significa que: \emph{o bien} el homomorfismo canónico $(f_s)^*((f_s)_*(\sh{L}_s))\to\sh{L}_s$ no es sobreyectivo; \emph{o bien} el homomorfismo precedente es sobreyectivo y el morfismo canónico $r:X_s\to\mathbf{P}((f_s)_*(\sh{L}_s))$ no es una inmersión.
Teniendo en cuenta el isomorfismo \sref{IV.9.6.3.1}, estas condiciones se escriben respectivamente de la forma: 1$^\circ$ el homomorfismo canónico $(f^*(f_*(\sh{L})))_s\to\sh{L}_s$ no es sobreyectivo; 2$^\circ$ el homomorfismo precedente es sobreyectivo y el morfismo canónico $X_s\to\mathbf{P}((f_*(\sh{L}))_s)$ no es una inmersión.

Supongamos primero que el homomorfismo canónico $(f^*(f_*(\sh{L})))_\eta\to\sh{L}_\eta$ no es sobreyectivo.
Como $f_*(\sh{L})$ es coherente, también lo es $f^*(f_*(\sh{L}))$, y se sigue entonces de \sref{IV.9.4.5} que para todo $s$ en un entorno de $\eta$, el homomorfismo $(f^*(f_*(\sh{L})))_s\to\sh{L}_s$ no es sobreyectivo, lo que demuestra en este caso que $S-E'$ es un entorno de $\eta$.

Supongamos ahora que el homomorfismo canónico $(f^*(f_*(\sh{L})))_\eta\to\sh{L}_\eta$ es sobreyectivo, pero que el morfismo $X_\eta\to\mathbf{P}((f_*(\sh{L}))_\eta)$ no es una inmersión.
Entonces el mismo razonamiento anterior muestra que para $s$ suficientemente próximo a $\eta$, el homomorfismo $(f^*(f_*(\sh{L})))_s\to\sh{L}_s$ es sobreyectivo; por otra parte, en virtud de \sref{IV.9.6.1}, (viii), el morfismo $X_s\to\mathbf{P}((f_*(\sh{L}))_s)$ no es una inmersión para $s$ en un entorno de $\eta$, lo que termina la demostración en el caso de $E'$.

\emph{Caso de $E$.} --- Consideremos primero el caso particular en que $Y=S$:

\begin{corollary}[9.6.4]
\label{IV.9.6.4}
Sea $f:X\to S$ un morfismo propio de presentación finita y sea $\sh{L}$ un $\sh{O}_X$-Módulo invertible.
Entonces el conjunto $E$ de los $s\in S$ tales que $\sh{L}_s$ es amplio (relativamente a $f_s$) es abierto en $S$, y $\sh{L}|f^{-1}(E)$ es amplio relativamente a la restricción $f^{-1}(E)\to E$ de $f$.
\end{corollary}

\begin{proof}
Como la condición \sref{IV.9.2.1},~(i) es verificada por la propiedad $\mathbf{P}$ definida anteriormente, el resultado de \sref{IV.9.2.2},~(iv) y el razonamiento de \sref{IV.9.2.3} muestran que basta reducirse al caso en que $S$ es \emph{noetheriano}; pero entonces el resultado se sigue de \sref[III]{III.4.7.1} y de la estabilidad de la amplitud por cambio de base \sref[II]{II.4.6.13}.
\end{proof}

\begin{corollary}[9.6.5]
\label{IV.9.6.5}
Bajo las hipótesis de \sref{IV.9.6.4}, para que $\sh{L}$ sea amplio relativamente a $f$, es necesario y suficiente que, para todo $s\in S$, $\sh{L}_s$ sea amplio relativamente a $f_s$.
\end{corollary}

\phantomsection\label{IV.9.6.6}%
\textbf{(9.6.6) Fin de la demostración de \sref{IV.9.6.3}.} ---
Volvamos al caso general, siendo $S$ noetheriano, íntegro, de punto genérico $\eta\in S-E$.
Como $f$ es propio, se sigue de \sref{IV.9.6.4} que el conjunto $V$ de los $y\in Y$ tales que $\sh{L}_y$ es un $\sh{O}_{X_y}$-Módulo amplio, relativamente al morfismo $f_y:X_y\to\operatorname{Spec}(\boldsymbol{k}(y))$, es abierto, y por tanto $F=Y-V$ es cerrado en $Y$.
Como $f_s$ es propio y, para todo $y\in Y$ situado sobre $s\in S$, se tiene
$\sh{L}_y=(\sh{L}_s)_y$, se sigue de \sref{IV.9.6.5} que $s$ pertenece a
$S-E$ si y solo si $F_s\neq\varnothing$.
Pero $F$ es cerrado en $Y$ y $F_\eta\neq\varnothing$; por tanto, \sref{IV.9.5.1} muestra
que el conjunto de los $s\in S$ tales que $F_s\neq\varnothing$ es un entorno de
$\eta$ en $S$.
\end{proof}

\begin{remarks}[9.6.7]
\label{IV.9.6.7}
(i) Para cada una de las propiedades $\mathbf{P}$ consideradas en \sref{IV.9.6.1}, es aplicable la proposición \sref{IV.9.3.3}, y estas propiedades (para los morfismos $f_s$) son, por tanto, ``estables'' al pasar del límite proyectivo de un sistema proyectivo esencialmente afín \sref{IV.8.13.4} de preesquemas $X_\lambda$ a uno conveniente de ellos.

(ii) Sea $Z$ una parte localmente construible de $X$ tal que, para todo $s\in S$, $Z\cap X_s$ sea \emph{abierto} en $X_s$, y denotemos por $Z_s$ el subpreesquema de $X_s$ inducido sobre el abierto $Z\cap X_s$.
\oldpage[IV-3]{76}
Entonces, en las proposiciones \sref{IV.9.6.1} y \sref{IV.9.6.2},~(I), podemos sustituir en todas partes $f_s$ por su restricción $f_s|Z_s:Z_s\to Y_s$ sin modificar las conclusiones.
En efecto, la verificación de \sref{IV.9.2.1},~(i) procede como en \sref{IV.9.6.1} y \sref{IV.9.6.2}.
Por otra parte, en la reducción al caso noetheriano, efectuada en \sref{IV.9.2.3}, si $Z=q^{-1}(Z_0)$, donde $q:X\to X_0$ es la proyección canónica y $Z_0$ una parte construible de $X_0$ \sref{IV.8.3.11}, el hecho de que $(Z_0)_{s_0}$ sea abierto en $(X_0)_{s_0}$, para $s_0=p(s)$, se sigue de \sref{IV.2.4.10} y del hecho de que la proyección $X_s\to(X_0)_{s_0}$ es sobreyectiva.
Nos reducimos, por tanto, a verificar \sref{IV.9.2.1},~(ii) bajo las nuevas hipótesis.
Ahora bien, como $Z_\eta$ es abierto en $X_\eta$, existe un abierto $Z'\subset X$ tal que $Z_\eta=Z'\cap X_\eta$; como $X$ es entonces noetheriano, $Z'$ es construible, y $Z$ también es construible por hipótesis; concluimos de \sref{IV.9.5.2} y de ($\mathbf{0}_{\mathrm{III}}$,~9.2.2) que existe un entorno $U$ de $\eta$ en $S$ tal que $Z_s=Z'_s$ para $s\in U$.
Sustituyendo $S$ por $U$, podemos, por tanto, reducirnos al caso en que $Z$ es abierto en $X$, y entonces el resultado ha sido demostrado en \sref{IV.9.6.1} y \sref{IV.9.6.2}.
\end{remarks}

\subsection{Construibilidad de las propiedades de separabilidad, irreducibilidad geométrica y conexidad geométrica}
\label{subsection:IV.9.7}
\label{IV.9.7}

\begin{lemma}[9.7.1]
\label{IV.9.7.1}
Sea $S$ un preesquema irreducible de punto genérico $\eta$, y sea $f:X\to S$ un morfismo de presentación finita.
Si $X_\eta$ tiene $n$ componentes irreducibles (resp.\ $n'$ componentes conexas), existe un entorno abierto $U$ de $\eta$ en $S$ tal que, para todo $s\in U$, $X_s$ tiene al menos $n$ componentes irreducibles (resp.\ $n'$ componentes conexas).
\end{lemma}

\begin{proof}
Podemos reducirnos al caso en que $S$ es afín, de modo que $X$ es cuasicompacto y cuasiseparado.
Sean $V_i$ ($1\leqslant i\leqslant n$) los interiores de las componentes irreducibles de $X_\eta$, que son dos a dos disjuntos y cuasicompactos, y cuya unión es densa en $X_\eta$; en virtud de \sref{IV.8.2.11}, aplicado según el método de \sref{IV.8.1.2},~\emph{a}), existe para cada $i$ un abierto cuasicompacto $W_i$ en $X$ tal que $W_i\cap X_\eta=V_i$; como $X$ es cuasiseparado, las intersecciones $W_i\cap W_j$ son cuasicompactas \sref{IV.1.2.7}, y por tanto, sustituyendo $S$ por un entorno de $\eta$, podemos suponer que $W_i\cap W_j=\varnothing$ para $i\neq j$ en virtud de \sref{IV.8.3.3}.
Además, como los $W_i$ son construibles, existe un entorno $U$ de $\eta$ en $S$ tal que para todo $s\in U$ los $(W_i)_s$ son no vacíos \sref{IV.9.5.1} y \sref{IV.9.2.3} y la unión de los $(W_i)_s$ es densa en $X_s$ \sref{IV.9.5.3} y \sref{IV.9.2.3}.
En estas condiciones, como las componentes irreducibles de cada $(W_i)_s$ (que son en número finito) son también componentes irreducibles de la unión de los $(W_i)_s$ \sref[0]{0.2.1.7}, las adherencias en $X_s$ de estas componentes son las componentes irreducibles de $X_s$ \sref[0]{0.2.1.6} y su número es evidentemente $\geqslant n$.

Sean ahora $C_j$ ($1\leqslant j\leqslant n'$) las componentes conexas de $X_\eta$, que son abiertos cuasicompactos de $X_\eta$, dos a dos disjuntos.
Si sustituimos $V_i$ por $C_j$ en el razonamiento precedente, vemos (utilizando dos veces \sref{IV.8.3.3}) que podemos suponer $X$ una unión de $n'$ abiertos cuasicompactos $M_j$ ($1\leqslant j\leqslant n'$), dos a dos disjuntos, y tales que, en un entorno de $\eta$, los $(M_j)_s$ son no vacíos.
Como la unión de los $(M_j)_s$ es $X_s$, las componentes conexas de los $(M_j)_s$ son las componentes conexas de $X_s$, y por tanto su número es $\geqslant n'$.
\end{proof}

\oldpage[IV-3]{77}

\begin{lemma}[9.7.2]
\label{IV.9.7.2}
Sea $S$ un preesquema irreducible de punto genérico $\eta$, y sea $f:X\to S$ un morfismo de presentación finita.
Si $X_\eta$ no es reducido, existe un entorno abierto $U$ de $\eta$ en $S$ tal que, para todo $s\in U$, $X_s$ no es reducido.
\end{lemma}

\begin{proof}
En efecto, sea $\mathscr{N}$ el nilradical de $\sh{O}_X$; se sigue de \sref{IV.8.2.13}, aplicado según el método de \sref{IV.8.1.2},~\emph{a}), que $\mathscr{N}_\eta$ es el nilradical de $\sh{O}_{X_\eta}$, y la hipótesis es que $\mathscr{N}_\eta\neq 0$.
Concluimos de \sref{IV.9.4.5} y \sref{IV.9.2.3} que existe un entorno $U$ de $\eta$ tal que, para todo $s\in U$, $\mathscr{N}_s$ se identifica con un ideal de $\sh{O}_{X_s}$ y $\mathscr{N}_s\neq 0$; como $\mathscr{N}_s$ está evidentemente contenido en el nilradical de $\sh{O}_{X_s}$, vemos que $X_s$ no es reducido para $s\in U$.
\end{proof}

\textbf{(9.7.3)}
\label{IV.9.7.3}
Dado un polinomio $F\in A[\mathrm{T}_1,\ldots,\mathrm{T}_n]$, donde $A$ es un anillo y los $\mathrm{T}_i$ son indeterminados, para un homomorfismo de anillos $\rho:A\to B$, denotamos por $F_{(\rho)}$ o $F_{(B)}$ el polinomio de $B[\mathrm{T}_1,\ldots,\mathrm{T}_n]$ obtenido sustituyendo cada coeficiente de $F$ por su imagen mediante $\rho$.
Si $k$ es un cuerpo, $F\in k[\mathrm{T}_1,\ldots,\mathrm{T}_n]$ un polinomio no constante y $X=\operatorname{Spec}(k[\mathrm{T}_1,\ldots,\mathrm{T}_n]/(F))$, decir que $X$ es íntegro (o que el ideal $(F)$ es primo) significa que $F$ es \emph{irreducible} (es decir, que en toda factorización $F=F_1 F_2$ en polinomios de $k[\mathrm{T}_1,\ldots,\mathrm{T}_n]$, $F_1$ o $F_2$ es de grado $0$); esto se sigue del hecho de que el anillo de polinomios $k[\mathrm{T}_1,\ldots,\mathrm{T}_n]$ es factorial.
Deducimos inmediatamente de \sref{IV.4.6.2} y \sref{IV.4.5.2}:

\begin{lemma}[9.7.4]
\label{IV.9.7.4}
Sea $k$ un cuerpo, $\Omega$ una extensión algebraicamente cerrada de $k$ y $F$ un polinomio no constante de $k[\mathrm{T}_1,\ldots,\mathrm{T}_n]$.
Las condiciones siguientes son equivalentes:
\begin{enumerate}[label=\emph{\alph*)}]
  \item $X=\operatorname{Spec}(k[\mathrm{T}_1,\ldots,\mathrm{T}_n]/(F))$ es geométricamente íntegro.
  \item $F_{(K)}$ es irreducible para toda extensión $K$ de $k$.
  \item $F_{(\Omega)}$ es irreducible.
\end{enumerate}
Diremos en este caso que $F$ es \emph{geométricamente irreducible}.
\end{lemma}

\begin{lemma}[9.7.5]
\label{IV.9.7.5}
Sea $A$ un anillo íntegro, $K$ el cuerpo de fracciones de $A$ y $F$ un polinomio no constante de $A[\mathrm{T}_1,\ldots,\mathrm{T}_n]$ de grado $d$ tal que $F_{(K)}$ es geométricamente irreducible.
Entonces existe $f\neq 0$ en $A$ tal que para todo $x\in D(f)$, $F_{(\boldsymbol{k}(x))}$ es geométricamente irreducible.
\end{lemma}

\begin{proof}
Pongamos $F=\sum_\alpha c_\alpha \mathrm{T}^\alpha$ como de costumbre, con $c_\alpha=c_{\alpha_1\alpha_2\ldots\alpha_n}$, $\mathrm{T}^\alpha=\mathrm{T}_1^{\alpha_1}\mathrm{T}_2^{\alpha_2}\cdots\mathrm{T}_n^{\alpha_n}$, $|\alpha|=\alpha_1+\alpha_2+\cdots+\alpha_n\leqslant d$ (siendo no nulo al menos uno de los $c_\alpha$ tales que $|\alpha|=d$).
Como los $c_\alpha$ no nulos son invertibles en $K$, podemos suponer, sustituyendo $A$ por $A_g$ ($g\neq 0$ en $A$), que los $c_\alpha$ no nulos son invertibles \emph{en} $A$.
Se sigue entonces que, para todo $x\in\operatorname{Spec}(A)$, $F_{(\boldsymbol{k}(x))}$ es de grado $d$.

Demostraremos el lema mediante un lema preliminar.

\begin{lemma}[9.7.5.1]
\label{IV.9.7.5.1}
Sea $A$ un anillo íntegro, $S=\operatorname{Spec}(A)$, $\eta$ el punto genérico de $S$, $B$ el anillo de polinomios $A[\mathrm{T}_1,\ldots,\mathrm{T}_m]$, $(P_i)_{i\in I}$ una familia finita de elementos de $B$ y $\mathfrak{a}$ el ideal de $B$ generado por los $P_i$.
Para todo $s\in S$, sea $\mathfrak{a}_s$ el ideal de $\boldsymbol{k}(s)[\mathrm{T}_1,\ldots,\mathrm{T}_m]$ generado por los polinomios $(P_i)_{(\boldsymbol{k}(s))}$ ($i\in I$).
Entonces, si $V(\mathfrak{a}_\eta)=\varnothing$, existe un entorno $U$ de $\eta$ en $S$ tal que $V(\mathfrak{a}_s)=\varnothing$ para todo $s\in U$.
\end{lemma}

\begin{proof}
En efecto, sea $Y=\operatorname{Spec}(B)$, y sea $Z$ la parte cerrada $V(\mathfrak{a})$ de $Y$; como $\mathfrak{a}$ es un ideal de tipo finito en $B$, $Z$ es construible en $Y$ \sref[0]{0.9.1.5} y, con las notaciones introducidas en \sref{IV.9.4.1}, tenemos $V(\mathfrak{a}_s)=Z_s$ para todo $s\in S$; como el morfismo estructural $Y\to S$ es de presentación finita, la conclusión del lema se sigue de \sref{IV.9.5.1} y \sref{IV.9.2.3}.
\end{proof}

\oldpage[IV-3]{78}
Sea $(p,q)$ un par de enteros $>0$ tales que $p+q=d$; introducimos indeterminados $\mathrm{T}'_\beta$, $\mathrm{T}''_\gamma$ para todos los sistemas de enteros $\beta$, $\gamma$ tales que $|\beta|\leqslant p$ y $|\gamma|\leqslant q$; para todo sistema de enteros $\alpha$ tal que $|\alpha|\leqslant d$, consideremos el polinomio de $B=A[\mathrm{T}'_\beta, \mathrm{T}''_\gamma \mid |\beta|\leqslant p, |\gamma|\leqslant q]: P_\alpha(\mathrm{T}'_\beta,\mathrm{T}''_\gamma)=\sum_{\beta+\gamma=\alpha}\mathrm{T}'_\beta \mathrm{T}''_\gamma-c_\alpha$.
Sea $\Omega$ una clausura algebraica de $K$; decir que existen dos polinomios $F_1$, $F_2$ de $\Omega[\mathrm{T}_1,\ldots,\mathrm{T}_n]$, de grados respectivos $p$ y $q$, y tales que $F_1 F_2=F_{(\Omega)}$ significa que el sistema de ecuaciones $P_\alpha(t'_\beta,t''_\gamma)=0$ ($|\alpha|\leqslant d$) admite una solución $(t'_\beta,t''_\gamma)$ ($|\beta|\leqslant p$, $|\gamma|\leqslant q$) formada por elementos de $\Omega$.
Sea $\mathfrak{a}$ el ideal de $B$ generado por los $P_\alpha$; la interpretación precedente y el teorema de los ceros de Hilbert muestran que la hipótesis sobre $F_{(K)}$ implica que $V(\mathfrak{a}_\eta)=\varnothing$, designando por $\eta$ el punto genérico de $\operatorname{Spec}(A)$; el lema \sref{IV.9.7.5.1} demuestra entonces que en un entorno de $\eta$, tenemos $V(\mathfrak{a}_s)=\varnothing$, y por consiguiente, para estos valores de $x$, $F_{(\boldsymbol{k}(x))}$ no admite ninguna factorización $F_{(\boldsymbol{k}(x))}=G_1 G_2$, donde $G_1$, $G_2$ son polinomios de grados respectivos $p$ y $q$, cuyos coeficientes pertenecen a una clausura algebraica de $\boldsymbol{k}(x)$.
Basta aplicar este resultado a todos los pares de enteros $(p,q)$ tales que $p>0$, $q>0$ y $p+q=d$, para obtener la conclusión del lema \sref{IV.9.7.5}.
\end{proof}

\begin{proposition}[9.7.6]
\label{IV.9.7.6}
Sea $S$ un preesquema noetheriano íntegro de punto genérico $\eta$, sea $f:X\to S$ un morfismo de tipo finito y sea $\sh{F}$ un $\sh{O}_X$-Módulo coherente.
Si $\sh{F}_\eta$ no tiene ningún ciclo primo inmerso asociado, existe un entorno $U$ de $\eta$ en $S$ tal que, para todo $s\in U$, $\sh{F}_s$ no tiene ningún ciclo primo inmerso asociado.
\end{proposition}

\begin{proof}
Como $X_\eta$ es noetheriano, existe un homomorfismo inyectivo $v:\sh{F}_\eta\to\bigoplus_{i=1}^m\sh{G}_i$, donde cada $\sh{G}_i$ es irredundante y $\operatorname{Ass}(\sh{F}_\eta)$ es el conjunto de los $x_i$, con $\{x_i\}=\operatorname{Ass}(\sh{G}_i)$ \sref{IV.3.2.6}; si $Z_i$ es la adherencia de $\{x_i\}$ en $X_\eta$, entonces $Z_i=\operatorname{Supp}(\sh{G}_i)$ \sref{IV.3.1.4}, y por hipótesis los $Z_i$ son las componentes irreducibles de $Z=\operatorname{Supp}(\sh{F}_\eta)$.
Se sigue de \sref{IV.8.5.2},~(i) y (ii) (restringiendo $S$ si es necesario a un entorno de $\eta$) y de \sref{IV.8.5.8} que existen $\sh{O}_X$-Módulos coherentes $\sh{F}_i$ y un homomorfismo inyectivo $u:\sh{F}\to\bigoplus_{i=1}^m\sh{F}_i$ tales que $(\sh{F}_i)_\eta=\sh{G}_i$ para todo $i$ y $v=u_\eta$.
Si $Y_i=\operatorname{Supp}(\sh{F}_i)$, tenemos $(Y_i)_s=\operatorname{Supp}((\sh{F}_i)_s)$ para todo $s\in S$ \sref[I]{I.9.1.13} y, en particular, $(Y_i)_\eta=Z_i$ para todo $i$.
La hipótesis implica que para $i\neq j$, $Z_i-(Z_i\cap Z_j)$ es abierto y denso en $Z_i$; deducimos, por tanto, de \sref{IV.9.5.3} y \sref{IV.9.5.4} que en un entorno de $\eta$, $(Y_i)_s-((Y_i)_s\cap(Y_j)_s)$ es abierto y denso en $(Y_i)_s$.
Supongamos demostrado que cada uno de los $(\sh{F}_i)_s$ carece de ciclo primo inmerso asociado para $s\in U$.
Entonces los elementos de $\operatorname{Ass}((\sh{F}_i)_s)$ son los puntos maximales de $(Y_i)_s$; ninguno de ellos puede, por tanto, pertenecer a $(Y_j)_s$ para $i\neq j$, y la proposición quedará demostrada. Podemos, pues, suponer que $\sh{F}_\eta$ es \emph{irredundante}; además, podemos reducirnos al caso en que $S=\operatorname{Spec}(A)$ es afín. Entonces $X$ es una unión de un número finito de abiertos afines $V_j$, y si $V_j\cap\operatorname{Supp}(\sh{F}_\eta)=\varnothing$, también tenemos $V_j\cap\operatorname{Supp}(\sh{F}_s)=\varnothing$ para $s$ próximo a $\eta$ \sref{IV.9.5.1}. Podemos, por tanto, reducirnos al caso en que $X=\operatorname{Spec}(B)$ es también afín y $f:X\to S$ es dominante. Así, $A$ es un anillo noetheriano íntegro, subanillo de una $A$-álgebra de tipo finito $B$, y $\sh{F}=\widetilde{M}$ con $M$ un $B$-módulo finito. Por hipótesis, si $K$ es el cuerpo de fracciones de $A$, el $B_{(K)}$-módulo $M_{(K)}$ es irredundante.
Sea $\mathfrak{q}$ el único

\oldpage[IV-3]{79}

elemento de $\operatorname{Ass}(M_{(K)})$, y sea $\mathfrak{p}$ el ideal primo de $B$ que es la imagen inversa de $\mathfrak{q}$ por la aplicación canónica $B\to B_{(K)}$.
Sabemos por \sref{IV.5.11.1.1} que existe una filtración finita $(N_h)_{0\leqslant h\leqslant m}$ de $M_{(K)}$ tal que $N_0=M_{(K)}$, $N_m=0$, y $N_h/N_{h+1}$ es isomorfo a un $B_{(K)}$-submódulo no nulo de $B_{(K)}/\mathfrak{q}$.
Sea $M_h$ la imagen inversa de $N_h$ por la aplicación canónica $M\to M_{(K)}$.
Se sigue de \sref{IV.9.4.4} que para $s$ suficientemente próximo a $\eta$, $(M_h)_{(\boldsymbol{k}(s))}$ se identifica con un sub-$B_{(\boldsymbol{k}(s))}$-módulo de $M_{(\boldsymbol{k}(s))}$, y el cociente $(M_h)_{(\boldsymbol{k}(s))}/(M_{h+1})_{(\boldsymbol{k}(s))}$ con un $B_{(\boldsymbol{k}(s))}$-submódulo no nulo de $B_{(\boldsymbol{k}(s))}/\mathfrak{p}_{(\boldsymbol{k}(s))}=(B/\mathfrak{p})_{(\boldsymbol{k}(s))}$.
Teniendo en cuenta \sref{IV.3.1.7}, con las mismas notaciones, queda demostrar que si $B$ es \emph{íntegro}, entonces $B_{(\boldsymbol{k}(s))}$ no tiene ideales primos asociados inmersos para $s$ próximo a $\eta$.
Ahora bien, sustituyendo $A$ por $A_g$ y $B$ por $B_g$ (donde $g$ es un elemento no nulo de $A$), podemos suponer que $B$ contiene una subálgebra de polinomios $C=A[\mathrm{T}_1,\ldots,\mathrm{T}_n]$ tal que $B$ es un $C$-módulo de tipo finito (Bourbaki, \emph{Alg.\ comm.}, cap.\ V, \S~3, n$^\circ$~1, cor.\ 1 del teorema\ 1).
Como $B_{(K)}$ es un $C_{(K)}$-módulo sin torsión, podemos aplicar una vez más el argumento precedente, sustituyendo $B$, $M$ y $\mathfrak{q}$ por $C$, $B$ y $(0)$ respectivamente; basta demostrar que para $s$ próximo a $\eta$, $C_{(\boldsymbol{k}(s))}$ no tiene ideales primos asociados inmersos.
Pero esto es evidente, puesto que $C_{(\boldsymbol{k}(s))}=\boldsymbol{k}(s)[\mathrm{T}_1,\ldots,\mathrm{T}_n]$ es un anillo íntegro.
\end{proof}

\begin{theorem}[9.7.7]
\label{IV.9.7.7}
Sea $f:X\to S$ un morfismo de presentación finita, y sea $E$ el conjunto de los $s\in S$ para los cuales $X_s$ posee una de las propiedades siguientes: ser
\begin{enumerate}[label=(\roman*)]
  \item geométricamente irreducible;
  \item geométricamente conexo;
  \item geométricamente reducido;
  \item geométricamente íntegro.
\end{enumerate}
Entonces $E$ es localmente construible en $S$.
\end{theorem}

\begin{proof}
Para ver que las propiedades consideradas son construibles, observemos primero que satisfacen trivialmente la condición \sref{IV.9.2.1},~(i), en virtud de \sref{IV.4.5.6},~(i) y \sref{IV.4.6.5},~(i).
Queda, por tanto, verificar \sref{IV.9.2.1},~(ii), de modo que podemos suponer $S=\operatorname{Spec}(A)$ afín, noetheriano, íntegro, de punto genérico $\eta$.
En virtud de \sref{IV.4.6.8}, existe una extensión finita $K'$ de $K=\boldsymbol{k}(\eta)$ tal que las componentes irreducibles (resp.\ componentes conexas) de $(X_\eta)_{(K')}$ son geométricamente irreducibles (resp.\ geométricamente conexas) y $((X_\eta)_{(K')})_\mathrm{red}$ es geométricamente reducido.
Además, existe una base de $K'$ sobre $K$ formada por elementos enteros sobre $A$; el anillo íntegro $A'$ generado por estos elementos es finito sobre $A$ y tiene $K'$ como cuerpo de fracciones.
Pongamos $S'=\operatorname{Spec}(A')$; el morfismo $g:S'\to S$ es finito y dominante, por tanto sobreyectivo \sref[II]{II.6.1.10}. Pongamos $X'=X_{(S')}$, de modo que, si $\eta'$ es el punto genérico de $S'$, entonces $X'_{\eta'}=(X_\eta)_{(K')}$. El conjunto $E'$ de los $s'\in S'$ tales que $X'_{s'}$ posee una de las propiedades (i), (ii), (iii) o (iv) es igual a $g^{-1}(E)$ \sref{IV.9.2.2},~(iv), donde $E$ corresponde, naturalmente, a la misma propiedad. Como $g$ es sobreyectivo, $E=g(E')$ y $S-E=g(S'-E')$. Además, $g$ es cerrado y $g^{-1}(\eta)=\{\eta'\}$ porque $A'$ es íntegro y finito sobre $A$ (Bourbaki, \emph{Alg.\ comm.}, cap.\ V, \S~2, n$^\circ$~1, cor.\ 1 de la prop.\ 1); por tanto, la imagen por $g$ de todo entorno de $\eta'$ es un entorno de $\eta$.
El teorema quedará, pues, demostrado si se prueba que $E'$ o $S'-E'$ es un entorno de $\eta'$.
En otras palabras, podemos suponer en lo sucesivo que las componentes irreducibles (resp.\ componentes conexas) de $X_\eta$ son geométricamente

\oldpage[IV-3]{80}

irreducibles (resp.\ geométricamente conexas) y que $(X_\eta)_\mathrm{red}$ es geométricamente reducido.

Supongamos primero que $\eta\in S-E$ para una de las propiedades (i) a (iv).
Si $X_\eta$ no es geométricamente irreducible (resp.\ geométricamente conexo), no es irreducible (resp.\ conexo) en virtud de la hipótesis precedente, y por tanto lo mismo vale para $X_s$ para $s$ en un entorno de $\eta$ \sref{IV.9.7.1}, y \emph{a fortiori} en este entorno $X_s$ no es geométricamente irreducible (resp.\ geométricamente conexo).
Por otra parte, si $X_\eta$ no es geométricamente reducido, no es reducido (pues de otro modo sería igual a $(X_\eta)_\mathrm{red}$, que es geométricamente reducido por hipótesis); por tanto, en un entorno de $\eta$, $X_s$ no es reducido \sref{IV.9.7.2} y, en consecuencia, no es geométricamente reducido.
Finalmente, si $X_\eta$ no es geométricamente íntegro, o bien no es reducido, en cuyo caso acabamos de ver que $X_s$ no es reducido (ni \emph{a fortiori} íntegro) en un entorno de $\eta$; o bien $X_\eta$ es reducido (por tanto geométricamente reducido por hipótesis), en cuyo caso no es geométricamente irreducible, y acabamos de ver más arriba que lo mismo vale entonces para $X_s$ para $s$ próximo a $\eta$; \emph{a fortiori} $X_s$ no es geométricamente íntegro para estos valores de $s$.

Supongamos ahora que $\eta\in E$ y examinemos por separado cada una de las propiedades consideradas.

$1^\circ$ Supongamos que $X_\eta$ es geométricamente íntegro.
Sea $L$ el cuerpo de funciones racionales sobre $X_\eta$; la hipótesis sobre $X_\eta$ implica que $L$ es una extensión separable de $K$ \sref{IV.4.6.3}, por tanto una extensión separable finita de una extensión pura $K(\mathrm{T}_1,\ldots,\mathrm{T}_n)$ (siendo los $\mathrm{T}_i$ indeterminados); existe entonces un elemento $x\in L$, entero sobre el anillo $K[\mathrm{T}_1,\ldots,\mathrm{T}_n]$, y tal que $L=K(\mathrm{T}_1,\ldots,\mathrm{T}_n)(x)$; sea $G\in K[\mathrm{T}_1,\ldots,\mathrm{T}_n,\mathrm{T}_{n+1}]$ su polinomio minimal.
Existe $g\neq 0$ de $A$ tal que todos los coeficientes $\neq 0$ de $G$ (que pertenecen todos a $K$) pertenecen a $A_g$; sustituyendo $A$ por $A_g$ (lo que equivale a sustituir $S$ por un entorno de $\eta$), podemos, por tanto, suponer que $G$ tiene sus coeficientes en $A$; designando por $F$ el polinomio $G$ considerado como elemento de $A[\mathrm{T}_1,\ldots,\mathrm{T}_n,\mathrm{T}_{n+1}]$, tenemos entonces $G=F_{(\boldsymbol{k}(\eta))}$.
Pongamos $Y=\operatorname{Spec}(A[\mathrm{T}_1,\ldots,\mathrm{T}_{n+1}]/(F))$, de modo que $Y_\eta=\operatorname{Spec}(K[\mathrm{T}_1,\ldots,\mathrm{T}_{n+1}]/(G))$; $Y_\eta$ es un esquema íntegro cuyo cuerpo de funciones racionales es $L$.
Como $X_\eta$ e $Y_\eta$ son noetherianos, existe un abierto no vacío $V\subset Y_\eta$ y una inmersión abierta (necesariamente dominante) $v:V\to X_\eta$ \sref[I]{I.6.5.1},~(ii) y \sref[I]{I.6.5.4},~(ii).
Sea $W$ un abierto de $Y$ tal que $W\cap Y_\eta=V$; aplicando \sref{IV.8.8.2},~(i) y \sref{IV.8.10.5},~(iii) según el método de \sref{IV.8.1.2},~\emph{a}), sustituyendo $S$ por un entorno de $\eta$, podemos suponer que $v=w_\eta$, donde $w:W\to X$ es una inmersión abierta.

Vimos en \sref{IV.4.6.3} que el criterio para que un preesquema algebraico íntegro sea geométricamente íntegro depende únicamente de su cuerpo de funciones racionales. Como $X_\eta$ es geométricamente íntegro por hipótesis, lo mismo vale para $Y_\eta$, y la definición de $G$ implica, por tanto, que este polinomio es geométricamente irreducible \sref{IV.9.7.4}.
Aplicando \sref{IV.9.7.5}, vemos que existe un entorno $U$ de $\eta$ en $S$ tal que $F_{(\boldsymbol{k}(s))}$ es geométricamente irreducible para todo $s\in U$, y por tanto $Y_s$ es geométricamente íntegro para $s\in U$ \sref{IV.9.7.4}; además, podemos también suponer que para $s\in U$, $W_s$ no es vacío \sref{IV.9.5.1}, y es por tanto geométricamente íntegro \sref{IV.4.6.3}; finalmente, podemos también suponer que

\oldpage[IV-3]{81}

para $s\in U$, $w_s:W_s\to X_s$ es una inmersión abierta \sref{IV.9.6.1},~(x), y que su imagen en $X_s$ es densa en todas partes \sref{IV.9.5.3}.
En otras palabras, para $s\in U$, existe en $X_s$ un abierto denso en todas partes $W'_s$ que es geométricamente íntegro; el criterio \sref{IV.4.5.9},~\emph{c}) muestra, por tanto, que $X_s$ es ya geométricamente irreducible para $s\in U$.
Finalmente, como $X_\eta$ es reducido y por tanto no tiene ningún ciclo primo inmerso asociado \sref{IV.3.2.1}, podemos también suponer que para $s\in U$, $X_s$ no tiene ningún ciclo primo inmerso asociado \sref{IV.9.7.6}; en estas condiciones (para $s\in U$ fijado), sea $\Omega$ una extensión algebraicamente cerrada de $\boldsymbol{k}(s)$, y sea $p:(X_s)_{(\Omega)}\to X_s$ la proyección canónica; $p^{-1}(W'_s)$ es un abierto denso en $(X_s)_{(\Omega)}$ \sref{IV.2.3.10} y es íntegro por hipótesis; además, $(X_s)_{(\Omega)}$ no tiene ningún ciclo primo inmerso asociado \sref{IV.4.2.7}, por tanto concluimos de \sref{IV.3.2.1} que $(X_s)_{(\Omega)}$ es reducido; esto termina la demostración de que $X_s$ es geométricamente íntegro \sref{IV.4.6.1}.

$2^\circ$ Supongamos que $X_\eta$ es geométricamente irreducible; como $X_{\mathrm{red}}$ es también de tipo finito sobre $S$ \sref{IV.1.5.4},~(vi), teniendo en cuenta \sref[I]{I.5.1.8}, podemos sustituir $X$ por $X_{\mathrm{red}}$; entonces $X_\eta$ es también íntegro, y como por hipótesis $X_\eta$ es geométricamente reducido, es geométricamente íntegro. Estamos entonces en las condiciones de $1^\circ$, y concluimos (volviendo a las hipótesis iniciales) que $X_s$ es geométricamente irreducible para $s$ próximo a $\eta$.

$3^\circ$ Supongamos $X_\eta$ geométricamente conexo, y sean $Z_i$ ($1\leqslant i\leqslant n$) las componentes irreducibles de $X_\eta$; existe (en virtud de \sref{IV.5.10.8.1} aplicado a $\mathfrak{T}=\{\varnothing\}$) una aplicación \emph{sobreyectiva} $j\mapsto\nu(j)$ de un intervalo $[1,m]$ de $\mathbf{N}$ sobre $[1,n]$ tal que $Z_{\nu(j)}\cap Z_{\nu(j+1)}\neq\varnothing$ para $1\leqslant j\leqslant m$.
Para todo $i$, sea $X_i$ la adherencia de $Z_i$ en $X$, y sea $Y$ la unión de los $X_i$; como $Y_\eta=X_\eta$ por definición, podemos, en virtud de \sref{IV.9.5.1}, suponer que $Y_s=X_s$ para todo $s\in S$, de modo que $X_s$ es la unión de los $(X_i)_s$.
Ahora bien, considerando los subpreesquemas reducidos de $X$ cuyos espacios subyacentes son los $X_i$, vemos por $2^\circ$ que existe un entorno $U$ de $\eta$ en $S$ tal que para $s\in U$ los $(X_i)_s$ son \emph{geométricamente irreducibles} (puesto que podemos suponer los $(X_i)_\eta$ geométricamente irreducibles, como vimos al principio).
Además, podemos también suponer que para $s\in U$, $(X_{\nu(j)})_s\cap(X_{\nu(j+1)})_s\neq\varnothing$ \sref{IV.9.5.1} para $1\leqslant j\leqslant m$; concluimos inmediatamente que $X_s$ es conexo, por tanto, por \sref{IV.4.5.13.1}, geométricamente conexo para $s\in U$.

$4^\circ$ Supongamos $X_\eta$ geométricamente reducido; sean $Z_i$ las componentes irreducibles de $X_\eta$, y sea $W_i$ el interior de $Z_i$ en $X_\eta$; existe para cada $i$ un abierto $V_i$ de $X$ tal que $W_i=V_i\cap X_\eta$ para todo $i$; como los $W_i$ son abiertos, dos a dos disjuntos, y su unión es densa en $X_\eta$, \sref{IV.9.5.1},~\sref{IV.9.5.3} y \sref{IV.9.5.4} nos permiten suponer que para $s$ próximo a $\eta$, los $(V_i)_s$ son abiertos, dos a dos disjuntos, en $X_s$ y su unión es densa en $X_s$.
Además, como los $W_i$ son geométricamente reducidos y se supusieron al principio geométricamente irreducibles, se sigue de $1^\circ$ que para $s$ próximo a $\eta$, los $(V_i)_s$ son geométricamente íntegros y, \emph{a fortiori}, reducidos.
Por otra parte, deducimos de \sref{IV.9.7.6} que para $s$ próximo a $\eta$, $X_s$ no tiene ningún ciclo primo inmerso asociado, puesto que lo mismo vale para el preesquema reducido $X_\eta$ \sref{IV.3.2.1}; concluimos, por tanto, de \sref{IV.3.2.1} que $X_s$ es reducido, y de \sref{IV.4.6.1} que $X_s$ es geométricamente reducido.
\end{proof}

\oldpage[IV-3]{82}

Las partes del enunciado de \sref{IV.9.7.7} relativas a las propiedades (i) y (ii) se generalizan como sigue.

\begin{proposition}[9.7.8]
\label{IV.9.7.8}
Sea $S$ un preesquema irreducible de punto genérico $\eta$, y sea $f:X\to S$ un morfismo de presentación finita. Sea $n$ (resp.\ $n'$) el número geométrico de componentes irreducibles (resp.\ componentes conexas) de $X_\eta$ \sref{IV.4.5.2}. Entonces existe un entorno $U$ de $\eta$ en $S$ tal que, para todo $s\in U$, el número geométrico de componentes irreducibles (resp.\ componentes conexas) de $X_s$ es igual a $n$ (resp.\ $n'$).
\end{proposition}

\begin{proof}
Teniendo en cuenta que el número geométrico de componentes irreducibles (resp.\ componentes conexas) de un preesquema algebraico es invariante por extensión del cuerpo base \sref{IV.4.5.6}, vemos, por el método de \sref{IV.9.2.3}, que podemos reducirnos al caso en que $S$ es afín, noetheriano e íntegro. Además, razonando como al principio de la demostración de \sref{IV.9.7.7}, vemos que podemos suponer que las componentes irreducibles (resp.\ componentes conexas) de $X_\eta$ son geométricamente irreducibles (resp.\ geométricamente conexas). Ya sabemos por \sref{IV.9.7.1} que, para todo $s$ próximo a $\eta$, $X_s$ tiene al menos $n$ componentes irreducibles y $n'$ componentes conexas.

Por otra parte, si $Z_i$ (resp.\ $Z'_j$) son las componentes irreducibles (resp.\ componentes conexas) de $X_\eta$, sea $X_i$ (resp.\ $X'_j$) el subpreesquema cerrado reducido de $X$ cuyo espacio subyacente es la adherencia de $Z_i$ (resp.\ $Z'_j$) en $X$. Se sigue de \sref{IV.9.5.1} que, en un entorno de $\eta$, los $(X_i)_s$ (resp.\ $(X'_j)_s$) recubren $X_s$, y que los $(X'_j)_s$ son dos a dos disjuntos. Como los $(X_i)_s$ (resp.\ $(X'_j)_s$) son geométricamente irreducibles (resp.\ geométricamente conexos) por \sref{IV.9.7.7} para $s$ próximo a $\eta$, se sigue que $X_s$ tiene a lo sumo $n$ componentes irreducibles y a lo sumo $n'$ componentes conexas, lo que demuestra la proposición.
\end{proof}

\begin{corollary}[9.7.9]
\label{IV.9.7.9}
Sea $f:X\to S$ un morfismo de presentación finita. Para todo $s\in S$, sea $n(s)$ (resp.\ $n'(s)$) el número geométrico de componentes irreducibles (resp.\ componentes conexas) de $f^{-1}(s)$ \sref{IV.4.5.2}. Entonces la función $s\mapsto n(s)$ (resp.\ $s\mapsto n'(s)$) es localmente construible sobre $S$.
\end{corollary}

\begin{proof}
Basta demostrar que la propiedad ``$X$ es un preesquema algebraico sobre un cuerpo $k$ y el número geométrico de componentes irreducibles (resp.\ componentes conexas) de $X$ es igual a $n$ (resp.\ $n'$)'' es construible. Se sigue de \sref{IV.4.5.6} que esta propiedad satisface la condición \sref{IV.9.2.1},~(i); nos reducimos, por tanto, al caso en que $S$ es afín, noetheriano e íntegro de punto genérico $\eta$, y la conclusión se sigue entonces de \sref{IV.9.7.8}.
\end{proof}

\begin{corollary}[9.7.10]
\label{IV.9.7.10}
Sea $X$ un preesquema localmente noetheriano tal que, si $X'$ es la normalización de $X_{\mathrm{red}}$, el morfismo canónico $f:X'\to X$ sea finito. Entonces el conjunto de los puntos $x\in X$ tales que $X$ es geométricamente unirrama en $x$ es localmente construible en $X$.
\end{corollary}

\begin{proof}
En efecto, este conjunto es, por definición, el conjunto de los puntos $x\in X$ tales que el número de puntos geométricos de $f^{-1}(x)$ es igual a $1$. Como $f$ es finito, este número es también el número geométrico de componentes irreducibles del espacio discreto $f^{-1}(x)$ (teniendo en cuenta la definición de la normalización \sref[II]{II.6.3.8} y \sref{IV.4.5.11}); la conclusión se sigue, por tanto, de \sref{IV.9.7.9}.
\end{proof}

\begin{remark}[9.7.11]
\label{IV.9.7.11}
Sea $Z$ un subconjunto localmente construible de $X$ tal que, para todo $s\in S$, $Z\cap X_s$ sea abierto en $X_s$, y denotemos por $Z_s$ el subpreesquema de $X_s$ inducido

\oldpage[IV-3]{83}

sobre el subconjunto abierto $Z\cap X_s$. Entonces, en \sref{IV.9.7.7}, se puede sustituir en todas partes $X_s$ por $Z_s$ sin modificar la conclusión; esto se sigue repitiendo el argumento utilizado en \sref{IV.9.6.7},~(ii).
\end{remark}

\begin{proposition}[9.7.12]
\label{IV.9.7.12}
Sea $f:X\to S$ un morfismo de presentación finita, y sea $g:S\to X$ una $S$-sección de $X$ \sref[I]{I.2.5.5}. Para todo $s\in S$, sea $X_s^0$ la componente conexa de $f^{-1}(s)$ que contiene $g(s)$. Entonces la unión $X^0$ de los $X_s^0$, para $s\in S$, es un subconjunto localmente construible de $X$.
\end{proposition}

\begin{proof}
Demostremos primero que podemos reducirnos al caso en que $S$ es afín y noetheriano. Siempre podemos suponer que $S=\operatorname{Spec}(A)$ es afín; con las notaciones de \sref{IV.9.2.3}, se tiene $f=(f_0)_{(S)}$, donde $f_0:X_0\to S_0$ es un morfismo de tipo finito, y podemos además suponer que existe una $S_0$-sección $g_0:S_0\to X_0$ tal que $g=(g_0)_{(S)}$ \sref{IV.8.9.1}. Si $p:S\to S_0$ es el morfismo canónico, entonces, para todo $s_0\in S_0$, la componente conexa $(X_0)^0_{s_0}$ de $f_0^{-1}(s_0)$ que contiene $g_0(s_0)$ es geométricamente conexa \sref{IV.4.5.13}; por consiguiente, si $s_0=p(s)$, se tiene $X_s^0=q^{-1}((X_0)^0_{s_0})$, donde $q:X_s\to (X_0)_{s_0}$ es la proyección canónica \sref{IV.4.5.8} y \sref{IV.4.4.1}. Nuestra afirmación se sigue, por tanto, de \sref{IV.1.8.2}.

Utilizamos ahora el criterio de construibilidad \sref[0\textsubscript{III}]{0.9.2.3}. Sea $x$ un punto de $X$, sea $Z$ el subpreesquema reducido de $X$ cuyo espacio subyacente es $\overline{\{x\}}$, y sea $Y$ el subpreesquema reducido de $S$ cuyo espacio subyacente es $\overline{f(Z)}$. Como la restricción de $f$ a $Z$ se factoriza como $Z\to Y\to S$ \sref[I]{I.5.2.2}, podemos sustituir $S$ por $Y$, es decir, suponer que $f(x)=\eta$, el punto genérico del preesquema íntegro $S$.

Por hipótesis, $X_\eta$ es la suma de dos subpreesquemas $X_\eta^0$ y $X_\eta^1$ inducidos sobre subconjuntos abiertos complementarios de $X_\eta$. Por \sref{IV.9.5.4} y \sref{IV.9.5.1}, tras sustituir $S$, si es necesario, por un entorno de $\eta$, podemos suponer que $X$ es la unión de dos subconjuntos abiertos disjuntos $X^0$ y $X^1$ tales que $(X^0)_\eta=X_\eta^0$ y $(X^1)_\eta=X_\eta^1$. Como $g:S\to X$ es continua e inyectiva, $S$ es la unión de los dos subconjuntos abiertos disjuntos $g^{-1}(X^0)$ y $g^{-1}(X^1)$; pero $S$ es irreducible y, \emph{a fortiori}, conexo, por lo que uno de estos dos subconjuntos abiertos es vacío. Como $g(\eta)\in X_\eta^0$ por definición, se tiene $g^{-1}(X^1)=\varnothing$.

En otras palabras, $g$ es una $S$-sección de $X^0$. Por otra parte, como $(X^0)_\eta$ es geométricamente conexo \sref{IV.4.5.13}, lo mismo vale para $(X^0)_s$ para todo $s$ próximo a $\eta$ por \sref{IV.9.7.7}; y como $g(s)\in (X^0)_s$, se tiene efectivamente $(X^0)_s=X_s^0$.
\end{proof}

\subsection{Descomposición primaria en un entorno de una fibra genérica}
\label{subsection:IV.9.8}
\label{IV.9.8}

\begin{proposition}[9.8.1]
\label{IV.9.8.1}
Sea $f:X\to S$ un morfismo de presentación finita y sea $\sh{F}$ un $\sh{O}_X$-Módulo cuasicoherente de presentación finita.
Entonces el conjunto $E$ de los $s\in S$ tales que $\sh{F}_s$ es un $\sh{O}_{X_s}$-Módulo sin ciclo primo inmerso asociado es localmente construible.
\end{proposition}

\begin{proof}
\label{IV.9.8.1-proof}
Sabemos que, para los Módulos cuasicoherentes sobre preesquemas algebraicos, la propiedad de carecer de ciclo primo inmerso asociado es invariante por cambio del cuerpo base \sref{IV.4.2.7}; podemos, por tanto, restringirnos al caso en que $S$ es afín, noetheriano e íntegro, y a demostrar que en este caso $E$ o $S-E$ es un entorno de $\eta$.
Vimos en \sref{IV.9.7.6} que si $\eta\in E$, $E$ es un entorno de $\eta$; queda, por tanto, considerar el caso en que $\sh{F}_\eta$ tiene ciclos primos inmersos asociados.
Podemos restringirnos al caso en que $X$ es afín, y en que existe un sub-$\sh{O}_{X_\eta}$-Módulo coherente $\sh{H}$ de $\sh{F}_\eta$ cuyo soporte $Z'$ es no vacío y \emph{raro} respecto del soporte $Z$ de $\sh{F}$ \sref{IV.3.1.3}; entonces, en virtud de \sref{IV.8.5.2},~(i) y~(ii),
\oldpage[IV-3]{84}
y \sref{IV.8.5.8}, aplicado según el método de \sref{IV.8.1.2},~$a$), existe un sub-$\sh{O}_X$-Módulo coherente $\sh{G}$ de $\sh{F}$ tal que $\sh{G}_\eta=\sh{H}$; si $Y=\operatorname{Supp}(\sh{F})$ e $Y'=\operatorname{Supp}(\sh{G})$, entonces, por consiguiente, $Y_s=\operatorname{Supp}(\sh{F}_s)$ e $Y'_s=\operatorname{Supp}(\sh{G}_s)$ para todo $s\in S$ \sref[I]{I.9.1.13}, y en particular $Z=Y_\eta$ y $Z'=Y'_\eta$; como $Z-Z'$ es denso en $Z$ y $Z'\neq\varnothing$, existe un entorno $U$ de $\eta$ tal que para $s\in U$, $Y_s-Y'_s$ es denso en $Y_s$ e $Y'_s\neq\varnothing$ \sref{IV.9.5.1} y \sref{IV.9.5.3}; considerando un punto genérico de una componente irreducible de $Y'_s$ y un entorno suficientemente pequeño de este punto en $X_s$, deducimos inmediatamente de \sref{IV.3.1.3} que $\sh{F}_s$ admite ciclos primos inmersos asociados.
\end{proof}

\medskip
\noindent
\emph{\sref{IV.9.8.2}}
\label{IV.9.8.2}
Sea $S$ un preesquema noetheriano, \emph{íntegro}, de punto genérico $\eta$, sea $f:X\to S$ un morfismo de tipo finito y sea $\sh{F}$ un $\sh{O}_X$-Módulo coherente.
Consideremos una descomposición reducida irredundante $(\sh{G}_i)_{i\in I}$ de $\sh{F}_\eta$ \sref{IV.3.2.6}; los $\sh{G}_i$ son, pues, cocientes de $\sh{F}_\eta$ y existe un homomorfismo inyectivo $h:\sh{F}_\eta\to\bigoplus_i\sh{G}_i$; además, $\operatorname{Ass}(\sh{G}_i)$ se reduce a un punto $x_i$.
Sea $Z_i$ la adherencia de $\{x_i\}$ en $X$, de modo que $(Z_i)_\eta=\operatorname{Supp}(\sh{G}_i)$.
Denotemos por $\sh{J}_i$ el Ideal coherente de $\sh{O}_X$ que define el subpreesquema cerrado de $X$ cuyo espacio subyacente es $Z_i$ (subpreesquema que seguimos denotando por $Z_i$).
En virtud de \sref{IV.8.5.2} y \sref{IV.8.5.8}, aplicados según el método de \sref{IV.8.1.2},~$a$), podemos (restringiendo $S$ a un entorno de $\eta$ si es necesario) suponer que existen cocientes $\sh{F}_i$ de $\sh{F}$ ($i\in I$) tales que $\sh{G}_i=(\sh{F}_i)_\eta$ para todo $i$, y un homomorfismo $g:\sh{F}\to\bigoplus_i\sh{F}_i$ tal que $h=g_\eta$.
Además, por \sref[I]{I.9.3.5}, existe un entero $m$ tal que $(\sh{J}_i^m\sh{F}_i)_\eta=(\sh{J}_i)_\eta^m\sh{G}_i=0$ y, restringiendo de nuevo $S$, podemos suponer que $\sh{J}_i^m\sh{F}_i=0$ \sref{IV.8.5.2.5}, de modo que el soporte de $\sh{F}_i$ está contenido en $Z_i$; pero como es cerrado y contiene $x_i$, es necesariamente igual a $Z_i$.

\begin{proposition}[9.8.3]
\label{IV.9.8.3}
Bajo las condiciones de \emph{\sref{IV.9.8.2}}, para todo $s\in S$ y todo $i\in I$, denotemos por $x_{is\alpha}$ ($\alpha\in J_{s,i}$) los puntos maximales de $(Z_i)_s$.
Existe un entorno $U$ de $\eta$ en $S$ tal que, para todo $s\in U$, los $x_{is\alpha}$ (para $i\in I$ y, para cada $i$, $\alpha\in J_{s,i}$) son dos a dos distintos y $\operatorname{Ass}(\sh{F}_s)$ es el conjunto de los $x_{is\alpha}$ (en otras palabras, los ciclos primos asociados a $\sh{F}_s$ son las componentes irreducibles de los $(Z_i)_s$).
Además, se puede tomar $U$ de modo que, para que la adherencia de $\{x_{is\alpha}\}$ en $X_s$ sea un ciclo primo asociado maximal de $\sh{F}_s$, sea necesario y suficiente que $(Z_i)_\eta$ (adherencia de $\{x_i\}$ en $X_\eta$) sea un ciclo primo asociado maximal de $\sh{F}_\eta$.
\end{proposition}

\begin{proof}
\label{IV.9.8.3-proof}
Se sigue de \sref{IV.3.1.3},~$c'$) que para cada $i$, existe un abierto $W_i$ en $X_\eta$ tal que $W_i\cap(Z_i)_\eta$ no es vacío, y un $\sh{O}_{X_\eta}$-Módulo coherente $\sh{H}_i$, de soporte $W_i\cap(Z_i)_\eta$, tal que existe un homomorfismo inyectivo $v_i:\sh{H}_i\to\sh{F}_\eta|W_i$.
Sea $V_i$ un abierto de $X$ tal que $V_i\cap X_\eta=W_i$; aplicando, como en \sref{IV.9.8.2}, los resultados de \sref{IV.8.5.2} y \sref{IV.8.5.8}, podemos (restringiendo $S$) suponer que existe un Módulo coherente $\sh{F}'_i$, de soporte $V_i\cap Z_i$, y un homomorfismo $u_i:\sh{F}'_i\to\sh{F}|V_i$ tales que $\sh{H}_i=(\sh{F}'_i)_\eta$ y $v_i=(u_i)_\eta$.
Mostraremos que existe un entorno $U$ de $\eta$ tal que para $s\in U$, se satisfacen las propiedades siguientes:
\begin{enumerate}[label=(\roman*)]
  \item Los $(\sh{F}_i)_s$ carecen de ciclo primo inmerso asociado.
  \item El homomorfismo $g_s:\sh{F}_s\to\bigoplus_i(\sh{F}_i)_s$ es inyectivo.
  \item $(Z_i)_s\cap(V_i)_s$ es denso en $(Z_i)_s$, $(\sh{F}'_i)_s$ tiene soporte $(Z_i)_s\cap(V_i)_s$, y $(u_i)_s:(\sh{F}'_i)_s\to\sh{F}_s|(V_i)_s$ es inyectivo.
\end{enumerate}
\oldpage[IV-3]{85}

\begin{enumerate}[label=(\roman*)]
  \item[(iv)] Para $i\neq j$, toda componente irreducible de $(Z_i)_s$ es distinta de toda componente irreducible de $(Z_j)_s$.
\end{enumerate}

Ahora bien, (i) ya se ha visto en \sref{IV.9.7.6}; (ii) es un caso particular de \sref{IV.9.4.5}; (iii) se sigue de \sref{IV.9.5.3} y \sref{IV.9.4.5}.
Finalmente, si $i\neq j$, $(Z_i)_\eta\cap(Z_j)_\eta=(Z_i\cap Z_j)_\eta$ es raro en $(Z_i)_\eta$ \emph{o} en $(Z_j)_\eta$; supongamos, por ejemplo, que $(Z_i\cap Z_j)_\eta$ es raro en $(Z_j)_\eta$; se sigue entonces de \sref{IV.9.5.3} y \sref{IV.9.5.4} que para $s$ en un entorno de $\eta$, $(Z_i\cap Z_j)_s=(Z_i)_s\cap(Z_j)_s$ es raro en $(Z_j)_s$, lo que muestra que ninguna componente irreducible de $(Z_j)_s$ puede estar contenida en una componente irreducible de $(Z_i)_s$, ni \emph{a fortiori} ser igual a ella.

En estas condiciones, se sigue de (ii) y \sref{IV.3.1.7} que para $s\in U$, tenemos
\[
  \operatorname{Ass}(\sh{F}_s)\subset\bigcup_i\operatorname{Ass}((\sh{F}_i)_s),
\]
de (i), que $\operatorname{Ass}((\sh{F}_i)_s)$ es el conjunto de los $x_{is\alpha}$ ($\alpha\in J_{s,i}$).
Por otra parte, en virtud de (iii) y del criterio \sref{IV.3.1.3},~$c'$), cada uno de los $x_{is\alpha}$ ($\alpha\in J_{s,i}$, $i\in I$) pertenece a $\operatorname{Ass}(\sh{F}_s)$.
Finalmente, (iv) significa que los $x_{is\alpha}$ son dos a dos distintos.

Queda demostrar la última afirmación del enunciado.
Ahora bien, se sigue de (i) que para $s$ y para $i$ dado, ninguno de los conjuntos $\{x_{is\alpha}\}$ puede estar contenido en otro para $\alpha\in J_{s,i}$.
Por otra parte, si $(Z_j)_\eta\subset(Z_i)_\eta$, podemos suponer $U$ tomado de modo que $(Z_j)_s\subset(Z_i)_s$ para $s\in U$ \sref{IV.9.5.1}, por lo que cada $x_{js\beta}$ pertenece a la adherencia de un $x_{is\alpha}$; por el contrario, si $(Z_j)_\eta\cap(Z_i)_\eta$ es raro en $(Z_j)_\eta$, hemos visto al demostrar (iv) que $(Z_j)_s\cap(Z_i)_s$ es raro en $(Z_j)_s$, por lo que ninguno de los $x_{js\beta}$ es adherente a un $x_{is\alpha}$.
En particular, si $(Z_j)_\eta$ es maximal, esto equivale a decir que $(Z_j)_\eta\cap(Z_i)_\eta$ es raro en $(Z_j)_\eta$ para \emph{todo} $i\neq j$, y concluimos que $(Z_j)_s\cap(Z_i)_s$ es raro en $(Z_j)_s$ para todo $i\neq j$, de modo que todo $x_{js\alpha}$ ($\alpha\in J_{s,j}$) es \emph{maximal}.
\end{proof}

\begin{corollary}[9.8.4]
\label{IV.9.8.4}
Siendo las notaciones e hipótesis las de \emph{\sref{IV.9.8.2}}, existe un entorno $U$ de $\eta$ en $S$ tal que, para todo $s\in U$, cada uno de los $(\sh{F}_i)_s$ carece de ciclo primo inmerso asociado; además, si $(\sh{F}_{is\alpha})_{\alpha\in J_{s,i}}$ es la única descomposición reducida irredundante de $(\sh{F}_i)_s$, entonces, para todo $s\in U$, la familia $(\sh{F}_{is\alpha})_{i\in I,\,\alpha\in J_{s,i}}$ es una descomposición reducida irredundante de $\sh{F}_s$.
\end{corollary}

\begin{proof}
\label{IV.9.8.4-proof}
La primera afirmación se sigue de \sref{IV.9.8.1} y de la definición de $\sh{F}_i$.
Por otra parte, hemos visto en \sref{IV.9.8.3} que el homomorfismo $\sh{F}_s\to\bigoplus_i(\sh{F}_i)_s$ es inyectivo y, por definición, lo mismo vale para cada uno de los homomorfismos $(\sh{F}_i)_s\to\bigoplus_\alpha\sh{F}_{is\alpha}$, por lo que el homomorfismo $\sh{F}_s\to\bigoplus_{i,\alpha}\sh{F}_{is\alpha}$ es inyectivo.
Como podemos suponer por \sref{IV.9.4.5} que para $s\in U$ cada uno de los $(\sh{F}_i)_s$ es un cociente de $\sh{F}_s$, los $\sh{F}_{is\alpha}$ son cocientes de $\sh{F}_s$, y queda verificar por \sref{IV.3.2.5} que los $x_{is\alpha}$ son dos a dos distintos y pertenecen a $\operatorname{Ass}(\sh{F}_s)$, lo que se ha demostrado en \sref{IV.9.8.3}.
\end{proof}

\begin{proposition}[9.8.5]
\label{IV.9.8.5}
Sea $f:X\to S$ un morfismo de presentación finita y sea $\sh{F}$ un $\sh{O}_X$-Módulo cuasicoherente de presentación finita.
Para todo $s\in S$, sea $E(s)$ (resp.~$E'(s)$) el conjunto (subconjunto finito de $\mathbf{Z}\cup\{-\infty\}$) de las dimensiones de los ciclos primos asociados a $\sh{F}_s$ (resp.~de los ciclos primos asociados maximales, es decir, las componentes irreducibles de $\operatorname{Supp}(\sh{F}_s)$).
Entonces las funciones $s\mapsto E(s)$ y $s\mapsto E'(s)$ son localmente construibles en $S$.
\end{proposition}

\begin{proof}
\label{IV.9.8.5-proof}
\oldpage[IV-3]{86}
Se sigue de \sref{IV.4.2.7} y \sref{IV.4.2.8} que la condición \sref{IV.9.2.1},~(i) se satisface para las propiedades cuya construibilidad queremos demostrar.
Nos reducimos así al caso en que $S$ es afín, noetheriano e íntegro, y a demostrar que $E$ y $E'$ son constantes en un entorno del punto genérico $\eta$ de $S$; pero esto se sigue de \sref{IV.9.8.3} y \sref{IV.9.5.5}.
\end{proof}

\begin{proposition}[9.8.6]
\label{IV.9.8.6}
Con las hipótesis y notaciones de \emph{\sref{IV.9.8.2}}, sea $i\in I$ tal que $(Z_i)_\eta$ sea maximal (en otras palabras, una componente irreducible de $\operatorname{Supp}(\sh{F}_\eta)$).
Entonces existe un entorno $U$ de $\eta$ en $S$ tal que para todo $s\in U$ y todo $\alpha\in J_{s,i}$, la longitud geométrica de $\sh{F}_s$ en $x_{is\alpha}$ (relativa a $\mathbf{k}(s)$) \sref{IV.4.7.5} es igual a la longitud geométrica de $\sh{F}_\eta$ en $x_i$ (relativa a $\mathbf{k}(\eta)$).
\end{proposition}

\begin{proof}
\label{IV.9.8.6-proof}
Podemos evidentemente restringirnos al caso en que $S=\operatorname{Spec}(A)$ es afín; demostremos primero que podemos reducirnos al caso en que el subpreesquema $(Z_i)_\eta$ de $X_\eta$, que es reducido, es \emph{geométricamente íntegro}.
En efecto, existe una extensión finita $K'$ de $K=\mathbf{k}(\eta)$ tal que $(((Z_i)_\eta)_{(K')})_{\mathrm{red}}$ es geométricamente reducido y que las componentes irreducibles de $((Z_i)_\eta)_{(K')}$ son geométricamente irreducibles \sref{IV.4.6.8}.
Procedemos como en la demostración de \sref{IV.9.7.7}, considerando una sub-$A$-álgebra $A'$ de $K'$ que tiene $K'$ como cuerpo de fracciones y es finita sobre $A$.
Ponemos $S'=\operatorname{Spec}(A')$ y consideramos el morfismo finito sobreyectivo $g:S'\to S$; sean $X'=X_{(S')}$ y $\sh{F}'=\sh{F}\otimes_{\sh{O}_S}\sh{O}_{S'}$, y sea $\eta'$ el punto genérico de $S'$.
Para todo $s'\in S'$, pongamos $s=g(s')$; si $T$ es una componente irreducible de $\operatorname{Supp}(\sh{F}_s)$, las componentes irreducibles $T'_j$ de $T_{(\mathbf{k}(s'))}$ son componentes irreducibles de $\operatorname{Supp}(\sh{F}'_{s'})$ y dominan $T$ \sref{IV.4.2.7}, y las multiplicidades radiciales de $T$ respecto de $\sh{F}_s$ y de cada $T'_j$ respecto de $\sh{F}'_{s'}$ son las mismas \sref{IV.4.7.9}.
El razonamiento de la primera parte de \sref{IV.9.7.7} muestra entonces que podemos restringirnos a demostrar la proposición para $X'$ y $\sh{F}'$; y en virtud de la elección de $K'$, los subpreesquemas reducidos cuyos espacios subyacentes son las componentes irreducibles de $((Z_i)_\eta)_{(K')}$ son geométricamente íntegros \sref{IV.4.6.1}.

Supongamos ahora que $(Z_i)_\eta$ es geométricamente íntegro; entonces, por \sref{IV.9.7.7}, lo mismo vale para $(Z_i)_s$ para $s$ en un entorno de $\eta$; la definición \sref{IV.4.7.5} muestra que bastará, por tanto, demostrar que la longitud del $\sh{O}_{X_\eta,x_i}$-módulo $(\sh{F}_\eta)_{x_i}$ es igual a la del $\sh{O}_{X_s,x_{is}}$-módulo $(\sh{F}_s)_{x_{is}}$ (donde hemos suprimido el índice $\alpha$, ahora innecesario por hipótesis).
Como la cuestión es evidentemente local sobre $X$, podemos suponer (restringiendo a un entorno de $x_i$) que $\sh{F}=\sh{F}_i$, de modo que $\sh{F}_\eta$ es \emph{irredundante} sobre $X=\operatorname{Spec}(B)$ afín, y escribimos $Z$ en lugar de $Z_i$, y $x$ para el punto genérico de $Z$ (y de $Z_\eta$).
El anillo noetheriano $B$ contiene $A$ como subanillo, y $\sh{F}=\widetilde{M}$, donde $M$ es un $B$-módulo de tipo finito; además, si $K$ es el cuerpo de fracciones de $A$, el $B_{(K)}$-módulo $M_{(K)}$ es $\neq 0$ e irredundante.
Sea $q$ el único elemento de $\operatorname{Ass}(M_{(K)})$ y sea $\mathfrak{p}$ el ideal primo de $B$, imagen inversa de $q$.
Utilizando \sref{IV.5.11.1.1} como en la demostración de \sref{IV.9.7.6}, nos reducimos al caso en que $B$ es \emph{íntegro} y $M$ un submódulo $\neq 0$ de $B$; entonces $\sh{F}_\eta$ es un sub-$\sh{O}_{X_\eta}$-Módulo no nulo de $\sh{O}_{Z_\eta}$, y en virtud de \sref{IV.9.4.5}, para $s$ en un entorno de $\eta$, $\sh{F}_s$ es isomorfo a un sub-$\sh{O}_{X_s}$-Módulo no nulo de $\sh{O}_{Z_s}$; como $Z_s$ es geométricamente íntegro, las longitudes de $(\sh{F}_\eta)_x$ y $(\sh{F}_s)_{x_s}$ son ambas iguales a $1$, lo que termina la demostración.
\end{proof}

\begin{corollary}[9.8.7]
\label{IV.9.8.7}
\oldpage[IV-3]{87}
Sea $f:X\to S$ un morfismo de presentación finita y sea $\sh{F}$ un $\sh{O}_X$-Módulo cuasicoherente de presentación finita.
Para todo $s\in S$, sea $M(s)$ el conjunto de los pares $(d,m)$ tales que existe una componente irreducible de $\operatorname{Supp}(\sh{F}_s)$ de dimensión $d$ y de multiplicidad radicial $m$ para $\sh{F}_s$ \sref{IV.4.7.8}.
Entonces la función $s\mapsto M(s)$ es localmente construible en $S$.
\end{corollary}

\begin{proof}
\label{IV.9.8.7-proof}
Se sigue de \sref{IV.4.2.7} y \sref{IV.4.7.9} que la condición \sref{IV.9.2.1},~(i) se satisface para la propiedad cuya construibilidad queremos demostrar.
Nos reducimos así al caso en que $S$ es afín, noetheriano e íntegro, y a demostrar que $M$ es constante en un entorno del punto genérico de $S$; pero esto se sigue de la proposición, a saber, \sref{IV.9.8.3} y \sref{IV.9.8.6}.
\end{proof}

\begin{proposition}[9.8.8]
\label{IV.9.8.8}
Sea $f:X\to S$ un morfismo de presentación finita y sea $\sh{F}$ un $\sh{O}_X$-Módulo cuasicoherente de presentación finita.
Para todo $s\in S$, sea $l(s)$ la suma de las multiplicidades totales para $\sh{F}_s$ (relativas a $\mathbf{k}(s)$) de los puntos genéricos de las componentes irreducibles de $\operatorname{Supp}(\sh{F}_s)$ \sref{IV.4.7.12}.
Entonces la función $s\mapsto l(s)$ es localmente construible en $S$.
\end{proposition}

\begin{proof}
\label{IV.9.8.8-proof}
Teniendo en cuenta \sref{IV.4.7.12}, la condición \sref{IV.9.2.1},~(i) se satisface para la propiedad cuya construibilidad queremos demostrar, y nos reducimos así al caso en que $S$ es afín, noetheriano e íntegro, y a demostrar que $l$ es constante en un entorno del punto genérico $\eta$.
Utilizando las notaciones de \sref{IV.9.8.2}, esto se sigue de la definición \sref{IV.4.7.12}, del hecho de que el número geométrico de componentes irreducibles de cada $(Z_i)_s$ es constante en un entorno de $\eta$ \sref{IV.9.7.8}, de que la longitud geométrica de $\sh{F}_s$ en $x_{is\alpha}$ es también igual a la de $\sh{F}_\eta$ en $x_i$ para cada $i$ tal que $(Z_i)_\eta$ es maximal \sref{IV.9.8.6}, y finalmente del hecho de que la adherencia de $\{x_{is\alpha}\}$ es un ciclo primo asociado maximal de $\sh{F}_s$ si y solo si $(Z_i)_\eta$ es un ciclo primo asociado maximal de $\sh{F}_\eta$ \sref{IV.9.8.3}.
\end{proof}

\begin{remark}[9.8.9]
\label{IV.9.8.9}
Se pueden precisar de diversas maneras las proposiciones precedentes; limitémonos a un enunciado dado a modo de ejemplo.
Diremos que un subconjunto finito $P$ de un $k$-preesquema algebraico $X$ es \emph{saturado} si, para todo par de puntos $x$, $y$ de $P$, los puntos genéricos de las componentes irreducibles de $\overline{\{x\}}\cap\overline{\{y\}}$ pertenecen también a $P$; para todo subconjunto finito $Q$ de $X$, existe un menor subconjunto finito $P$ de $X$ que contiene $Q$ y es saturado; decimos que $P$ es el \emph{saturado} de $Q$.
Para todo $\sh{O}_X$-Módulo coherente $\sh{F}$, llamaremos \emph{esqueleto primario} de $\sh{F}$ al sistema $(P,Q,\omega,d,m)$ donde $Q=\operatorname{Ass}(\sh{F})$, $P$ es el saturado de $Q$, $\omega$ la relación de orden $\overline{\{x\}}\subset\overline{\{y\}}$ sobre $P$, $d$ la función $x\mapsto\dim\overline{\{x\}}$ sobre $P$, y $m$ la función $x\mapsto\operatorname{length}_{\sh{O}_x}\sh{F}_x$ definida sobre el conjunto de los elementos maximales de $Q$ para la relación $\omega$.
Llamamos \emph{esqueleto virtual} a todo sistema $(P,Q,\omega,d,m)$ donde $P$ es un conjunto, $Q$ un subconjunto de $P$, $\omega$ una relación de orden sobre $P$, $d$ una aplicación de $P$ en $\mathbf{N}$, y $m$ una aplicación del conjunto de los elementos maximales de $Q$ en $\mathbf{N}$; se define de manera evidente la noción de \emph{isomorfismo} de dos esqueletos virtuales.
Finalmente, con las notaciones precedentes, llamamos \emph{tipo primario} de $\sh{F}$ a la clase (para la relación de isomorfismo de esqueletos virtuales) del esqueleto primario de $\sh{F}$.
Se sigue de \sref{IV.4.2.6}, \sref{IV.4.2.7}, \sref{IV.4.2.8}, \sref{IV.4.5.1} y \sref{IV.4.7.9} que si, para una extensión algebraicamente cerrada $k'$ de $k$, se ponen $X'=X\otimes_k k'$ y $\sh{F}'=\sh{F}\otimes_k k'$, el \emph{tipo primario} de $\sh{F}'$ es independiente de la
\oldpage[IV-3]{88}
extensión algebraicamente cerrada $k'$ de $k$ considerada; decimos que este es el \emph{tipo primario geométrico} de $\sh{F}$.
Dadas las definiciones, el enunciado que tenemos en vista es el siguiente:

\medskip
\noindent
\emph{\sref{IV.9.8.9.1}}
\label{IV.9.8.9.1}
Sea $f:X\to S$ un morfismo de presentación finita y sea $\sh{F}$ un $\sh{O}_X$-Módulo cuasicoherente de presentación finita.
Para todo $s\in S$, sea $T(s)$ el tipo primario geométrico de $\sh{F}_s$.
Entonces la función $s\mapsto T(s)$ es localmente construible en $S$.

\medskip
Teniendo en cuenta las observaciones precedentes, nos reducimos como de costumbre a demostrar que si $S$ es afín, noetheriano e íntegro, de punto genérico $\eta$, entonces $T(s)$ es constante en un entorno de $\eta$.
Razonando como al principio de \sref{IV.9.7.7}, podemos suponer que todas las componentes irreducibles de $X_\eta$ que intervienen son geométricamente irreducibles, y entonces la proposición se sigue de \sref{IV.9.5.1}, \sref{IV.9.5.5}, \sref{IV.9.8.3} y \sref{IV.9.8.6}; dejamos los detalles al lector.
Se podría generalizar considerando varios Módulos coherentes y definiendo su ``esqueleto primario simultáneo'', etc.
La conclusión general de este parágrafo es que para todas las propiedades del tipo considerado (y para $S$ irreducible), \emph{las propiedades válidas sobre la ``fibra genérica'' son también válidas sobre todas las fibras vecinas}.
\end{remark}

\subsection{Construibilidad de propiedades locales de las fibras}
\label{subsection:IV.9.9}
\label{IV.9.9}

\begin{proposition}[9.9.1]
\label{IV.9.9.1}
Sea $f:X\to S$ un morfismo localmente de presentación finita, sea $Z$ un subconjunto localmente construible de $X$ tal que para todo $s\in S$, $Z_s$ es cerrado en $X_s$, y sea $\Phi$ un subconjunto finito de $\mathbf{Z}\cup\{\pm\infty\}$.
Entonces los subconjuntos siguientes de $X$ son localmente construibles:
\begin{enumerate}[label=(\roman*)]
  \item El conjunto de los $x\in X$ tales que $\dim_x(Z_{f(x)})$ pertenece a $\Phi$.
  \item El conjunto de los $x\in X$ tales que $\operatorname{codim}_x(Z_{f(x)},X_{f(x)})$ pertenece a $\Phi$.
  \item El conjunto de los $x\in X$ tales que el anillo local $\sh{O}_{X_{f(x)},x}$ es equidimensional.
\end{enumerate}

Observamos que las propiedades (i) y (ii) pueden expresarse también diciendo que las funciones $x\mapsto\dim_x(Z_{f(x)})$ y $x\mapsto\operatorname{codim}_x(Z_{f(x)},X_{f(x)})$ son \emph{localmente construibles} en $X$ ($0_{\mathrm{III}}$, \sref{IV.9.3.1}).
\end{proposition}

\begin{proof}
\label{IV.9.9.1-proof}
Como las cuestiones son locales sobre $X$, podemos restringirnos al caso en que $S=\operatorname{Spec}(A)$ y $X=\operatorname{Spec}(B)$ son afines y en que $f$ es un morfismo de presentación finita; existe entonces un subanillo $A_0$ de $A$ que es una $\mathbf{Z}$-álgebra de tipo finito, un $A_0$-preesquema de tipo finito $X_0$ y un subconjunto construible $Z_0$ de $X_0$ tales que $X=X_0\otimes_{A_0}A$ y $Z=h^{-1}(Z_0)$, donde $h:X\to X_0$ es la proyección canónica (\sref{IV.8.9.1} y \sref{IV.8.3.11}).
Además, para todo $s\in S$, si $s_0$ es la proyección de $s$ en $S_0=\operatorname{Spec}(A_0)$, tenemos
\[
  X_s=(X_0)_{s_0}\otimes_{\mathbf{k}(s_0)}\mathbf{k}(s),
\]
y, si $h_s$ es la proyección $X_s\to(X_0)_{s_0}$, tenemos $Z_s=h_s^{-1}((Z_0)_{s_0})$.
Como el morfismo $h_s$ es fielmente plano y cuasicompacto, la hipótesis de que $Z_s$ es cerrado en $X_s$ implica que $(Z_0)_{s_0}$ es cerrado en $(X_0)_{s_0}$ \sref{IV.2.3.12}.

En estas condiciones, la transitividad de las fibras \sref[I]{I.3.6.4} y la prop.~\sref{IV.4.2.7} implican que el conjunto de las dimensiones de las componentes irreducibles de $Z_s$ que contienen $x$ es el mismo que el conjunto de las dimensiones de las componentes irreducibles de $(Z_0)_{s_0}$ que contienen $x_0=h_s(x)$.
En particular, tenemos $\dim_x(Z_s)=\dim_{x_0}((Z_0)_{s_0})$.
Por otra parte, si $Z_s^{(\beta)}$ son las componentes irreducibles de $Z_s$ que contienen $x$ y $X_s^{(\alpha)}$ las componentes irreducibles de $X_s$
\oldpage[IV-3]{89}
que contienen $x$, tenemos
\[
  \operatorname{codim}_x(Z_s,X_s)
  =\inf_{\beta}\sup_{\alpha}\operatorname{codim}(Z_s^{(\beta)},X_s^{(\alpha)}),
\]
donde $(\alpha,\beta)$ recorre los pares tales que $x\in Z_s^{(\beta)}\subset X_s^{(\alpha)}$ ($0$, \sref{IV.14.2.6}).
Como los preesquemas algebraicos irreducibles son biequidimensionales \sref{IV.5.2.1}, podemos escribir, en virtud de ($0$, \sref{IV.14.3.3.1})
\begin{equation}
\tag{9.9.1.1}
\label{IV.9.9.1.1}
  \operatorname{codim}_x(Z_s,X_s)=\inf_{\beta}\sup_{\alpha}\bigl(\dim(X_s^{(\alpha)})-\dim(Z_s^{(\beta)})\bigr)
\end{equation}
con la misma elección de pares $(\alpha,\beta)$.
Como $h_s$ es fielmente plano y cuasicompacto, para todo par formado por una componente irreducible $(X_0)_{s_0}^{(\lambda)}$ de $(X_0)_{s_0}$ que contiene $x_0$ y una componente irreducible $(Z_0)_{s_0}^{(\mu)}$ de $(Z_0)_{s_0}$ que contiene $x_0$ y está contenida en $(X_0)_{s_0}^{(\lambda)}$, existe un par $(X_s^{(\alpha)},Z_s^{(\beta)})$ del tipo descrito anteriormente tal que $Z_s^{(\beta)}$ domina $(Z_0)_{s_0}^{(\mu)}$ y $X_s^{(\alpha)}$ domina $(X_0)_{s_0}^{(\lambda)}$ \sref{IV.2.3.5}.
La fórmula \sref{IV.9.9.1.1} (y la fórmula análoga aplicada a $(X_0)_{s_0}$) muestra entonces, en virtud de \sref{IV.4.2.7}, que tenemos
\[
  \operatorname{codim}_x(Z_s,X_s)=\operatorname{codim}_{x_0}((Z_0)_{s_0},(X_0)_{s_0}).
\]

Vemos así que si $E$ (resp.~$E_0$) es el conjunto de los $x\in X$ (resp.~$x_0\in X_0$) que satisfacen una de las condiciones (i), (ii), (iii) del enunciado (resp.~la misma condición), entonces $E=h^{-1}(E_0)$; por \sref{IV.1.8.2}, podemos, por tanto, restringirnos al caso en que $A$, y por consiguiente también $B$, es noetheriano.
Teniendo en cuenta ($0_{\mathrm{III}}$, \sref{IV.9.2.3}) y \sref{IV.9.9.1.1}, queda demostrar que para todo $x\in X$ existe un entorno $V$ de $x$ en $\overline{\{x\}}$ tal que, para todo $x'\in V$, los conjuntos de dimensiones de las componentes irreducibles de $X_{f(x')}$ y de $Z_{f(x')}$ que contienen $x'$ son constantes, y asimismo el conjunto de pares
\[
  \bigl(\dim Z_{f(x')}^{(\beta)},\dim X_{f(x')}^{(\alpha)}\bigr),
\]
donde $X_{f(x')}^{(\alpha)}$ es una componente irreducible de $X_{f(x')}$ y $Z_{f(x')}^{(\beta)}$ es una componente irreducible de $Z_{f(x')}$ contenida en ella y que contiene $x'$.
Para ello podemos sustituir $S$ por el subpreesquema cerrado reducido $S'$ cuyo espacio subyacente es $\overline{\{f(x)\}}$, y $X$ por $X'=f^{-1}(S')$, puesto que las fibras de $X$ y $X'$ en los puntos de $S'$ son las mismas.
Así, podemos suponer que $S$ es íntegro y que $\eta=f(x)$ es su punto genérico.

Por hipótesis, $Z_\eta$ es cerrado en $X_\eta$; como $Z$ es construible, se sigue de \sref{IV.8.3.11}, aplicado según el método de \sref{IV.8.1.2},~(a), que tras sustituir $S$ por un entorno abierto de $\eta$ podemos suponer que $Z$ es \emph{cerrado} en $X$.
Sean $X_i$ (resp.~$Z_j$) las componentes irreducibles de $X$ (resp.~$Z$) que contienen $x$.
Por ($0_I$, \sref{IV.2.1.8}), los $X_i\cap X_\eta$ (resp.~$Z_j\cap X_\eta$) son las componentes irreducibles de $X_\eta$ (resp.~$Z_\eta$) que contienen $x$; por \sref{IV.9.5.1}, tras restringir $S$ alrededor de $\eta$, podemos suponer además que los $X_i$ (resp.~$Z_j$) son exactamente las componentes irreducibles de $X$ (resp.~$Z$) que encuentran $\overline{\{x\}}$, y que $(X_i)_s\cap\overline{\{x\}}$ y $(Z_j)_s\cap\overline{\{x\}}$ son no vacíos para todo $s\in S$.
En estas condiciones, se sigue aún de \sref{IV.9.7.1} que podemos suponer, restringiendo $S$, que los pares $(i,j)$ tales que $(Z_j)_s\subset(X_i)_s$ son los mismos para todo $s\in S$.
La conclusión se sigue entonces de \sref{IV.9.5.6}: para todo $s$ suficientemente próximo a $\eta$, todas las componentes irreducibles de $(X_i)_s$ (resp.~$(Z_j)_s$) tienen la misma dimensión, igual a la de $(X_i)_\eta$ (resp.~$(Z_j)_\eta$).
Además, si $(i,j)$ es un par tal que $Z_j\not\subset X_i$, $(X_i)_\eta$ no contiene el
\oldpage[IV-3]{90}
punto genérico de $(Z_j)_\eta$, de donde $\dim((X_i\cap Z_j)_\eta)<\dim((Z_j)_\eta)$; por consiguiente, la dimensión común de las componentes irreducibles de $(X_i)_s\cap(Z_j)_s$ es, para todo $s\in S$, estrictamente menor que la dimensión común de las componentes irreducibles de $(Z_j)_s$, lo que demuestra que ninguna de las componentes irreducibles de $(Z_j)_s$ está contenida en una componente irreducible de $(X_i)_s$ para $s\in S$.
\end{proof}

\begin{proposition}[9.9.2]
\label{IV.9.9.2}
Sea $f:X\to S$ un morfismo localmente de presentación finita, sea $\sh{F}$ un $\sh{O}_X$-Módulo cuasicoherente de presentación finita y sea $\Phi$ un subconjunto finito de $\mathbf{Z}\cup\{\pm\infty\}$.
Los subconjuntos siguientes de $X$ son localmente construibles:
\begin{enumerate}[label=(\roman*)]
  \item El conjunto de los $x\in X$ tales que $\sh{F}_{f(x)}$ es equidimensional en el punto $x$ \sref{IV.5.1.12}.
  \item El conjunto de los $x\in X$ tales que las longitudes geométricas de $\sh{F}_{f(x)}$ relativas a $\mathbf{k}(f(x))$, en los puntos genéricos de las componentes irreducibles de $\operatorname{Supp}(\sh{F}_{f(x)})$ que contienen $x$ \sref{IV.4.7.5}, pertenecen a $\Phi$.
  \item El conjunto de los $x\in X$ tales que las dimensiones de los ciclos primos asociados a $\sh{F}_{f(x)}$ y que contienen $x$ pertenecen a $\Phi$.
  \item El conjunto de los $x\in X$ tales que $\sh{F}_{f(x)}$ es geométricamente reducido en el punto $x$ \sref{IV.4.6.17}.
  \item El conjunto de los $x\in X$ tales que $\sh{F}_{f(x)}$ es geométricamente íntegro en el punto $x$ \sref{IV.4.6.22}.
  \item El conjunto de los $x\in X$ tales que $\dim.\operatorname{proj}((\sh{F}_{f(x)})_x)\in\Phi$.
  \item El conjunto de los $x\in X$ tales que $\operatorname{coprof}((\sh{F}_{f(x)})_x)\in\Phi$.
  \item El conjunto de los $x\in X$ tales que $\sh{F}_{f(x)}$ posee la propiedad ($S_k$) en el punto $x$ \sref{IV.5.7.2}.
  \item El conjunto de los $x\in X$ tales que $\sh{F}_{f(x)}$ es un $\sh{O}_{X_{f(x)}}$-Módulo de Cohen-Macaulay en el punto $x$ \sref{IV.5.7.1}.
\end{enumerate}
\end{proposition}

\begin{proof}
\label{IV.9.9.2-proof}
(i) El soporte $Z$ de $\sh{F}$ es localmente construible y cerrado en $X$ \sref{IV.8.9.1}, y considerando un subpreesquema de $X$ cuyo espacio subyacente es $Z$ y que es de presentación finita sobre $S$ \sref{IV.8.9.1}, vemos que la afirmación (i) es un caso particular de \sref{IV.9.9.1},~(iii).

(ii) Todas las propiedades consideradas son locales sobre $X$, y podemos, por tanto, suponer que $X=\operatorname{Spec}(B)$ y $S=\operatorname{Spec}(A)$ son afines y que $f$ es de presentación finita.
Conservamos las notaciones del comienzo de la demostración de \sref{IV.9.9.1}, y elegimos $A_0$ de modo que exista un $\sh{O}_{X_0}$-Módulo coherente $\sh{F}_0$ con $\sh{F}\simeq\sh{F}_0\otimes_{A_0}A$.
Entonces \sref{IV.4.2.7} y \sref{IV.4.7.9} muestran que el conjunto de longitudes geométricas de $(\sh{F}_0)_{f(x_0)}$ en los puntos genéricos de las componentes irreducibles de su soporte que contienen $x_0$ es el mismo que el conjunto correspondiente para $\sh{F}_{f(x)}$ en $x$.
Así, si $E$ (resp.~$E_0$) denota el conjunto de los $x\in X$ (resp.~$x_0\in X_0$) que satisfacen la condición (ii), entonces $E=h^{-1}(E_0)$, y por \sref{IV.1.8.2} podemos restringirnos al caso en que $A$ es noetheriano.
Como en la demostración de \sref{IV.9.9.1}, queda demostrar que para todo $x\in X$ existe un entorno $V$ de $x$ en $\overline{\{x\}}$ sobre el cual es constante el conjunto de longitudes geométricas de $\sh{F}_{f(x')}$ en los puntos genéricos de las componentes irreducibles de su soporte que contienen $x'$.
Si $S'$ es el subpreesquema cerrado reducido de $S$ cuyo espacio subyacente es $\overline{\{f(x)\}}$, y si $X'=f^{-1}(S')$, las fibras de $X$ y $X'$ en los puntos de $S'$ son las mismas,
\oldpage[IV-3]{91}
por lo que podemos sustituir $S$ por $S'$ y $X$ por $X'$; en otras palabras, suponer que $S$ es íntegro y que $\eta=f(x)$ es su punto genérico.
Si $Z=\operatorname{Supp}(\sh{F})$ y los $Z_i$ son las componentes irreducibles de $Z$ que contienen $x$, el argumento de \sref{IV.9.9.1} nos permite suponer que estas son exactamente las componentes irreducibles de $Z$ que encuentran $\overline{\{x\}}$ y que $(Z_i)_s\cap\overline{\{x\}}\neq\varnothing$ para todo $s\in S$.
La conclusión se sigue de \sref{IV.9.8.3} y \sref{IV.9.8.6}.

(iii) Utilizando \sref{IV.4.2.7}, nos reducimos como en (ii) al caso en que $S$ es noetheriano e íntegro y $\eta=f(x)$ es su punto genérico.
Como en la demostración de \sref{IV.9.9.1}, queda encontrar un entorno $V$ de $x$ en $\overline{\{x\}}$ sobre el cual sea constante el conjunto de dimensiones de los ciclos primos asociados a $\sh{F}_{f(x')}$ y que contienen $x'$.
Sean $Z'_i$ las adherencias en $X$ de los ciclos primos asociados a $\sh{F}_\eta$ que contienen $x$ (cf.~\sref{IV.9.8.2}).
Por \sref{IV.9.8.3} y \sref{IV.9.5.1}, para todo $s$ suficientemente próximo a $\eta$, los ciclos primos asociados a $\sh{F}_s$ que encuentran $\overline{\{x\}}$ son componentes irreducibles de los $(Z'_i)_s$; por \sref{IV.9.5.6}, todas estas componentes tienen dimensión $\dim((Z'_i)_\eta)$, lo que demuestra la afirmación.

(iv) Razonamos como en (iii), utilizando ahora \sref{IV.4.7.11}.
Con las mismas notaciones, los ciclos primos asociados a $\sh{F}_s$ que son componentes irreducibles de $(Z'_i)_s$ están inmersos si y solo si $(Z'_i)_\eta$ es un ciclo primo asociado inmerso de $\sh{F}_\eta$.
Se sigue que, según que $x$ pertenezca a un ciclo primo inmerso asociado de $\sh{F}_\eta$ o a ninguno, los puntos $x'\in\overline{\{x\}}$ que poseen la propiedad análoga para $\sh{F}_{f(x')}$ forman un entorno de $x$ en $\overline{\{x\}}$.
La conclusión se sigue de esta observación, de la caracterización de los Módulos geométricamente reducidos en \sref{IV.4.7.10}, y de (ii).

(v) Teniendo en cuenta \sref{IV.4.7.11}, nos reducimos de nuevo al caso en que $S$ es noetheriano e íntegro y en que $\eta=f(x)$ es su punto genérico.
Sea $Y$ un subpreesquema cerrado de $X$ cuyo espacio subyacente es $\operatorname{Supp}(\sh{F})$.
Como al principio de la demostración de \sref{IV.9.7.7}, existe una $A$-álgebra íntegra finita $A'$ (de modo que, si $S'=\operatorname{Spec}(A')$, el morfismo $S'\to S$ es finito y sobreyectivo) tal que, si $Y'=Y_{(S')}$ y $\eta'$ es el punto genérico de $S'$, las componentes irreducibles de $Y'_{\eta'}$ son geométricamente irreducibles.
Si $X'=X_{(S')}$, la proyección $h:X'\to X$ es finita y sobreyectiva, por tanto cerrada; por consiguiente, para un punto $x'$ situado sobre $x$, se tiene $h(\overline{\{x'\}})=\overline{\{x\}}$, y la imagen por $h$ de un entorno de $x'$ en $\overline{\{x'\}}$ es un entorno de $x$ en $\overline{\{x\}}$.
Escribiendo $\sh{F}'=\sh{F}_{(S')}$ y utilizando \sref{IV.4.7.11}, nos reducimos a demostrar que, según que $\sh{F}'_{\eta'}$ sea o no geométricamente íntegro en $x'$, los puntos $y'\in\overline{\{x'\}}$ donde $\sh{F}'_{f'(y')}$ es o no geométricamente íntegro en $y'$ forman un entorno de $x'$ en $\overline{\{x'\}}$.
Podemos, por tanto, suponer que $X'=X$ y que las componentes irreducibles de $Y_\eta$ son geométricamente irreducibles.
Elijamos subconjuntos cerrados $Z_i$ de $X$ tales que los $Z_i\cap X_\eta$ sean las
  \oldpage[IV-3]{92}
  componentes irreducibles de $Y_\eta$.
Por \sref{IV.9.7.7}, \sref{IV.9.7.8} y \sref{IV.9.5.3}, para todo $s$ próximo a $\eta$ los $(Z_i)_s$ son las componentes irreducibles de $Y_s$ y son geométricamente irreducibles.
En un punto $y\in X_s$, el Módulo $\sh{F}_s$ es geométricamente íntegro exactamente cuando es geométricamente reducido allí y $y$ pertenece a uno solo de los $(Z_i)_s$ \sref{IV.4.6.22}.
La conclusión se sigue de (iv) y de \sref{IV.9.5.1} aplicado a $\overline{\{x\}}\cap Z_i$ para cada $i$.

(vi) Conservando las notaciones de (ii), \sref{IV.6.2.1} da
\[
  \dim.\operatorname{proj}((\sh{F}_s)_x)
  =\dim.\operatorname{proj}(((\sh{F}_0)_{s_0})_{x_0}),
\]
de modo que podemos restringirnos de nuevo al caso en que $A$ es noetheriano. Queda demostrar que para todo $x\in X$ existe un entorno $V$ de $x$ en $\overline{\{x\}}$ tal que
\[
  \dim.\operatorname{proj}((\sh{F}_{f(x')})_{x'})
  =\dim.\operatorname{proj}((\sh{F}_{f(x)})_x)
\]
para todo $x'\in V$. Como antes, podemos suponer que $S$ es íntegro y que $\eta=f(x)$ es su punto genérico, de modo que $(\sh{F}_\eta)_x=\sh{F}_x$. Por planitud genérica \sref{IV.6.9.1}, tras sustituir $S$ por un entorno abierto de $\eta$, podemos suponer que $f$ es plano y que $\sh{F}$ es $f$-plano. Entonces \sref{IV.6.2.3} da
\[
  \dim.\operatorname{proj}((\sh{F}_{f(x')})_{x'})
  =\dim.\operatorname{proj}(\sh{F}_{x'})
\]
para todo $x'\in X$. Por \sref{IV.6.11.1}, tras sustituir $X$ por un entorno abierto de $x$, podemos suponer que
\[
  \dim.\operatorname{proj}(\sh{F}_{x'})
  \leqslant\dim.\operatorname{proj}(\sh{F}_x)
\]
para todo $x'\in X$.

Si $\dim.\operatorname{proj}(\sh{F}_x)=n$, existe un $\sh{O}_x$-módulo finito $M$ tal que $\operatorname{Ext}^n_{\sh{O}_x}(\sh{F}_x,M)\neq 0$ \sref[0]{0.17.2.4}.
Tras sustituir $X$ por un entorno abierto de $x$, existe un $\sh{O}_X$-módulo coherente $\sh{G}$ tal que $M=\sh{G}_x$ \sref[0]{0.5.3.8};
en virtud de (T, \sref{IV.4.2.2}), tenemos $(\shExt^n_{\sh{O}_X}(\sh{F},\sh{G}))_x\neq 0$.
Pero $\shExt^n_{\sh{O}_X}(\sh{F},\sh{G})$ es un $\sh{O}_X$-Módulo coherente ($0_{\mathrm{III}}$, \sref{IV.12.3.3}), por tanto su soporte $Y$ es cerrado ($0_I$, \sref{IV.5.2.2}); como contiene $x$, contiene también $\overline{\{x\}}$, de donde concluimos (aplicando de nuevo (T, \sref{IV.4.2.2})) que $\dim.\operatorname{proj}(\sh{F}_{x'})\geqslant n$ para todo $x'\in\overline{\{x\}}$, lo que termina la demostración de la afirmación en el caso (vi).

(vii) Como $B$ es una $A$-álgebra de tipo finito, $X$ es $S$-isomorfo a un \emph{subesquema cerrado} de un $S$-esquema de la forma $Y=\operatorname{Spec}(A[T_1,\ldots,T_r])$; si $j:X\to Y$ es la inyección canónica y $\sh{G}=j_*(\sh{F})$, entonces $\sh{G}_s=(j_s)_*(\sh{F}_s)$ para todo $s\in S$, y, en virtud de ($0$, \sref{IV.16.4.11}), podemos restringirnos a demostrar la afirmación para $Y$ y $\sh{G}$.
En otras palabras, podemos suponer que $B=A[T_1,\ldots,T_r]$, de modo que los esquemas $X_s=\operatorname{Spec}(\mathbf{k}(s)[T_1,\ldots,T_r])$ son \emph{regulares} ($0$, \sref{IV.17.3.7}).
Sea $W=\operatorname{Supp}(\sh{F})$ y $W_s=\operatorname{Supp}(\sh{F}_s)$ \sref[I]{I.9.1.13}; tenemos, por \sref{IV.6.11.2.1},
\begin{equation}
\tag{9.9.2.1}
\label{IV.9.9.2.1}
  \operatorname{coprof}((\sh{F}_{f(x)})_x)=\dim.\operatorname{proj}((\sh{F}_{f(x)})_x)-\operatorname{codim}_x(W_{f(x)},X_{f(x)}).
\end{equation}

Pero como $W$ es construible \sref{IV.8.9.1} y cada $W_s$ es cerrado, se sigue de (vi) y \sref{IV.9.9.1},~(ii), que las dos funciones del segundo miembro de \sref{IV.9.9.2.1} son construibles; por tanto, también lo es su diferencia, lo que demuestra la afirmación en el caso (vii).

(viii) Sea $U_n$ el conjunto de los $x\in X$ tales que $\operatorname{coprof}((\sh{F}_{f(x)})_x)\leqslant n$, y pongamos $Z_n=X-U_n$; se sigue de (vii) que los $Z_n$ son \emph{construibles}.
\oldpage[IV-3]{93}
Además, como la función $x\mapsto\dim_x(W_{f(x)})$ es construible por \sref{IV.9.9.1},~(i), toma solo un número finito de valores, y los números $\dim(W_{f(x)})$ tienen, por tanto, una cota superior finita $m$ cuando $x$ recorre $X$.
Como $\operatorname{coprof}((\sh{F}_{f(x)})_x)\leqslant\dim((\sh{F}_{f(x)})_x)\leqslant\dim(W_{f(x)})$, tenemos $Z_n=\varnothing$ para $n\geqslant m$.
Finalmente, \sref{IV.6.11.2},~(i), muestra que $(Z_n)_s$ es cerrado en $X_s$ para todo $n$ y todo $s\in S$.
Por \sref{IV.5.7.4}, el conjunto de los $x\in X$ donde $(\sh{F}_{f(x)})_x$ posee la propiedad ($S_k$) es el conjunto de los $x\in X$ que satisfacen todas las relaciones
\begin{equation}
\tag{9.9.2.2}
\label{IV.9.9.2.2}
  \operatorname{codim}_x((Z_n)_{f(x)},W_{f(x)})>n+k
\end{equation}
para todo $n\geqslant 0$; como esta relación se satisface automáticamente para $n\geqslant m$, basta considerar las relaciones \sref{IV.9.9.2.2} para $0\leqslant n<m$.
Pero, en virtud de \sref{IV.9.9.1},~(ii), el conjunto $V_{n,k}$ de los $x$ que satisfacen \sref{IV.9.9.2.2} es construible, y lo mismo vale para la intersección de estos conjuntos para $0\leqslant n<m$.

(ix) La afirmación se sigue inmediatamente de (vii), pues el conjunto de los $x\in X$ donde $(\sh{F}_{f(x)})_x$ es un módulo de Cohen-Macaulay está definido por la relación $\operatorname{coprof}((\sh{F}_{f(x)})_x)=0$.
\end{proof}

\begin{corollary}[9.9.3]
\label{IV.9.9.3}
Sea $f:X\to S$ un morfismo de presentación finita, sea $\sh{F}$ un $\sh{O}_X$-Módulo cuasicoherente de presentación finita y sea $\mathbf{P}$ una de las propiedades (i) a (ix) de \sref{IV.9.9.2}.
Entonces el conjunto de los $s\in S$ tales que la propiedad $\mathbf{P}$ es cierta en todos los puntos $x\in X_s$, es localmente construible en $S$.
\end{corollary}

\begin{proof}
\label{IV.9.9.3-proof}
En efecto, su complemento es la imagen por $f$ del complemento del conjunto $E$ de los puntos donde $\mathbf{P}$ es cierta.
Como $E$ es localmente construible en $X$, también lo es $X-E$, y por tanto $f(X-E)$ es localmente construible en $S$ en virtud del teorema de Chevalley \sref{IV.1.8.4}.
\end{proof}

\begin{proposition}[9.9.4]
\label{IV.9.9.4}
Sea $f:X\to S$ un morfismo localmente de presentación finita.
El conjunto de los $x\in X$ tales que $X_{f(x)}$ posee en el punto $x$ una de las propiedades (determinadas) siguientes:
\begin{enumerate}[label=(\roman*)]
  \item ser geométricamente regular;
  \item poseer la propiedad ($R_k$) geométrica;
  \item ser geométricamente normal;
  \item ser geométricamente reducido (es decir, separable);
  \item ser geométricamente puntualmente íntegro;
\end{enumerate}
es un subconjunto localmente construible de $X$.
\end{proposition}

\begin{proof}
\label{IV.9.9.4-proof}
Para las propiedades (iv) y (v), esto se sigue de \sref{IV.9.9.2},~(iv) y~(v), aplicado a $\sh{F}=\sh{O}_X$.
Para las demás propiedades, teniendo en cuenta \sref{IV.6.7.8}, nos reducimos, como al principio de \sref{IV.9.9.2},~(ii), al caso en que $S$ es noetheriano e íntegro y de punto genérico $\eta=f(x)$.
Además, en virtud del teorema de planitud genérica \sref{IV.6.9.1}, podemos, sustituyendo $S$ por un entorno abierto de $\eta$, suponer que $f$ es un morfismo \emph{plano}.
Decir que $X_{f(y)}$ es geométricamente regular en el punto $y$ significa que $f$ es \emph{regular} en el punto $y$ \sref{IV.6.8.1}, y sabemos que el conjunto $L$ de estos puntos es \emph{abierto} en $X$ \sref{IV.6.8.7}, lo que demuestra la proposición en el caso (i).

Para demostrar el caso (ii), pongamos $F=X-L$, que es \emph{cerrado} en $X$.
Decir que $X_s$
\oldpage[IV-3]{94}
satisface la propiedad geométrica ($R_k$) en el punto $y\in f^{-1}(s)$ significa o bien que $y\in L$, o bien que los puntos genéricos $z_i$ de las componentes irreducibles del conjunto cerrado $F_s$ que contienen $y$ satisfacen la relación $\dim(\sh{O}_{X_s,z_i})\geqslant k+1$ (teniendo en cuenta \sref{IV.4.2.7} y \sref{IV.5.2.3}); en otras palabras, los puntos $y\in F_s$ donde $X_s$ satisface la propiedad ($R_k$) geométrica son aquellos tales que $\operatorname{codim}_y(F_s,X_s)\geqslant k+1$ \sref{IV.5.1.2}.
La conclusión se sigue, pues, de (i) y \sref{IV.9.9.1},~(ii).

Finalmente, en el caso (iii), la conclusión se sigue de (ii), de \sref{IV.9.9.2},~(viii), del hecho de que la propiedad ($S_k$) es estable por extensión del cuerpo base \sref{IV.6.7.1} y, finalmente, del criterio de Serre \sref{IV.5.8.6}.
\end{proof}

\begin{corollary}[9.9.5]
\label{IV.9.9.5}
Sea $f:X\to S$ un morfismo de presentación finita, y denotemos por $\mathbf{P}$ una de las propiedades (i) a (v) de \sref{IV.9.9.4}.
Entonces el conjunto de los $s\in S$ tales que la propiedad $\mathbf{P}$ es cierta en todos los puntos $x\in X_s$, es localmente construible en $S$.
\end{corollary}

\begin{proof}
\label{IV.9.9.5-proof}
La demostración a partir de \sref{IV.9.9.4} es la misma que la de \sref{IV.9.9.3} a partir de \sref{IV.9.9.2}.
\end{proof}

\begin{proposition}[9.9.6]
\label{IV.9.9.6}
Sea $f:X\to S$ un morfismo localmente de presentación finita, y sea $\sh{L}_\bullet$ un complejo formado por $\sh{O}_X$-Módulos cuasicoherentes de presentación finita; para todo $s\in S$, sea $(\sh{L}_\bullet)_s$ el complejo $\sh{L}_\bullet\otimes_{\sh{O}_S}\mathbf{k}(s)$ de $\sh{O}_{X_s}$-Módulos coherentes de tipo finito.
Entonces, para un entero dado $n$, el conjunto de los $x\in X$ tales que $(\sh{H}_n((\sh{L}_\bullet)_{f(x)}))_x=0$ es localmente construible en $X$.
\end{proposition}

\begin{proof}
\label{IV.9.9.6-proof}
Podemos restringirnos al caso en que $\sh{L}_i=0$ salvo para $i=0$, $1$ o $2$, y en que $n=1$.
Además, como la cuestión es local sobre $X$, podemos restringirnos al caso en que $S=\operatorname{Spec}(A)$ y $X=\operatorname{Spec}(B)$ son afines, siendo $B$ una $A$-álgebra de presentación finita.
Existe entonces un subanillo noetheriano $A_0$ de $A$, un $A_0$-preesquema de tipo finito $X_0$ y un complejo $\sh{L}_\bullet^{(0)}$ de $\sh{O}_{X_0}$-Módulos coherentes, nulo salvo en dimensiones $0$, $1$ y $2$ y tal que $X=X_0\otimes_{A_0}A$ y $\sh{L}_\bullet=\sh{L}_\bullet^{(0)}\otimes_{A_0}A$.
Para todo $s\in S$, si $s_0$ es la proyección de $s$ en $S_0=\operatorname{Spec}(A_0)$, entonces
\[
  X_s=(X_0)_{s_0}\otimes_{\mathbf{k}(s_0)}\mathbf{k}(s),
\]
y la proyección $X_s\to(X_0)_{s_0}$ es fielmente plana.
Por consiguiente,
\[
  \sh{H}_n((\sh{L}_\bullet)_s)
  =\sh{H}_n((\sh{L}_\bullet^{(0)})_{s_0})\otimes_{\mathbf{k}(s_0)}\mathbf{k}(s).
\]
Si $E$ (resp.~$E_0$) es el conjunto de los $x\in X$ (resp.~$x_0\in X_0$) para los cuales
\[
  \begin{aligned}
  (\sh{H}_n((\sh{L}_\bullet)_{f(x)}))_x&=0\\
  \bigl(\text{resp. }(\sh{H}_n((\sh{L}_\bullet^{(0)})_{f_0(x_0)}))_{x_0}&=0\bigr),
  \end{aligned}
\]
entonces $E=h^{-1}(E_0)$, donde $h:X\to X_0$ es la proyección canónica.
Por \sref{IV.1.8.2}, podemos, por tanto, restringirnos al caso en que $A$ es noetheriano.
Queda demostrar ($0_{\mathrm{III}}$, \sref{IV.9.2.3}) que si $x\in X$ es tal que $(\sh{H}_n((\sh{L}_\bullet)_{f(x)}))_x=0$ (resp.~$\neq 0$), existe un entorno abierto $V$ de $x$ en $\overline{\{x\}}$ tal que para todo $x'\in V$ se tiene $(\sh{H}_n((\sh{L}_\bullet)_{f(x')}))_{x'}=0$ (resp.~$\neq 0$).
Sustituyendo $S$ por el subpreesquema cerrado reducido cuyo espacio subyacente es $\overline{\{f(x)\}}$, podemos suponer que $S$ es íntegro y que $f(x)=\eta$ es su punto genérico.
Entonces, restringiendo $S$ a un entorno abierto de $\eta$, podemos suponer que para todo $s\in S$,
$(\sh{H}_n(\sh{L}_\bullet))_s=\sh{H}_n((\sh{L}_\bullet)_s)$ \sref{IV.9.4.3}; por consiguiente, si $Z$ es el soporte de $\sh{H}_n(\sh{L}_\bullet)$, entonces el soporte de $\sh{H}_n((\sh{L}_\bullet)_s)$ es $Z_s=Z\cap X_s$ \sref[I]{I.9.1.13.1}.
La hipótesis es que $Z_\eta\cap\overline{\{x\}}=\varnothing$ (resp.~$\neq\varnothing$).
Como $Z\cap\overline{\{x\}}$ es cerrado en el espacio noetheriano $X$, concluimos de \sref{IV.9.5.1} que existe un entorno $U$ de $\eta$ en $S$ tal que, para todo $s\in U$, tenemos $Z_s\cap\overline{\{x\}}=\varnothing$ (resp.~$\neq\varnothing$).
El conjunto $V=f^{-1}(U)\cap\overline{\{x\}}$ responde así a la cuestión.
\end{proof}
